% concatenated from DPRSRM_arXiv.tex
% SIAM Article Template
% \pdfoutput=1
\documentclass[onefignum,onetabnum]{siamart220329}
\usepackage{graphics,graphicx,epstopdf,url}
\usepackage{arydshln}
\usepackage{color}
%\usepackage{algorithm2e}
\usepackage{geometry}
\usepackage{algorithmic}
\usepackage{url,cases,supertabular}
\usepackage{amssymb,mathtools}
\usepackage{enumerate,makecell}
\usepackage{amsfonts}
%\usepackage{amsthm}
\usepackage{graphicx,epstopdf}
\usepackage{color,verbatim}
\usepackage{pifont}
\usepackage{mathrsfs,subeqnarray,subfigure}
\usepackage{listings,array}
\usepackage{booktabs}
\usepackage{dashrule}
\usepackage{arydshln}
\usepackage{graphicx}
\usepackage{amsfonts}
\usepackage{mathrsfs}
%\usepackage[normalem]{ulem}   \usepackage{enumerate}
\usepackage{graphicx}  \usepackage{epstopdf}
\usepackage{subfigure}  \usepackage{lineno}
\usepackage{multirow,bm}
\usepackage{color}
\usepackage{algorithm}
\usepackage{algorithmic}
\usepackage{longtable}
\usepackage{mdwtab}
\usepackage{mdwmath}
\usepackage{bbding}
\usepackage{booktabs}
\usepackage{rotating}
\usepackage{enumitem}
\usepackage{lineno}
%\usepackage{changes}
\usepackage{url}
\usepackage{pdflscape}
\usepackage{makecell}
%\usepackage{hyperref}
\newcommand{\comm}[1]{{\color{red}#1}}
\newcommand{\revise}[1]{#1}
\newcommand{\abstractrevision}[1]{{\color{blue}#1}}
\newenvironment{revisionblock}{}{}


%\usepackage[left=1in,top=1in,right=1in,bottom=1in,a4paper]{geometry}


\numberwithin{equation}{section}

\newtheorem{Theorem}{Theorem}[section]
% \newtheorem{corollary}[theorem]{Corollary}
\newtheorem{Lemma}[theorem]{Lemma}
\newtheorem{Proposition}[theorem]{Proposition}
\newtheorem{Definition}[theorem]{Definition}
\newtheorem{example}{Example}[section]
\newtheorem{assumption}[theorem]{Assumption}
 \newtheorem{remark}[theorem]{Remark}
% \newtheorem{question}{Question}[section]


\def\red#1{\textcolor{red}{#1}}
\def\black#1{#1}

\newcommand{\T}{\top}
\newcommand{\E}{\mathbb{E}}
\newcommand\numberthis{\addtocounter{equation}{1}\tag{\theequation}}
\newcommand{\etal}{ et al. }
\newcommand{\br}{\mathbb{R}}

\newcommand{\ba}{\begin{array}}
\newcommand{\ea}{\end{array}}
\newcommand{\Snn}{\mathbb{S}^{n\times n}}

\newcommand{\upcite}[1]{\textsuperscript{\textsuperscript{\cite{#1}}}}
\newcommand{\bit}{\begin{itemize}}
\newcommand{\eit}{\end{itemize}}
\newcommand{\be}{\begin{equation}}
\newcommand{\ee}{\end{equation}}
\newcommand{\bee}{\begin{equation*}}
\newcommand{\eee}{\end{equation*}}
\newcommand{\bea}{\begin{eqnarray}}
\newcommand{\eea}{\end{eqnarray}}
\newcommand{\zero}{\textbf{0}~}
\newcommand{\mvec}{\mbox{vec}}
\newcommand{\mat}{\mbox{mat}}
\newcommand{\orth}{\mathsf{St}_{n,p}}
\newcommand{\stief}{\mathsf{St}_{d,r}}
\newcommand{\grass}{\mathsf{Gr}_{d,r}}
\newcommand{\mskew}{\mathrm{skew}}
\newcommand{\mF}{f}
\newcommand{\mL}{\mathcal{L}}
\newcommand{\Rbb}{\mathbb{R}}
\newcommand{\mbC}{\mathbb{C}}
\newcommand{\mT}{\mathcal{T}}
\newcommand{\teigs}{\textbf{eigs}}
\newcommand{\mN}{\mathcal{N}}
\newcommand{\mbsign}{\mbox{\bf sign}}
\newcommand{\st}{\mathrm{s.t.}}
\newcommand{\mE}{\mathcal{E}}
\newcommand{\mR}{\mathcal{R}}
\newcommand{\wR}{\widetilde{R}}
\newcommand{\wN}{\widetilde{N}}
\newcommand{\mM}{\mathcal{M}}
\newcommand{\mD}{\mathcal{D}}
\newcommand{\Drho}{D_{\rho}}
\newcommand{\argmin}{\mathop{\mathrm{arg\,min}}}
\newcommand{\argmax}{\mathop{\mathrm{arg\,max}}}
\newcommand{\odr}{\overline{\delta r}}
\newcommand{\udr}{\underline{\delta r}}
\newcommand{\oq}{\overline{q}}
\newcommand{\uq}{\underline{q}}
\newcommand{\ola}{\overline{\lambda}}
\newcommand{\ula}{\underline{\lambda}}
\newcommand{\oom}{\overline{\omega}}
\newcommand{\uom}{\underline{\omega}}
%\DeclareMathOperator{\argmin}{\operatorname*{arg\,min}}
%\DeclareMathOperator{\argmax}{\operatorname*{arg\,max}}
\newcommand{\myarg}{\mathop{\mathrm{arg}}}
\newcommand{\Drhok}{D_{\rho, k}}
\newcommand{\mP}{\mathcal{P}}
\newcommand{\mTr}{\text{Tr}}
\newcommand{\mtr}{\mbox{tr}}
\newcommand{\tW}{\widetilde{W}}
\newcommand{\mK}{\mathcal{K}}
%\newcommand{P_XD}{\widetilde{g}}
\newcommand{\tJ}{\widetilde{J}}
\newcommand{\tF}{\widetilde{F}}
\newcommand{\eqreff}[1]{(\eqref{#1})}
\newcommand{\nii}{\noindent}
\newcommand{\qf}{\text{qf}}
\newcommand{\nn}{\nonumber}
\newcommand{\diag}{\text{diag}}
\newcommand{\Diag}{\text{Diag}}
\newcommand{\proj}{\mathcal{P}}
\newcommand{\by}{\mathrm{\bf y}}
\newcommand{\bs}{\mathrm{\bf s}}
\newcommand{\bx}{\mathbf{x}}
\newcommand{\bd}{\mathbf{d}}
\newcommand{\bq}{\mathbf{q}}
\newcommand{\bu}{\mathbf{u}}
\newcommand{\bv}{\mathbf{v}}
\newcommand{\bW}{\mathbf{W}}
\newcommand{\mfS}{\mathbf{S}}
\newcommand{\mfY}{\mathbf{Y}}
\newcommand{\mfR}{\mathbf{R}}
\newcommand{\LBB}{\mathrm{LBB}}
\newcommand{\SBB}{\mathrm{SBB}}
\newcommand{\mfD}{\mathbf{D}}
\newcommand{\mfL}{\mathbf{L}}
\newcommand{\mscrL}{\mathscr{L}}
\newcommand{\nF}{\nabla \mathcal{F}}
\newcommand{\rmn}[1]{\romannumeral#1}
\newcommand{\Rmn}[1]{\uppercase\expandafter{\romannumeral#1}}
\renewcommand{\theenumi}{\roman{enumi}}
\renewcommand{\labelenumi}{\theenumi}
\newcommand{\rev}[1]{{\color{red}{#1}}}
\newcommand{\rblue}[1]{#1}
\newcommand{\rblack}[1]{{\color{black}{#1}}}
%\renewcommand{\smartqed}{\hfill{\qed}}
\newcommand{\Fsf}{\mathsf{F}}
\newcommand{\Tsf}{\mathsf{T}}
\newcommand{\Xcal}{\mathcal{X}}
\newcommand{\Ebb}{\mathbb{E}}
\newcommand{\BB}{\mathrm{BB}}
\newcommand{\Bsf}{\mathsf{B}}
\newcommand{\Tcal}{\mathcal{T}}
\newcommand{\Rcal}{\mathcal{R}}
\newcommand{\llangle}{\left \langle}
\newcommand{\rrangle}{\right \rangle}
\newcommand{\tx}{\tilde{X}}
\newcommand{\tX}{\tilde{X}}
\newcommand{\Gcal}{\mathcal{G}}
\newcommand{\Dcal}{\Gcal^{\mathrm{R}}}
\newcommand{\Pcal}{\mathcal{P}}
\newcommand{\Fcal}{\mathcal{F}}
\newcommand{\Kcal}{\mathcal{K}}
\newcommand{\Ocal}{\mathcal{O}}
\newcommand{\Tbf}{\mathbf{T}}
\newcommand{\Rscr}{\mathcal{R}}
\newcommand{\ud}{\mathrm{d}}
\numberwithin{equation}{section}
\newcommand{\Pbb}{\mathbb{P}}
\newcommand{\Dbf}{\mathbf{D}}
\newcommand{\Drm}{\mathrm{D}}
\newcommand{\drm}{\mathrm{d}}
\newcommand{\rgrad}{\mathrm{grad}}
\newcommand{\up}{\mathrm{up}}
\newcommand{\qr}{\mathrm{qr}}
\newcommand{\Mcal}{\mathcal{M}}
\newcommand{\xbf}{\mathbf{x}}
\newcommand{\ubf}{\mathbf{u}}
\newcommand{\abf}{\mathbf{a}}
\newcommand{\bbf}{\mathbf{b}}
\newcommand{\cbf}{\mathbf{c}}
\newcommand{\dbf}{\mathbf{d}}
\newcommand{\fbf}{\mathbf{f}}
\newcommand{\Lbf}{\mathbf{L}}
\newcommand{\Nsf}{\mathsf{N}}
\newcommand{\TR}{\#\mathrm{N}_{\Rscr}}
\newcommand{\TCG}{\#\mathrm{N}_{\mathrm{cg}}}
\newcommand{\TT}{\#\mathrm{N}_{\mathrm{\mathcal{T}}}}
\newcommand{\1}{\mathbf{1}}
\newcommand{\grad}{\mathrm{grad}}
\newcommand{\sym}{\mathrm{sym}}
\newcommand{\R}{\mathbb{R}}
\newcommand{\hess}{\mathrm{Hess}}
%\renewcommand\arraystretch{10}
\numberwithin{theorem}{section}
\newcommand{\iprod}[2]{\left \langle #1, #2 \right \rangle }
\newcommand{\V}{\mathcal{V}}
\newcommand{\tr}{\mathrm{tr}}
\newcommand{\Diff}{\mathrm{Diff}}
\providecommand{\xx}{\mathbf{x}}
\providecommand{\cD}{\mathcal{D}}
% Information that is shared between the article and the supplement
% (title and author information, macros, packages, etc.) goes into
% ex_shared.tex. If there is no supplement, this file can be included
% directly.

%\input{SISC_shared}

% SIAM Shared Information Template
% This is information that is shared between the main document and any
% supplement. If no supplement is required, then this information can
% be included directly in the main document.


% Packages and macros go here
\usepackage{lipsum}
\usepackage{clrscode}
\usepackage{appendix}
\usepackage{bm}
\usepackage{booktabs}
\usepackage{url}
%\usepackage{graphicx,epstopdf}
%\usepackage[caption=false]{subfig}
\usepackage{multirow}
\usepackage{textcomp}
\usepackage{amsmath}
\usepackage{amsfonts}
\usepackage{amssymb}
\usepackage{mathrsfs}
\usepackage{hyperref}
\usepackage{graphicx,graphics,subfigure}
\usepackage{hyperref,url}
\usepackage{epsf,epstopdf}
%\usepackage[bookmarks=false]{hyperref}
% package `geometry' implicily changes the page layout (margins)
%\usepackage{geometry}
\usepackage{algorithm}
\usepackage{algorithmic}
%\usepackage{essay-def-siam}
%\usepackage{clrscode} % this is incompatiable with algpseudocode
%\usepackage{algpseudocode}
\usepackage{bbm}
\usepackage{dsfont}
%\algnewcommand{\IIf}[1]{\State\algorithmicif\ #1\ \algorithmicthen}
%\algnewcommand{\EndIIf}{\unskip\ \algorithmicend\ \algorithmicif}
%\algnewcommand{\IfThenElse}[3]{% \IfThenElse{<if>}{<then>}{<else>}
%  \State \algorithmicif\ #1\ \algorithmicthen\ #2\ \algorithmicelse\ #3 \algorithmicend\ \algorithmicif}
%\newcommand{\red}{\color{red}}
%\usepackage{algorithmic}
\renewcommand{\algorithmicensure}{ \textbf{Output:}} %UseOutput in the format of Algorithm
\renewcommand{\algorithmicrequire}{ \textbf{Input:}} %Use Input in the format of Algorithm
\usepackage{indentfirst}
% \usepackage{natbib}
\usepackage{pdfpages}
\usepackage{pdflscape}
\usepackage{longtable}
\ifpdf
  \DeclareGraphicsExtensions{.eps,.pdf,.png,.jpg}
\else
  \DeclareGraphicsExtensions{.eps}
\fi
\allowdisplaybreaks[1]

% Add a serial/Oxford comma by default.
\newcommand{\creflastconjunction}{, and~}

% Used for creating new theorem and remark environments
%\newsiamremark{remark}{Remark}
%\newsiamremark{hypothesis}{Hypothesis}
%\crefname{hypothesis}{Hypothesis}{Hypotheses}
%\newsiamthm{claim}{Claim}

% Sets running headers as well as PDF title and authors
\headers{Decentralized projected Riemannian gradient method}{Kangkang Deng, Jiang Hu}

% Title. If the supplement option is on, then "Supplementary Material"
% is automatically inserted before the title.
\title{Decentralized projected Riemannian stochastic recursive momentum method for smooth optimization on compact submanifolds}
% \thanks{Submitted to the editors DATE.
%\funding{Z. Wen was supported in part by the NSFC grant 11831002.}
%}

% Authors: full names plus addresses.
\author{Kangkang Deng\thanks{ Department of Mathematics,  National University of Defense Technology, Changsha, 410073,
China (\email{freedeng1208@gmail.com}).}
\and Jiang Hu\thanks{Corresponding author. Department of Mathematics, University of California, Berkeley, CA 94720, US
(\email{hujiangopt@gmail.com}).}
}


\usepackage{amsopn}

% Optional PDF information
\ifpdf
\hypersetup{
  pdftitle={Decentralized projected Riemannian gradient method},
  pdfauthor={Kangkang Deng, Jiang Hu}
}
\fi
% The next statement enables references to information in the
% supplement. See the xr-hyperref package for details.
% \makeatletter
% \newcommand*{\addFileDependency}[1]{% argument=file name and extension
%   \typeout{(#1)}
%   \@addtofilelist{#1}
%   \IfFileExists{#1}{}{\typeout{No file #1.}}
% }
% \makeatother

%\newcommand*{\myexternaldocument}[1]{%
 %\externaldocument{#1}%
 %\addFileDependency{#1.tex}%
 %\addFileDependency{#1.aux}%
%}
%\myexternaldocument{SISC_supplement}
%\externaldocument{SISC_supplement}

% FundRef data to be entered by SIAM
%<funding-group specific-use="FundRef">
%<award-group>
%<funding-source>
%<named-content content-type="funder-name">
%</named-content>
%<named-content content-type="funder-identifier">
%</named-content>
%</funding-source>
%<award-id> </award-id>
%</award-group>
%</funding-group>
\allowdisplaybreaks[2]
\geometry{left=4.0cm, right = 4.0cm}
\begin{document}

\maketitle

\begin{abstract}
This paper studies decentralized optimization over a compact submanifold within a communication network of $n$ nodes, where each node possesses a smooth non-convex local cost function, and the goal is to jointly minimize the sum of these local costs. We focus particularly on the online setting, where local data is processed in real-time as it streams in, without the need for full data storage. We propose a decentralized projected Riemannian stochastic recursive momentum (DPRSRM) method that employs local hybrid stochastic gradient estimators and uses the network to track the global gradient. DPRSRM achieves an oracle complexity of  \(\mathcal{O}(\epsilon^{-\frac{3}{2}})\), outperforming existing methods that have at most \(\mathcal{O}(\epsilon^{-2})\) complexity. Our method requires only $\mathcal{O}(1)$ gradient evaluations per iteration for each local node and does not require restarting with a large batch gradient.   Furthermore, we demonstrate the effectiveness of our proposed methods compared to state-of-the-art ones through numerical experiments on principal component analysis problems and low-rank matrix completion. \abstractrevision{This revised version corrects the mean-square oracle condition and the estimator/tracking analysis, and replaces the original separate error summation by a coupled Lyapunov argument; Algorithm~\ref{alg:drgta} and the $\mathcal O(\epsilon^{-3/2})$ oracle-complexity order are unchanged.}



% This paper studies decentralized optimization over a compact submanifold within a communication network of $n$ nodes, where each node possesses a smooth non-convex local cost function, and the goal is to jointly minimize the sum of these local costs. We focus particularly on the online setting, where local data is processed in real-time as it streams in, without the need for full data storage. We propose a decentralized projected Riemannian stochastic recursive momentum (DPRSRM) method that employs local hybrid stochastic gradient estimators and uses the network to track the global gradient. DPRSRM achieves an oracle complexity of  \(\mathcal{O}(\epsilon^{-\frac{3}{2}})\), outperforming existing methods that have at most \(\mathcal{O}(\epsilon^{-2})\) complexity. Our method requires only $\mathcal{O}(1)$ gradient evaluations per iteration for each local node and does not require restarting with a large batch gradient.   Furthermore, we demonstrate the effectiveness of our proposed methods compared to state-of-the-art ones through numerical experiments on principal component analysis problems and low-rank matrix completion.
\end{abstract}
\begin{keywords}
decentralized optimization, Riemannian manifold, gradient tracking, consensus
\end{keywords}

% REQUIRED
\begin{AMS}
  90C06, 90C22, 90C26, 90C56
\end{AMS}


\section{Introduction}


Decentralized optimization has emerged as a prominent area of research, particularly for its application in large-scale systems, such as sensor networks, distributed computing, and machine learning. In these contexts, data is often partitioned across numerous nodes, rendering centralized optimization approaches impractical due to challenges such as privacy limitations and restricted computational resources. In this work, we are concerned with the distributed smooth optimization on a compact submanifold
\be \label{prob:original}
\begin{aligned}
  \min \quad &  \frac{1}{n}\sum_{i=1}^n f_i(x_i), \\
  \st \quad & x_1 = \cdots = x_n, \;\;\\
  & x_i \in \Mcal, \;\; \forall~ i=1,2,\ldots, n,
\end{aligned}
\ee
where $n$ is the number of nodes, $f_i$ is the smooth nonconvex local objective at the $i$-th node, and $\Mcal$ is a (nonconvex) compact smooth submanifold {embedded in} $\R^{d\times r}$ {with the extrinsic dimensions $(d,r)$}, e.g., the Stiefel manifold ${\rm St}(d,r):=\{ x \in \R^{d\times r} : x^\top x = I_r \}$.  
%We consider $n$ nodes, such as machines or edge devices, communicating over a decentralized network described by a undirected graph $\mathcal{G} = (\mathcal{V},\mathcal{E})$, where $\mathcal{V} = \{1,\cdots,n\}$ is the set of node indices and $\mathcal{E} \subset \mathcal{V} \times \mathcal{V}$ is the collection of pairs $(i,j), i,j\in \mathcal{V}$. 
Problem \eqref{prob:original} is prevalent in machine learning, signal processing, and deep learning, see, e.g., the principal component analysis \cite{ye2021deepca}, the low-rank matrix completion \cite{boumal2015low,kasai2019riemannian,deng2023decentralized}, the low-dimensional subspace learning \cite{ando2005framework,mishra2019riemannian}, and the deep neural networks with batch normalization \cite{cho2017riemannian,hu2022riemannian}. One challenge in solving \eqref{prob:original} comes from the nonconvexity of the manifold constraint, which causes difficulty in achieving the consensus \cite{chen2021decentralized,deng2023decentralized,hu2023decentralized,deng2023decentralized11}. 


In this paper, we investigate an online setting where each node \(i\) interacts with its local cost function \(f_i\) through a stochastic first-order oracle (SFO).  This SFO setting is particularly relevant in various online learning and expected risk minimization problems, where the noise introduced by the SFO stems from the variability of sampling over streaming data received at each node. A notable example is online principal component analysis \cite{cardot2018online}. Our primary focus is on the oracle complexity, defined as the total number of SFO queries required at each node to compute an \(\epsilon\)-stationary point tuple \(\{x_1, \cdots, x_n\}\), as formalized in Definition \ref{def:station}.


% In this paper,
% we focus on an online setting where each
% node $i$ only accesses $f_i$ by querying a local stochastic first-order oracle (SFO)  that returns a stochastic gradient, i.e., a noisy version of the exact gradient, at the queried point. As a concrete example of practical
% interest, the SFO mechanism applies to many online learning
% and expected risk minimization problems where the noise in SFO  arises from the uncertainty of sampling from the
% underlying streaming data received at each node, e.g., online principal component analysis \cite{cardot2018online}. We are interested in the oracle complexity, i.e., the total number of queries to SFO required
% at each node to find an $\epsilon$-stationary
% point tuple $\{x_1,\cdots,x_n\}$ defined in Definition \ref{def:station}.
% in term of Definition \ref{def:station}.  
% We
% are interested in the case where each local function $f_i$ takes the following expectation form: 
% $$f_i(x) = \mathbb{E}_{\xi\in \mathcal{D}_i}[f_i(x;\xi)]$$
% with a slight abuse of notation for ease of exposition. In such
% a case, nodes locally compute stochastic gradients of $f_i$. 




\subsection{Related works}




% Decentralized optimization in the Euclidean space (i.e., $\Mcal = \R^{d\times r}$) has been extensively studied over the past few decades. The decentralized (sub)-gradient descent (DGD) method \cite{nedic2009distributed,tsitsiklis1986distributed,yuan2016convergence} provides a straightforward approach to combine local gradient update and consensus update.  When the objective functions are nonconvex, there are also extensive studies, see, e.g., \cite{bianchi2012convergence,di2016next,hong2017prox,tatarenko2017non,scutari2019distributed,sun2020improving,wai2017decentralized,zeng2018nonconvex}.  
% However, since problem \eqref{prob:original} includes a manifold $\Mcal$, which is often nonlinear and nonconvex, these works will fail when applied to solve problem \eqref{prob:original}. 

% As a special case of problem \eqref{prob:original}, the Riemannian consensus
% problem (defined in \eqref{prob:consensus}) is well-studied \cite{chen2021local,markdahl2020high,sarlette2009consensus,tron2012riemannian}. The analysis of decentralized gradient algorithms for solving \eqref{prob:original} often relies on the linear convergence of consensus iteration for solving such consensus problem.
Decentralized optimization in Euclidean space (i.e., $\mathcal{M} = \mathbb{R}^{d \times r}$) has been extensively studied over the past few decades (see, e.g., \cite{bianchi2012convergence,shi2015extra,xu2015augmented,qu2017harnessing,di2016next,tatarenko2017non, wai2017decentralized,yuan2018exact,hong2017prox,zeng2018nonconvex, scutari2019distributed, sun2020improving}). However, since problem \eqref{prob:original} involves a manifold $\Mcal$, which is often nonlinear and nonconvex, these works may fail when directly applied to solve problem \eqref{prob:original}. 

% One basic component in decentralized optimization is to achieve fast consensus across local nodes by solving a corresponding consensus problem. 
%  In the Euclidean setting, due to its linear and convex structure, the linear convergence of the consensus can be straightforwardly obtained by applying the (projected) gradient method.
% % As explained above, the nonconvex nature of the problem in \eqref{prob:original} significantly distinguishes it from Euclidean decentralized composite optimization problems, thereby making consensus construction and algorithmic design more challenging. 
% To address the nonconvexity of the manifold $\Mcal$, the authors \cite{tron2012riemannian} introduce a Riemannian consensus based on the geodesic distance on the manifold. Recently, for the case where $\Mcal$ is the Stiefel manifold, a more tractable consensus has been defined in \cite{chen2021local} by using the Euclidean distance.  For general compact submanifolds, \cite{hu2023achieving} prove that the Riemannian gradient descent with a
% unit step size applied to the consensus problem has locally linear convergence.  


%\revise{split the literature review on deterministic and stochastic}   

Perhaps the earliest works for solving \eqref{prob:original} are  \cite{mishra2019riemannian,shah2017distributed}. However, these methods require an asymptotically infinite number of consensus steps for convergence, which limits their practical applicability.
For the case where $\Mcal$ is the Stiefel manifold, \cite{chen2021decentralized} propose a decentralized Riemannian gradient descent method and its gradient-tracking version.   To use a single step of consensus, augmented Lagrangian methods \cite{wang2022decentralized,wang2022variance} are also investigated, where a different stationarity is used.  \cite{sun2024global} propose a decentralized retraction-free gradient tracking algorithm, and show that it exhibits ergodic
$\mathcal{O}(1/K)$ convergence rate.   
However, these studies rely on the orthogonal structure of the Stiefel manifold.  Recently, \cite{deng2023decentralized} used the projection operators instead of retractions and expanded the distributed Riemannian gradient descent algorithm and the gradient tracking version to the compact submanifolds of Euclidean space. 
Moreover, the integration of decentralized manifold optimization with other algorithms has also been proposed, including the conjugate gradient algorithm \cite{chen2023decentralized} and the natural gradient method \cite{hu2023decentralized}. 
Furthermore, \cite{hu2024improving} achieves single-step consensus for the general compact submanifold by carefully elaborating on the smoothness structure and the asymptotic 1-Lipschitz continuity of the projection operator associated with the submanifold geometry.  



Several studies have focused on the finite-sum setting of problem \eqref{prob:original}, where $f_i = \frac{1}{m}\sum_{r = 1}^{m_i} f_{i,r}$.  
\cite{chen2021decentralized} propose a decentralized  Riemannian stochastic gradient descent method. By combining the variable sample size gradient approximation method with the gradient tracking dynamic, 
\cite{zhao2024distributed} propose a distributed Riemannian stochastic optimization algorithm on the Stiefel manifold. Although both methods can also be used in the online setting,  the oracle complexities are $\mathcal{O}(\epsilon^{-2})$, which is not optimal. 
 It is worth noting that the decentralized variance reduced method \cite{wang2022variance} has been studied. However, they need to periodically calculate the full gradient, which is not suitable for the online setting.
% they are explicitly designed for empirical minimization where each local cost $f_i$ is decomposed as a finite sum of component functions, i.e., $f_i = \frac{1}{m}\sum_{r = 1}^m f_{i,r}$; it is therefore unclear whether these algorithms can be adapted to the online SFO setting.



% Another line of work is the variance reduction technique in stochastic optimization algorithms.  Many methods, including
% Riemannian stochastic variance reduced gradient (RSVRG) \cite{zhang2016riemannian,sato2019riemannian}, stochastic recursive gradient (RSRG) \cite{kasai2018riemannian}, stochastic path-integrated differential estimator (RSPIDER) \cite{zhou2019faster,zhang2018r} are generalized from their Euclidean versions.  However, VR
% methods are originally designed for finite-sum optimization
% where full gradient access is possible. Some recent studies
% have extended Riemannian VR to online setting \cite{zhou2019faster,han2020variance}. 
% Nevertheless, these online VR methods still require computing
% a large batch gradient for each epoch. To address this issue, \cite{han2020riemannian} introduces a novel online variance reduction method, inspired by a recently proposed recursive momentum estimator 
% \cite{cutkosky2019momentum}.









\subsection{Contribution} 
In this paper, we propose DPRSRM, a novel online variance-reduced method for decentralized non-convex manifold optimization with stochastic first-order oracles (SFO). 
% In the following, we emphasize the key advantages of DPRSRM compared with the
% existing decentralized stochastic (variance-reduced) manifold optimization algorithms, from both theoretical and practical aspects.




\begin{itemize}


   \item To achieve
fast and robust performance, the DPRSRM algorithm is built upon gradient tracking \cite{chen2021decentralized,deng2023decentralized} and a stochastic gradient momentum estimator \cite{cutkosky2019momentum,han2020riemannian}, which can be viewed as online variance reduction method. The only existing decentralized stochastic variance-reduced manifold optimization algorithm is the VRSGT proposed by \cite{wang2022variance}. Note that VRSGT is a double-loop algorithm that requires very large minibatch sizes. Conversely, the proposed DPRSRM is a single-loop algorithm with $\mathcal{O}(1)$ oracle queries per update. Numerical experiments demonstrate the effectiveness of the proposed methods compared to state-of-the-art ones through eigenvalue problems and low-rank matrix completion.

    
    \item  Our algorithm achieves an oracle complexity of $\mathcal{O}(\epsilon^{-\frac{3}{2}})$.   A comparison of the oracle complexity of DPRSRM with related algorithms is provided in Table \ref{tab:my_label}, where DPRSRM achieves a lower oracle complexity than the existing decentralized stochastic manifold optimization algorithms \cite{deng2023decentralized,chen2021decentralized,wang2022variance,zhao2024distributed}. 
    Moreover, DPRSRM uses a single step of consensus to achieve communication, compared to other project/retraction algorithms \cite{chen2021decentralized,deng2023decentralized,zhao2024distributed} that need $\log_{\sigma_2}(\frac{1}{2\sqrt{n}})$ rounds of consensus, where $\sigma_2$ is the second largest singular value of the communication graph matrix.  
% \item The only existing decentralized stochastic variance-reduced manifold optimization algorithm is the VRSGT proposed by \cite{wang2022variance}. Note that VRSGT is double loop
% algorithms that require very large minibatch sizes. In particular, during each inner loop they execute a fixed
% number of minibatch stochastic gradient type iterations with $\mathcal{O}(\epsilon^{-1})$ oracle queries per update per node, while
% at every outer loop they obtain a stochastic gradient with
% mega minibatch size by $\mathcal{O}(\epsilon^{-2})$ oracle queries at each node. Clearly, querying the oracles exceedingly, i.e., obtaining
% a large amount of samples, at each node and every iteration in online streaming data scenarios substantially jeopardizes
% the actual wall-clock time. This is because the next iteration cannot be performed until all nodes complete the
% sampling process. \revise{Move to introduction} Moreover, the double-loop implementation may incur periodic network synchronizations. These
% issues are especially significant when the working environments
% of the nodes are heterogeneous. Conversely, the
% proposed DPRSRM is a single-loop algorithm with $\mathcal{O}(1)$ oracle queries per update and only requires a large minibatch size with $\mathcal{O}(1)$ oracle queries once in the initialization
% phase, i.e., before the update recursion is executed; see Algorithm \ref{alg:drgta} for details. Numerical experiments demonstrate the effectiveness of the proposed methods compared to state-of-the-art ones through eigenvalue problems and low-rank matrix completion.
\end{itemize}

\begin{table*}
    \centering
    \begin{tabular}{|c|c|c|c|c|c|}
  \hline
        Algorithm  & Manifold types & Communication & Tracking & VR  & Oracle \\ \hline
        % \multicolumn{6}{*}{Deterministic} \\ \hline
           \cite{chen2021decentralized}     &   Stiefel manifold &  multiple  & \ding{55} & \ding{55} & $\mathcal{O}(\epsilon^{-2})$  \\ \hline
         % \cite{wang2022decentralized}     &  Stiefel & single  & full & no & $\mathcal{O}(\epsilon^{-4})$\\\hline
         % \multirow{2}{*}{\cite{deng2023decentralized}}    &   \multirow{2}{*}{compact}  &  \multirow{2}{*}{multiple}  &  \multirow{2}{*}{full} & no & $\mathcal{O}(\epsilon^{-4})$\\
         % &    &   &  & yes & $\mathcal{O}(\epsilon^{-2})$ \\\hline
%\cite{wang2022variance}     &  Stiefel & single  & \ding{55} & \checkmark  & $\mathcal{O}(\epsilon^{-2})$\\\hline
 % \cite{sun2024global}     &  Stiefel & single  & full & yes & $\mathcal{O}(\epsilon^{-4})$\\\hline
 % \cite{chen2023decentralized}     &  Stiefel & multiple  & full & no  & $-$\\\hline
 %\cite{sun2024global}     &  Stiefel & single  & full & yes & $\mathcal{O}(\epsilon^{-4})$\\\hline
 \cite{zhao2024distributed}     &  Stiefel manifold & multiple  & \ding{55} & 
 \ding{55} & $\mathcal{O}(\epsilon^{-2})$\\\hline
         This paper  &  compact submanifold & single  & \checkmark & \checkmark & $\mathcal{O}(\epsilon^{-3/2})$\\ \hline
    \end{tabular}
    \caption{Comparison of the oracle complexity results of Riemannian online decentralized methods.  ``Communication'' means rounds of communications per iteration.  ``Tracking'' denotes the gradient tracking, ``VR'' denotes variance reduction.  We do not list the work in \cite{wang2022variance} since they focus on the finite-sum setting and are not applicable to the online setting. }%\revise{Add a column on whether needs compute local full gradient Periodically, Require accessibility to full gradient?}}
    \label{tab:my_label}
\end{table*}

\subsection{Notation}
For the compact submanifold $\Mcal$ of $\mathbb{R}^{d \times r}$, we always take the Euclidean metric $\iprod{\cdot}{\cdot}$ as the Riemannian metric. We use $\|\cdot \|$ to denote the Euclidean norm. We denote the $n$-fold Cartesian product of $\Mcal$ as $\Mcal^n = \Mcal \times \cdots \times \Mcal$.
For any $x\in \Mcal$, the tangent space and normal space of $\Mcal$ at $x$ are denoted by $T_x\Mcal$ and $N_x\Mcal$, respectively. For a differentiable function $h: \R^{d\times r} \rightarrow \R$, we denote its Euclidean gradient by $\nabla h(x)$ and its Riemannian gradient by $\grad h(x)$.  
For a positive integer $n$, define $[n] =\{1, \dots, n\}$. Let $\mathbf{1}_n\in \mathbb{R}^n$ be a vector where all entries are equal to 1. Define $J := \frac{1}{n}\mathbf{1}_n\mathbf{1}_n^{\top}$.   Unless otherwise explicitly defined, we now provide explanations for all lowercase variables used in this paper. Take $x$ as an example, we denote $x_i$ as a local variable at $i$-th node; $\hat{x} = \frac{1}{n}\sum_{i=1}^nx_i$ is the Euclidean average. Moreover, we use the bold notations $\bx := [x_1^\top, \dots, x_n^\top]^\top \in \R^{nd \times r},~~\hat{\bx} := [\hat{x}^\top, \dots, \hat{x}^\top]^\top \in \R^{nd \times r},$ where $\bx$ denotes the collection of all local variables $x_i$ and $\hat{\bx}$ is $n$ copies of $\hat{x}$. When applied to the iterative process, in $k$-th iteration, we use $x_{i,k}$ to denote a local variable at $i$-th node and $\hat{x}_k = \frac{1}{n}\sum_{i=1}^nx_{k,i}$. Similarly, we also 
denote 
$\bx_k := [x_{1,k}^\top, \dots, x_{n,k}^\top]^\top \in \R^{nd \times r},~~~ \hat{\bx}_k = [\hat{x}_k^\top, \dots, \hat{x}_k^\top]^\top \in \R^{nd \times r}.$ 
Other lowercase variables can also be denoted similarly as $x$. \revise{Whenever the projection is single-valued, we write $\bar{x}:=\Pcal_{\Mcal}(\hat{x})$ and $\bar{\bx}:=\mathbf 1_n\otimes\bar{x}$.} Define the function $f(\bx): = \sum_{i=1}^n f_i(x_i)$. \revise{For a stacked argument, $\grad f(\bx)$ denotes $[\grad f_1(x_1)^\top,\ldots,\grad f_n(x_n)^\top]^\top$; for a single point $z\in\Mcal$, $f(z):=\frac1n\sum_{i=1}^n f_i(z)$.} Let $\bW := W \otimes I_d \in \R^{nd \times nd}$, where $\otimes$ denotes the Kronecker product. 
% \begin{table}
%     \centering
%         \caption{Caption}
%     \begin{tabular}{|c|l|} \hline
%         $x_i$ & The $i$-th local variable \\ \hline
%         $\bar{x}, \hat{x}$ & The Euclidean average and the induced arithmetic mean on $\mathcal{M}$ for $\{x_i\}_{i=1}^n$ \\ \hline
%         $\bx, \bar{\bx}, \hat{\bx}$ &   $\bx := [x_1^\top, \cdots, x_n^\top]^\top, \hat{\bx} := [\hat{x}^\top, \cdots, \hat{x}^\top]^\top, \bar{\bx} := [\bar{x}^\top, \cdots, \bar{x}^\top]^\top$ \\\hline
        
%         $W$ & The mixing matrix associated with the network $G$ \\\hline
%         $\nabla \phi$  &  The objective function of the consensus problem \\\hline
%          & \\\hline
%          & \\\hline
%          & \\\hline
%          & \\\hline
%          & \\\hline
%          & \\\hline
%     \end{tabular}
%     \label{tab:my_label}
% \end{table}



\section{Preliminary}
This section introduces the definition of stationary point for problem \eqref{prob:original}, and gives a key property for compact submanifolds, i.e., proximal smoothness. 
\subsection{Stationary point}
Let $x_1,\cdots,x_n\in \Mcal$ represent the local copies of each node. 
Let $\Pcal_{\Mcal}$ be the orthogonal projection onto $\Mcal$. Note that for $\{x_i\}_{i=1}^n \subset \Mcal$,
\[ {\rm argmin}_{y\in\Mcal}\sum_{i=1}^n \|y - x_i\|^2 = \Pcal_{\Mcal}(\hat{x}). \]

Any element $\bar{x}$ in $\Pcal_{\Mcal}(\hat{x})$ is the induced arithmetic mean of $\{x_i\}_{i=1}^n$ on $\Mcal$ \cite{sarlette2009consensus}.
\revise{Thus, for a single point $z\in\Mcal$, $\grad f(z)=\frac1n\sum_{i=1}^n\grad f_i(z)$.} The $\epsilon$-stationary point of problem \eqref{prob:original} is defined as follows.
\begin{definition}[\cite{deng2023decentralized}] \label{def:station}
The set of points $\{x_1,x_2,\cdots,x_n\} \subset \Mcal$ is called an $\epsilon$-stationary point of \eqref{prob:original} if there exists a $\bar{x} \in \Pcal_{\Mcal}(\hat{x})$ such that
\[\mathbb{E}\left[ \frac{1}{n}\sum_{i=1}^n \| x_i - \bar{x}\|^2 \right] \leq \epsilon \quad {\rm and} \quad \mathbb{E}[\|\grad f(\bar{x})\|^2] \leq \epsilon. \]
\end{definition}
We refer to these two terms as consensus error and optimality error, respectively. 
%In the following development, we always assure that $\hat{x} \in \bar{U}_{\Mcal}({\delta})$. Consequently, $\Pcal_{\Mcal}(\hat{x})$ is a singleton and we have $\bar{x} = \Pcal_{\Mcal}(\hat{x})$. 



\subsection{Proximal smoothness}\label{sec:proximal}
Proximal smoothness is an effective tool for addressing the nonconvex nature of manifold constraints within decentralized optimization settings \cite{deng2023decentralized}. Define the distance from a point \(x \in \mathbb{R}^{d \times r}\) to the manifold \(\mathcal{M}\) and the nearest-point projection of \(x\) onto \(\mathcal{M}\) as $
\text{dist}(x, \mathcal{M}) := \inf_{y \in \mathcal{M}} \|y - x\| \;\; \text{and} \;\; \mathcal{P}_{\mathcal{M}}(x) := \arg\min_{y \in \mathcal{M}} \|y - x\|,$
respectively. For a given number \(R > 0\), the \(R\)-tube around \(\mathcal{M}\) is defined as the set $U_{\mathcal{M}}(R) := \{x : \text{dist}(x, \mathcal{M}) < R\}.$

A closed set \(\mathcal{M}\) is said to be \(R\)-proximally smooth if the projection \(\mathcal{P}_{\mathcal{M}}(x)\) is unique whenever \(\text{dist}(x, \mathcal{M}) < R\). Following \cite{clarke1995proximal},  an $R$-proximally smooth set $\mathcal{M}$ satisfies that for any real ${\delta}\in (0,R)$,
    \be \label{proximally-1}
    \left\| \Pcal_{\mathcal{M}} (x) -\Pcal_{\mathcal{M}} (y)\right\| \leq \frac{R}{R-{\delta}}\|x - y\|,~~ \forall x,y \in \bar{U}_{\mathcal{M}}({\delta}),
    \ee
where \(\bar{U}_{\mathcal{M}}(\delta) := \{x : \text{dist}(x, \mathcal{M}) \leq \delta\}\). For instance, the Stiefel manifold is known to be a 1-proximally smooth set \cite{balashov2021gradient}, {which also provides additional examples of rank manifolds along with their specific proximal smoothness radii.}


% \begin{itemize}
%     \item[(i)] For any real ${\delta}\in (0,R)$, the estimate holds:
%     \be \label{proximally-1}
%     \left\| \Pcal_{\mathcal{M}} (x) -\Pcal_{\mathcal{M}} (y)\right\| \leq \frac{R}{R-{\delta}}\|x - y\|,~~ \forall x,y \in \bar{U}_{\mathcal{M}}({\delta}),
%     \ee
%     where $\bar{U}_{\Mcal}({\delta}):=\{x: {\rm dist}(x, \Mcal) \leq {\delta} \}$.
%     \item[(ii)] For any point $x\in \mathcal{M}$ and a normal $v\in N_x\Mcal$, the following inequality holds for all $y\in\Mcal$:
%     \be \label{proximally-2}
%     \iprod{v}{y-x} \leq \frac{\|v\|}{2R}\|y-x\|^2.
%     \ee
% \end{itemize}

 
% In the design of manifold optimization algorithms \cite{absil2009optimization,boumal2023introduction,hu2020brief,huang2017intrinsic}, a key concept is the retraction operator. A smooth mapping $\Rcal:T\mathcal{M}:=\cup_{x \in \Mcal} T_{x} \Mcal \rightarrow \Mcal$ is called a retraction operator if  
% \begin{itemize}
% 	\item $\Rcal_{x}(0) = x$,
% 	\item $\mathrm{D}\Rcal_{x}(0)[\xi]:=\frac{\mathrm{d}}{\mathrm{d}{\tau}} \Rcal_{x}({\tau}\xi) \mid_{{{\tau}}=0} = \xi$, for all $\xi \in T_{x}\Mcal$. 
% \end{itemize}
% Note that the retraction operator may not be unique, e.g., the exponential maps \cite{absil2009optimization}, (certain types of) projections \cite{absil2012projection}.   
% The following Lipschitz-type inequality of the retraction operator plays an important role in the analysis of retraction-based methods.
% \begin{proposition}[\cite{boumal2019global}] \label{eq:lip-retr}
% 	Let $\mathcal{M}$ be a compact  submanifold of $\mathbb{R}^{d \times r}$. For all $x \in \mathcal{M}$ and $\xi \in T_{x} \mathcal{M}$, there exists a constant $M_1 >0$ such that the following inequality holds:
% 	\be \label{eq:bound-retraction1} \|\Rcal_{x}(\xi) - x - \xi \| \leq M_1 \|\xi \|^2, \; \forall x
% 1	\in \Mcal, \; \forall \xi \in T_{x} \Mcal. \ee
% \end{proposition} 



\section{Problem setup and the proposed DPRSRM}\label{sec:main-alg}
In this section, we present the problem setup considered in this paper, outlining the assumptions for the objective function and the communication network. Building on this setup, we then develop a decentralized algorithm to solve problem \eqref{prob:original} and provide the main convergence rate results. 
\subsection{Problem setup}
Let us start with {assumptions on problem \eqref{prob:original}, based on the fact that any compact $C^2$-submanifolds in Euclidean space belong to proximally smooth set \cite{balashov2021gradient,clarke1995proximal,davis2020stochastic}.}  
% Without loss of generality, we assume that $\Mcal$ is $R$-proximally smooth. For the objective function $f_i$, we make the following assumption.
\begin{assumption}\label{assum-f}
    {Assume that $\Mcal$ is proximally smooth with radius $R$.} Each objective function $f_i$ is gradient Lipschitz continuous with modulus $L_{f}$ on the convex hull of $\Mcal$, denoted by ${\rm conv}(\Mcal)$. \revise{Moreover, the aggregate objective $f(x):=\frac1n\sum_{i=1}^n f_i(x)$ is bounded below on $\Mcal$; denote its minimum by $f_*$.}  
%Moreover, the Euclidean gradient is bounded by $L_G$, i.e., $ \max_{x \in \Mcal} \|  \nabla f_i(x)\| \leq L_G, \; i\in [n].$
\end{assumption}
Under this assumption, the following Riemannian quadratic upper bound for $f_i$ has been estabilshed in \cite{deng2023decentralized}. 
\begin{lemma}[\cite{deng2023decentralized}, Lemma 2]\label{lemma:lipsctz}
Under Assumption \ref{assum-f}, there exists $L_g$, for any $x,y\in \Mcal$, the following inequality holds:
\begin{equation}\label{f-riemannian-Lip}
    f_i(y) - f_i(x) \leq \left<\grad f_i(x),y-x\right>+\frac{L_g}{2}\|y-x\|^2, ~ i\in [n].
\end{equation}   
Moreover, there exists a constant $L_G>0$ such that
\begin{equation}\label{g-riemannian-lip}
    \| \grad f_i(x) - \grad f_i(y) \| \leq L_G \|x - y\|,~i\in [n].
\end{equation}
\end{lemma}

We now present the assumption for the communication network. 
Denote by the undirected graph $G:=\{ \mathcal{V}, \mathcal{E}\}$, where $\mathcal{V}=\{1,2,\ldots, n\}$ is the set of all nodes and $\mathcal{E}$ is the set of edges. {Let $W$ be the mixing matrix associated with $G$ and $W_{ij} \ne 0$ if there is an edge $(i,j) \in \mathcal{E}$ and $W_{ij} = 0$ otherwise.} We use the following standard assumptions on $W$, see, e.g., \cite{chen2021decentralized,zeng2018nonconvex}.
% Let $W$ be the adjacency matrix of $G$. Then $W_{ij} = W_{ji}$ and $W_{ij} > 0$ if an edge $(i,j) \in \mathcal{E}$ and otherwise $W_{ij} = 0$.
\begin{assumption} \label{assum-w}
    We assume that the undirected graph $G$ is connected and $W$ is doubly stochastic, i.e., (i) $W=W^{\top}$; (ii) $W_{i j} \geq 0$ and $1>W_{i i}>0$ for all $i, j$; (iii) the eigenvalues of $W$ lie in $(-1,1]$. \revise{In addition, the second-largest singular value $\sigma_2$ of $W$ satisfies $\sigma_2\in[0,1)$.}
\end{assumption}
Consider a sufficiently rich probability space \(\{\Omega, \mathbb{P}, \mathcal{F}\}\). For a given decentralized algorithm, we assume it generates an iterative sequence \(\{x_{i,k}\}_{k \geq 0}\), where \(x_{i,k}\) denotes the \(k\)-th iteration at node \(i\). At each step, node \(i\) observes a random vector \(\xi_{i,k}\). We then define an increasing sequence of sub-\(\sigma\)-algebras within \(\mathcal{F}\), constructed from the random vectors observed in succession by the network nodes: for all $k \geq 1$,
$
\mathcal{F}_0 := \{\Omega, \emptyset\}, \quad \mathcal{F}_k := \sigma\left(\{\xi_{i,0}, \xi_{i,1}, \ldots, \xi_{i,k-1} : i \in \mathcal{V}\}\right),
$
where \(\emptyset\) represents the empty set. The following assumptions are made regarding the stochastic gradient \(\nabla f_i(x, \xi_{i,k})\):
\begin{assumption}\label{assum:stochatis}
\begin{revisionblock}
For any $\mathcal{F}_k$--measurable variable $x\in\Mcal$ and $k\geq 0$, the algorithm generates a sample $\xi_{i,k}\sim \Omega$ for each node $i$ and returns a stochastic gradient $\grad f_i(x,\xi_{i,k})$. There are constants $\nu_i\geq0$ and $\bar L>0$; we write $\nu^2:=\frac{1}{n}\sum_{i=1}^n\nu_i^2$. The oracle satisfies
\begin{align}
  &  \mathbb{E}\left[ \grad f_i(x,\xi_{i,k}) \vert \mathcal{F}_k \right] = \grad f_i(x), \label{assm:mean} \\ 
 &    \mathbb{E}\left[ \|\grad f_i(x,\xi_{i,k}) - \grad f_i(x) \|^2 \vert \mathcal{F}_k\right] \leq \nu_i^2. \label{assm:vari}
\end{align}
Moreover, the collection $\{\xi_{i,k}:i\in \mathcal{V},k\geq0\}$ consists of independent random variables and the same-sample stochastic gradients are mean-square $\bar L$-Lipschitz:
\begin{equation}\label{assm:ms-lipschitz}
   \mathbb{E}\!\left[ \| \grad f_i(x,\xi_{i,k}) - \grad f_i(y,\xi_{i,k}) \|^2 \mid \mathcal{F}_k\right] \leq \bar{L}^2\|x-y\|^2.
\end{equation}
Finally, the un-clipped initialization used in Algorithm~\ref{alg:drgta} satisfies $\|\grad f_i(x_{i,0},\xi_{i,0})\|\leq B$ almost surely.
\end{revisionblock}
\end{assumption}




 


To measure the oracle complexity,  we give the definition of a stochastic first-order oracle (SFO) for \eqref{prob:original}.
\begin{definition}[\textbf{stochastic first-order oracle}]\label{def:stochastic}
    For the problem \eqref{prob:original}, a stochastic first-order oracle for each node $i$ is defined as follows: compute the Riemannian gradient $\grad f_i(x,\xi_i)$ given a sample $\xi_i\in \Omega$.  
\end{definition}

 













 





  
\subsection{The Algorithm}




In this subsection, we introduce the Decentralized Projected Riemannian Stochastic Recursive Momentum (DPRSRM) method for addressing \eqref{prob:original}, along with its associated convergence results. Inspired by the notable effectiveness of variance reduction and gradient tracking in decentralized frameworks \cite{wang2022variance,deng2023decentralized,chen2021decentralized}, we seek a novel combination of variance reduction and gradient tracking for decentralized online problems on compact submanifolds for improving the oracle complexity. In particular, we focus on the following stochastic gradient estimator using the momentum technique introduced in \cite{cutkosky2019momentum,han2020riemannian,deng2024oracle}:
\begin{equation}\label{update:qik}
\begin{aligned}
    q_{i,k} = &   \grad f_{i}(x_{i,k},\xi_{i,k}) \\
   & +(1-\tau)(d_{i,k-1} - \grad f_{i}(x_{i,k-1},\xi_{i,k})).
\end{aligned}
\end{equation}

We note that \eqref{update:qik} can be rewritten as 
\begin{equation}\label{update:qik1}
\begin{aligned}
    q_{i,k} = &  \tau \grad f_{i}(x_{i,k},\xi_{i,k})  +(1-\tau)(d_{i,k-1} \\
    &+\grad f_{i}(x_{i,k},\xi_{i,k})- \grad f_{i}(x_{i,k-1},\xi_{i,k})),
\end{aligned}
\end{equation}
which hybrids stochastic gradient $\grad f_{i}(x_{i,k},\xi_{i,k})$ with the recursive gradient estimator in RSARAH/SRG \cite{han2020variance,zhou2019faster} for $\tau \in[0,1]$.  Since the direction $q_{i,k}$ may be unbounded, we introduce a clipped gradient estimator $d_{i,k}$
\begin{equation}\label{update:dik}
    d_{i,k} = \left\{
    \begin{array}{cc}
        q_{i,k} & \text{if} ~ \| q_{i,k} \| \leq B, \\
        B \frac{q_{i,k}}{\|q_{i,k}\|} & \text{otherwise},    \end{array}
    \right.
\end{equation}
where $B>0$ is a user-defined constant.   To further reduce the variance of the gradient estimator in different nodes, we compute the gradient tracking iteration as follows:
\be\label{update:sik}
s_{i,k}  = \sum_{j=1}^n W_{ij}s_{j,k-1} + d_{i,k} - d_{i,k-1}, ~i\in [n].
\ee

 A crucial advantage of gradient tracking-type methods lies in the applicability of the use of a constant step size $\alpha$. Since $s_{i,k}$ may not remain in the tangent space $T_{x_{i,k}}\Mcal$, we introduce {its tangent space projection} $v_{i,k} = \Pcal_{T_{x_{i,k}} (s_{i,k})}$ and update the new iterate $x_{i,k+1}$ as follows:
 \begin{equation}\label{updata-new-}
     x_{i,k+1} = \Pcal_{\Mcal}(\sum_{j=1}^n W_{ij}x_{j,k} - \alpha
v_{i,k} ), ~i\in [n],
 \end{equation}
 {where the projection-based average \cite{deng2023decentralized,hu2024improving}, i.e., $\Pcal_{\Mcal}(\sum_{j=1}^n W_{ij}x_{j,k})$, is adopted to ensure the decrease in the consensus error among nodes.}
For ease of analysis, we stack variables in each node and rewrite \eqref{update:qik}, \eqref{update:sik} and \eqref{updata-new-} as 
\begin{equation}
\left\{
    \begin{aligned}
        \bq_k  =& \grad f(\bx_k,\mathbf{\xi}_k) \\
        &+ (1-\tau) (\bd_{k-1} - \grad f(\bx_{k-1},\mathbf{\xi}_k)) \\
        \bs_{k} =& \bW\bs_{k-1} + \bd_k - \bd_{k-1} \\
        \bv_k  = &\Pcal_{T_{\bx}\Mcal^n}( \bs_k ) \\
        \bx_{k+1}  =&  \Pcal_{\Mcal^n}(\bW \bx_k - \alpha \bv_k),
    \end{aligned}
    \right.
\end{equation}
where  $\Pcal_{T_{\bx}\Mcal^n}:= \Pcal_{T_{x_{1,k}}\Mcal} \times \dots \times \Pcal_{T_{x_{n,k}}\Mcal}$ and $\mathbf{\xi}_k = \{\xi_{i,k}\}_{i\in \mathcal{V}}$. The overall algorithm is given in Algorithm \ref{alg:drgta}.






    


\begin{algorithm}[htbp]
\caption{The DPRSRM for solving \eqref{prob:original}} \label{alg:drgta}
\begin{algorithmic}[1]
\REQUIRE  \revise{Initial point $\bx_{0} = \bar{\bx}_0 \in \Mcal^n$, $\bs_{-1} =  \bd_{-1} = \mathbf{0}$, $\alpha,\tau$, and clipping radius $B$.}
\STATE Sample $\xi_{i,0}$,  let $s_{i,0} = d_{i,0} = \grad f_i(x_{i,0},\xi_{i,0})$.
\STATE $v_{i,0} = \Pcal_{T_{x_{i,0}}\Mcal}( s_{i,0})$.
\STATE $ x_{i,1} = \Pcal_{\Mcal}(\sum_{j=1}^n W_{ij}x_{j,0} - \alpha
v_{i,0} )$.
\FOR{$k = 1,2,\cdots,K$}
\STATE Update stochastic gradient estimator $q_{i,k}$ via \eqref{update:qik}
\STATE Update the clipped gradient estimator $d_{i,k}$ via \eqref{update:dik}.
\STATE Update Riemannian gradient tracking $s_{i,k}$ via \eqref{update:sik}.
\STATE Project onto tangent space: $v_{i,k} = \Pcal_{T_{x_{i,k}}\Mcal}( s_{i,k})$.
\STATE Update new iterate $x_{i,k+1}$ via \eqref{updata-new-}. 
%\STATE Set $k=k+1$.
\ENDFOR
\end{algorithmic}
\end{algorithm}
   
    
%\end{remark}

%\subsubsection{DPRSRM}
%Next, we investigate a gradient tracking method for solving \eqref{prob:original} by leveraging the gradient tracking techniques introduced in \cite{chen2021decentralized,nedic2017achieving,qu2017harnessing}  to get a better estimate for the full gradient. 


\subsection{Main results}\label{sec:main-resilt}
   

The convergence analysis of Algorithm \ref{alg:drgta} can be divided into two parts: the consensus error and the optimality error, as defined in Definition  \ref{def:station}.  Let us first focus on the consensus error.  We define a neighborhood around $\bx\in \Mcal^n$ as follows:
\begin{equation}\label{def:N-neiborhood-1}
    \mathcal{N}(\delta):=\{ \bx\in \Mcal^n: \|\bx - \bar{\bx} \| \leq \delta \}.
\end{equation}
Note that if $\bx\in \mathcal{N}(\delta)$, then $x_{i}\in \bar{U}_{\Mcal}(\delta)$ for any $i\in [n]$. %We also give the conditions for $\alpha$ and $\delta$:

\revise{By appropriately selecting the step size $\alpha$ and initializing with $\bx_0 \in \mathcal{N}(\delta)$, we can establish an explicit relationship between the consensus error and the step size. The proof is given in Section \ref{sec:linear-consensus}.}
% The proof follows from \cite[Lemma D.3]{chen2021decentralized}.
\begin{revisionblock}
Define
\begin{equation}\label{def:repaired-constants}
 L:=\max\{L_g,L_G,\bar L\},\qquad
 G_{\max}:=\max_{i\in[n],\,x\in\Mcal}\|\grad f_i(x)\|,
\end{equation}
where $G_{\max}<\infty$ follows from compactness. Throughout the convergence
analysis we take $B\geq G_{\max}$ and set
\[
 D_B:=\left(\frac{3B}{1-\sigma_2}\right)^2,
 \qquad \kappa_M:=\frac{R}{R-2\delta},
 \qquad \rho:=\kappa_M\sigma_2.
\]
\begin{theorem}[Consensus error] \label{lem:consensus}
Let $\{\bx_k\}_k$ be generated by Algorithm~\ref{alg:drgta} and suppose that
Assumptions~\ref{assum-f}--\ref{assum:stochatis} hold. Let $B\geq G_{\max}$
and suppose that
   \begin{equation}\label{cond-alpha-delta}
\begin{aligned}
    \alpha & \leq  \min\left\{ \frac{\delta}{\sqrt{nD_B}},
      \frac{(R(1-\sigma_2)-2\delta)\delta}{R\sqrt{nD_B}} \right\}, \\
    \delta &<  \frac{1}{2}\min\left\{ R,  R(1-\sigma_2)\right\}.
    \end{aligned}
\end{equation}
If $\bx_0=\bar\bx_0$, then every projection in Algorithm~\ref{alg:drgta} is
single-valued in the prescribed tube and, for all $k\geq0$,
    \begin{equation}\label{eq:xk-barxk-bound}
    \frac{1}{n} \|\bar{\bx}_k - \bx_k \|^2 \leq C_B\alpha^2,
    \qquad C_B:=\frac{\kappa_M^2D_B}{(1-\rho)^2}.
    \end{equation}
\end{theorem}
   
We now state the repaired optimality result. Its proof is given in
Section~\ref{sec:proof:dprgt}.
\begin{theorem}[Optimality error] \label{thm:dprgt}
Let the hypotheses of Theorem~\ref{lem:consensus} hold, with $\alpha$ chosen
below. There exists a
problem-dependent constant $\kappa_0>0$. Choose a fixed
$\kappa\geq\kappa_0$ and then a sufficiently small problem-dependent constant
$\theta>0$. There are constants $K_0,C_{\rm opt}>0$, independent of the
horizon $K$, such that, for every $K\geq K_0$, setting
\begin{equation}\label{eq:repaired-parameters}
 \alpha=\theta(K+1)^{-1/3},\qquad \tau=\kappa\alpha^2\leq1,
\end{equation}
with $\alpha$ also satisfying \eqref{cond-alpha-delta}. Then
\begin{equation}\label{eq:repaired-optimality-rate}
 \min_{0\leq k\leq K}\mathbb{E}\|\grad f(\bar x_k)\|^2
 \leq C_{\rm opt}(K+1)^{-2/3}.
\end{equation}
\end{theorem}




\begin{corollary}\label{coll}
Under the conditions of Theorem~\ref{thm:dprgt}, there is an index
$k_*\leq K$ such that
\[
 \mathbb{E}\|\grad f(\bar x_{k_*})\|^2=\mathcal O(K^{-2/3}),
 \qquad
 \frac1n\mathbb{E}\|\bx_{k_*}-\bar\bx_{k_*}\|^2
 =\mathcal O(K^{-2/3}).
\]
Since each node evaluates only a constant number of stochastic gradients per
iteration, DPRSRM has per-node SFO complexity
$\mathcal O(\epsilon^{-3/2})$ for Definition~\ref{def:station}.
\end{corollary}
\end{revisionblock}








According to Corollary \ref{coll}, DPRSRM achieves an oracle complexity of  $\mathcal{O}(\epsilon^{-\frac{3}{2}})$, outperforming existing methods \cite{chen2021decentralized,zhao2024distributed}, which have an oracle complexity of at most \(\mathcal{O}(\epsilon^{-2})\).

% {add comparison with [12]: target more general compact submanifold, develop new tools: linear consensus of PGD, although some results can be inferred from the Stiefel manifold case, it needs careful elaboration of the proximal smoothness and the essential inequalities of the manifold, which are nontrivial.
% }


\section{Outline of convergence analysis}\label{sec4}
As shown in Section \ref{sec:main-resilt},  the convergence analysis consists of two critical components: the consensus error and the optimality measure, corresponding to Theorem \ref{lem:consensus} and Theorem \ref{thm:dprgt}, respectively. In this section, we will outline the proof. 
% To ensure that the projection operator in Algorithm \ref{alg:drgta} is well-defined, we need to demonstrate that the points being projected always remain within a certain neighborhood. This is due to our assumption that $\Mcal$ is an $R$-proximally smooth set. Based on this result, we provide the iteration complexity result for Algorithm \ref{alg:drgta}.
 \subsection{Consensus error}  \label{sec:linear-consensus}
% To achieve the stationarity given in Definition \ref{def:station}, it is necessary to consider consensus on the compact submanifold $\Mcal$. 

% Specifically, we consider the following consensus problem over $\Mcal$: 
% \be \label{prob:consensus}
% \min_{\bx} \phi(\bx): = \frac{1}{4} \sum_{i=1}^n \sum_{j=1}^n W_{ij}\|x_i - x_j\|^2,~ \text{s.t.}~x_i \in \Mcal, i\in [n],
% \ee
% {where $t$ is a given integer and $W_{ij}$ denotes the $(i,j)$-th element of $W$ (the $t$-th power of $W$).}
% % A large $t$ means that the graph associated with $W$ is more connectivity, which is crucial for the later convergence analysis of our algorithm.
% The Euclidean gradient of $\phi(\bx)$ is given as follows:
% $$\nabla \phi(\bx): = [\nabla \phi_1(\bx)^T, \nabla \phi_2(\bx)^T, \cdots, \nabla \phi_n(\bx)^T ]^T = (I_{nd} - \bW) \bx,$$ where $\nabla \phi_i(\bx): = x_i - \sum_{j=1}^n W_{ij}x_{j}, i \in [n]$.   

% {Note that if the graph is connected, a matrix $W$ that satisfies the above assumption can be constructed through various approaches \cite{shi2015extra}. For any such $W$, $W$ will also satisfy Assumption \ref{assum-w} for any integer $t \geq 1$. Furthermore, as $t$ increases, $W$ becomes denser, indicating an increase in the connectivity of the corresponding graph.} 

% Let us denote $\bar{\bx}_k = \mathbf{1}_n \otimes \bar{x}_k$ and $\hat{\bx}_k = \mathbf{1}_n \otimes \hat{x}_k$. For the Euclidean case (i.e., $\Mcal = \R^{d\times r}$), the gradient descent method with a unit step size has a locally linear convergence rate, in which the $k$-th iteration is given by $\bx_{k+1} = \bx_k - \nabla \phi(\bx_k)$. In fact, we have
% \[ \begin{aligned}
%     \| \bx_{k + 1} - \bar{\bx}_{k+1} \| & \leq \| \bx_{k+1} - \bar{\bx}_k \| = \| \bx_k - (I_{nd} - \bW)\bx_k - \bar{\bx}_k\| \\
%     & = \| (\bW - J)(\bx_k - \bar{\bx}_k) \| \leq \sigma_2 \|\bx_k - \bar{\bx}_k\|. 
% \end{aligned} \]
% When $\Mcal$ is the Stiefel manifold, the iterates generated by the Riemannian gradient descent method converge linearly under suitable step size and initialization, as shown in \cite[Theorem 2]{chen2021local}. However, it is not clear whether such linear convergence result holds for general compact submanifolds. Note that the linear convergence of (Riemannian) gradient methods is crucial in analyzing the convergence of decentralized gradient-type methods \cite{chen2021decentralized,sun2020improving}. This raises the question of whether we can design a gradient-type algorithm to solve \eqref{prob:consensus} with a locally linear convergence rate.


%The main contribution of this section is to show the locally linear convergence of the scheme \eqref{eq:consensus-pg-iter} with an explicit characterization of the local neighborhood. The main technique utilized is the proximal smoothness of the compact submanifold $\Mcal$ \cite{clarke1995proximal}. 


% Without loss of generality, we assume that $\Mcal$ is $2{\delta}$-proximally smooth. By \eqref{eq:lip-proj} and \eqref{proximally-0}, the projection operator $\Pcal_{\Mcal}(\cdot)$ has the following properties:
%   \begin{align}
% \left\| \Pcal_{\mathcal{M}} (x) -\Pcal_{\mathcal{M}} (y)\right\| & \leq 2\|x - y\|,~~ \forall x,y \in \bar{U}_{\Mcal}({\delta}), \label{proximally-1} \\
%  \left<v, y-x\right> &\leq \frac{\|v\|}{4{\delta}}\|y-x\|^2,~ \forall x,y\in \Mcal, v\in N_x\Mcal. \label{proximally-2}
% \end{align}
% We will use the above inequality to characterize the locally linear convergence of the projected gradient method and the associated neighborhood.

This subsection addresses the consensus error in Theorem \ref{lem:consensus}. 
In Algorithm \ref{alg:drgta}, the update of the main iterate $x_{k+1}$ involves the projection on $\Mcal$. Due to the nonconvex nature of compact submanifolds, the projection is not always unique. Before proving Theorem \ref{lem:consensus}, we will first demonstrate that the projection operator in Algorithm \ref{alg:drgta} is well-defined, meaning that the points being projected are always ensured to be within a neighborhood that belongs to  $\bar{U}_{\Mcal}(R)$.  %This implies that the projection is unique and possesses Lipschitz continuity.

\begin{revisionblock}
We first investigate the uniform boundedness of $\|\bs_k\|$ in the following lemma.
\begin{lemma}\label{lemma:neibohood}
Let $\{\bx_k\}_k$ be generated by Algorithm~\ref{alg:drgta}. Under
Assumptions~\ref{assum-w}--\ref{assum:stochatis}, for every $k\geq0$,
\be\label{eq:sk-bound}
 \|\bs_k\|^2 \leq nD_B.
\ee
\end{lemma}

 The following lemma shows that the iterates $\{\bx_k\}$ remain in
 $\mathcal{N}(\delta)$ under the stated conditions.

\begin{lemma} \label{lem:stay-neighborhood}

Suppose that Assumptions~\ref{assum-f}--\ref{assum:stochatis} hold and
$B\geq G_{\max}$. Let $x_{i,k}$ be generated by Algorithm~\ref{alg:drgta},
and let $\alpha$ and $\delta$ satisfy \eqref{cond-alpha-delta}.
If $\bx_0=\bar\bx_0$, then for every $k\geq0$,
$\bx_k\in\mathcal N(\delta)$ and
\begin{align}
    \sum_{j=1}^n {W_{ij}} x_{j,k} -\alpha v_{i,k} & \in \bar{U}_{\Mcal}(2\delta), ~ i=1,\cdots,n.
\end{align}
\end{lemma}


Now we give the proof of Theorem \ref{lem:consensus}. 

\begin{proof}[Proof of Theorem \ref{lem:consensus}]
Since $\bx_0=\bar\bx_0$, Lemma~\ref{lem:stay-neighborhood} gives, for every
$k\geq0$, $\bx_k\in\mathcal N(\delta)$ and
\begin{equation}\label{eq:neibohood3}
    \begin{aligned}
       &  \sum_{j=1}^n {W_{ij}} x_{j,k} -\alpha v_{i,k}  \in \bar{U}_{\Mcal} (2\delta), ~ i = 1,\cdots,n.
    \end{aligned}
\end{equation}
By the definition of $\bar{\bx}_{k+1}$, we have
\[ \begin{aligned}
    \|\bx_{k+1} - \bar{\bx}_{k+1}\| & \leq \| \bx_{k+1} - \bar{\bx}_k \| \\
    & = \| \Pcal_{\Mcal^n}(\bW \bx_k - \alpha  \bv_k) - \Pcal_{\Mcal^n}( \hat{\bx}_k) \| \\
    & \leq \frac{R}{R-2\delta} \|\bW \bx_k - \alpha  \bv_k -\hat{\bx}_k  \| \\
    & \leq \rho\|\bx_k-\bar{\bx}_k\|
      +\kappa_M\sqrt{nD_B}\,\alpha,
\end{aligned}
\]
where the first inequality follows from the optimality of $\bar{\bx}_{k+1}$,
and the second uses \eqref{eq:neibohood3} and the
$\kappa_M$-Lipschitz continuity of $\Pcal_{\Mcal}$ over
$\bar U_{\Mcal}(2\delta)$. Hence
 \be \label{eq:temp:pgd:consensus}
    \begin{aligned}
        \|\bar{\bx}_{k+1}-\bx_{k+1}\|
        &\leq \rho^{k+1}\|\bar{\bx}_0-\bx_0\|
        +\frac{\kappa_M\sqrt{nD_B}}{1-\rho}\alpha.
    \end{aligned}
    \ee   
The initialization $\bx_0=\bar\bx_0$ proves \eqref{eq:xk-barxk-bound}.
\end{proof}
\end{revisionblock}



% Lemma \ref{lem:stay-neighborhood} is established based on a general iterative scheme with bounded $v_{i,k}$.  

% We consider the projected gradient method with a unit step size for solving \eqref{prob:consensus}, namely,
% \be \label{eq:consensus-pg-iter}
% x_{i,k+1} = \Pcal_{\Mcal}(x_{i,k} - \nabla \phi_i(\bx_k))  = \Pcal_{\Mcal}\left(\sum_{j=1}^n W_{ij}x_{j,k} \right),\quad i\in [n].
% \ee
% Let us denote $\Pcal_{\Mcal^n}(\bx) = [\Pcal_{\Mcal}(x_1)^\top, \ldots, \Pcal_{\Mcal}(x_n)^\top]^\top$. Based on Lemma \ref{lem:stay-neighborhood}, we show that the projected gradient method \eqref{eq:consensus-pg-iter} converges to the consensus set at a locally linear rate. 
% \begin{theorem}[Linear convergence of consensus error] \label{theo:linear-con-consensus}
%    Let $\{\bx_k\}$ be the iterate sequences generated by \eqref{eq:consensus-pg-iter}.  Suppose that Assumption \ref{assum-w} holds. If $\bx_{0} \in \mathcal{N}$,  $t > \max\left\{ \left\lceil \log_{\sigma_2}\left(\frac{{\delta}}{24\sqrt{n}\zeta}\right) \right \rceil, \left\lceil \log_{\sigma_2}(1/2) \right \rceil \right\}$, then, $\bx_k \in \mathcal{N}$ for all $k\geq 0 $ and the following linear convergence with  rate $\rho_t: = 2\sigma_2 < 1$ holds,
%    \bee
%     \|\bx_{k+1} - \bar{\bx}_{k+1}\| \leq  \rho_t \|\bx_k -\bar{\bx}_k\|.
%    \eee
   
% \end{theorem}

% \begin{proof}
% Since $\|\hat{x}_0 - \bar{x}_0\| \leq \frac{1}{2}{\delta}$, by invoking Lemma \ref{lem:stay-neighborhood} with $\alpha = 0, v_{i,k} = 0, i\in [n]$, it follows that for any $k>0$, $\bx_k \in \mathcal{N}$ and
% \begin{equation}\nonumber
%     \begin{aligned}
%        &  \sum_{j=1}^n {W_{ij}} x_{j,k}  \in \bar{U}_{\Mcal} ({\delta}), ~ i\in [n].
%     \end{aligned}
% \end{equation}
%      Based on the iterative scheme in \eqref{eq:consensus-pg-iter}, we have
% \be\label{eq:wx-hatx} \begin{aligned}
%     \|\bx_{k+1} - \bar{\bx}_{k+1}\| & \leq \| \bx_{k+1} - \bar{\bx}_k \|  = \| \Pcal_{\Mcal^n}(\bW \bx_k) - \Pcal_{\Mcal^n}( \hat{\bx}_k) \| \\
%     & \leq 2 \|\bW \bx_k -\hat{\bx}_k  \|  = 2 \| (W \otimes I_d) \bx_k - \hat{\bx}_k \| \\
%     & = 2 \| ((W -J ) \otimes I_d) (\bx_k - \hat{\bx}_k) \|  \\
%     % & \leq 2 \| \bW \bx_k - \hat{\bx}_k \| + \alpha  L \\
%     & \leq 2 \sigma_2 \|\bx_k -\hat{\bx}_k\| \leq  2 \sigma_2 \|\bx_k -\bar{\bx}_k\|,
% \end{aligned}
% \ee
% where the first inequality utilizes \eqref{proximally-1}.  The proof is completed.
% \end{proof}

% The result above is also applicable to the case when $\Mcal$ is a convex set. In this scenario, by utilizing the 1-Lipschitz continuity of the projection operator, the linear rate can be improved to $\sigma_2$, which is consistent with the results in \cite{nedic2018network,nedic2010constrained}. It is worth noting that the rate in Theorem \ref{theo:linear-con-consensus} will asymptotically go to $\sigma_2$ as ${\Pcal}_{\Mcal}$ is approximately $1$-Lipschitz continuously near $\Mcal$ (i.e., the second inequality in \eqref{eq:wx-hatx} can be tighter when $\bW \bx_k$ and $\hat{\bx}_k$ are closer to $\Mcal^n$). {Unlike the results for the Stiefel manifold \cite{chen2021local}, which focuses on the Riemannian gradient descent, we show the linear convergence of projected gradient descent method for general compact submanifolds. It is worth noting that extending the analysis in \cite{chen2021local} to general submanifolds is challenging due to the absence of Stiefel manifold-specific geometries, such as the Riemannian gradient and projections to tangent spaces.}
% we establish the locally linear convergence of the projected gradient descent. Furthermore, our findings allow for the use of unit step size, which may not be acceptable in their case.}
% with a rate of $\sigma_2$ (where $t$ can be any positive integer) \cite{nedic2010constrained,nedic2018network} can be obtained. 
% Furthermore, our findings allow for the use of unit step size, which may not be acceptable in their case.




%\subsection{Two key lemmas}

% \begin{proof}
% Since $\| \nabla \phi(\bx_k)\|  = \|(I_{nd} - \bW) \bx_k \| = \|(I_{nd} - \bW) (\bx_k - \bar{\bx}_k) \| \leq 2 \|\bx_k - \bar{\bx}_k \|$, we have
% \be
% \begin{aligned}
%     & \| \hat{x}_{k+1} - \hat{x}_k \| \\
%     \leq & \| \hat{x}_{k+1} - \hat{x}_k + \frac{1}{n} \sum_{i=1}^n (\grad \phi_i(\bx_k) + \alpha v_{i,k}  ) \| + \|\frac{1}{n} \sum_{i=1}^n (\grad \phi_i(\bx_k) + \alpha v_{i,k}  )\| \\
%     \overset{\eqref{projec-second-order1} }{\leq} & \frac{Q}{n}\sum_{i=1}^n \| \nabla \phi_i(\bx_k) + \alpha v_{i,k} \|^2 +  \| \frac{1}{n} \sum_{i=1}^n\grad \phi_i(\bx_k)\|  + \alpha \| \hat{u}_{k}  \| \\
%     \leq &\frac{2Q}{n}\| \nabla \phi(\bx_k)\|^2 + \frac{2Q\alpha^2}{n} \| \bu_k \|^2 +  \| \frac{1}{n} \sum_{i=1}^n\grad \phi_i(\bx_k)\|  + \alpha \| \hat{u}_{k}  \| \\
%    \overset{\eqref{eq:sum-consen-grad}}{ \leq} & \frac{8Q + \sqrt{n}L_2}{n} \|\bx_k - \bar{\bx}_k\|^2 +\frac{2Q\alpha^2}{n} \| \bu_k \|^2  + \alpha \| \hat{u}_{k}  \|.
% \end{aligned}
% \ee
% Therefore, we have
% \begin{equation}
%     \begin{aligned}
%  &   \| \bar{x}_{k+1} - \bar{x}_k \|  \leq \| \hat{x}_{k+1} - \hat{x}_k  \| + \| \hat{x}_{k+1} - \bar{x}_{k+1}  \| + \| \hat{x}_{k} - \bar{x}_{k}  \| \\
%     \overset{\eqref{eq:distance-an-rm}}{ \leq} &   \frac{8Q + \sqrt{n}L_2}{n} \|\bx_k - \bar{\bx}_k\|^2 +\frac{2Q\alpha^2}{n} \| \bu_k \|^2  + \alpha \| \hat{u}_{k}  \|  \\
%     &+ \frac{M_2}{n} (  \| \bx_k  - \bar{\bx}_k \|^2 + \| \bx_{k+1}  - \bar{\bx}_{k+1} \|^2 ) \\
%      \leq & \frac{8Q + \sqrt{n}L_2 + M_2}{n} \|\bx_k - \bar{\bx}_k\|^2 +\frac{2Q\alpha^2}{n} \| \bu_k \|^2  + \alpha \| \hat{u}_{k}  \|  \\
%      &+ \frac{M_2}{n} \| \bx_{k+1}  - \bar{\bx}_{k+1} \|^2.
%   %& \leq \left( \frac{36}{n} \|\bx_k - \bar{\bx}_k\|^2 + \frac{16\alpha^2}{n}\| \grad f(\bx_k) \|^2  \right) + 8M_2^2CL^2  ( \alpha^2 \| \bx_k  - \bar{\bx}_k \|^2 +\alpha^2 \| \bx_{k+1}  - \bar{\bx}_{k+1} \|^2 )
%     \end{aligned}
% \end{equation}
% The proof is completed.
% \end{proof}



















\subsection{Optimality error}\label{sec:proof:dprgt}

\begin{revisionblock}
We replace the separate summation of the estimation and tracking errors by a
single coupled argument. Define the normalized deterministic sequences
\begin{equation}\label{eq:five-errors}
\begin{aligned}
 \mathsf{G}_k&:=\mathbb{E}\|\grad f(\bar x_k)\|^2,&
 \mathsf{S}_k&:=\frac1n\mathbb{E}\|\bs_k\|^2,\\
 \mathsf{E}_k&:=\frac1n\mathbb{E}
   \|\bd_k-\grad f(\bx_k)\|^2,&
 \mathsf{T}_k&:=\frac1n\mathbb{E}\|\bs_k-\hat{\bd}_k\|^2,\\
 \mathsf{X}_k&:=\frac1n\mathbb{E}\|\bx_k-\bar\bx_k\|^2.&
\end{aligned}
\end{equation}

\begin{lemma}[Coupled one-step inequalities]\label{lem:coupled-one-step}
Suppose that Assumptions~\ref{assum-f}--\ref{assum:stochatis} hold,
$B\geq G_{\max}$, $0<\tau\leq1$, and the local projection conditions in
Theorem~\ref{lem:consensus} are satisfied. There are constants $d_X,C_D>0$,
independent of $k,K,\alpha,\tau$, such that
\begin{equation}\label{eq:coupled-E-main}
 \mathsf{E}_{k+1}\le(1-\tau)^2\mathsf{E}_k+64\bar L^2\mathsf{X}_k
      +16\bar L^2\alpha^2\mathsf{S}_k+2\tau^2\nu^2.
\end{equation}
\begin{equation}\label{eq:coupled-T-main}
 \mathsf{T}_{k+1}\le q_T\mathsf{T}_k+h_T(96\bar L^2\mathsf{X}_k
      +24\bar L^2\alpha^2\mathsf{S}_k
      +3\tau^2\mathsf{E}_k+3\tau^2\nu^2).
\end{equation}
\begin{equation}\label{eq:coupled-X-main}
 \mathsf{X}_{k+1}\le q_X\mathsf{X}_k+c_X\alpha^2\mathsf{S}_k.
\end{equation}
\begin{equation}\label{eq:coupled-D-main}
 \begin{aligned}
 \mathbb{E} f(\bar x_{k+1})
 &\le\mathbb{E} f(\bar x_k)-\frac\alpha4\mathsf{G}_k
      -\frac\alpha4\mathsf{S}_k
      +\frac{3\alpha}{2}(\mathsf{E}_k+\mathsf{T}_k)\\
 &\quad+d_X\alpha\mathsf{X}_k+C_D\alpha^3.
 \end{aligned}
\end{equation}
where
\[
 q_T:=\frac{1+\sigma_2^2}{2},\quad
 h_T:=\frac{1+\sigma_2^2}{1-\sigma_2^2},\quad
 q_X:=\frac{1+\rho^2}{2},\quad
 c_X:=\kappa_M^2\frac{1+\rho^2}{1-\rho^2}.
\]
In particular, $q_T,q_X<1$. The proof is given in
Appendix~\ref{appen}.
\end{lemma}

\begin{proof}[Proof of Theorem~\ref{thm:dprgt}]
For brevity set
\[
 a_X=64\bar L^2,\quad a_S=16\bar L^2,\quad
 b_E=3h_T,\quad b_X=96h_T\bar L^2,\quad
 b_S=24h_T\bar L^2,
\]
and let $\delta_T=1-q_T$ and $\delta_X=1-q_X$. Choose constants in the
following order:
\begin{equation}\label{eq:potential-weights}
 A=\frac72,\qquad B_0=\frac{5}{2\delta_T},\qquad
 C_0=\frac{Aa_X+2}{\delta_X},\qquad
 \kappa\ge16(Aa_S+C_0c_X).
\end{equation}
After these choices, take the constant $\theta$ in
\eqref{eq:repaired-parameters} sufficiently small so that all local
projection, descent, and coefficient restrictions below hold. Define
\begin{equation}\label{eq:coupled-potential}
 \Phi_k:=\mathbb{E} f(\bar x_k)
 +A\frac\alpha\tau\mathsf{E}_k+B_0\alpha\mathsf{T}_k
 +C_0\frac\alpha\tau\mathsf{X}_k.
\end{equation}
Using Lemma~\ref{lem:coupled-one-step}, $\tau=\kappa\alpha^2$, and
$1-(1-\tau)^2=\tau(2-\tau)$, the coefficients of
$\mathsf{E}_k,\mathsf{T}_k,\mathsf{X}_k$, and $\mathsf{S}_k$ in
$\Phi_{k+1}-\Phi_k$ are bounded above by
\begin{align*}
 &\alpha\{\tfrac32-A(2-\tau)+B_0b_E\tau^2\},
 &&-\alpha\{B_0\delta_T-\tfrac32\},\\
 &\frac\alpha\tau\{-C_0\delta_X+Aa_X
       +\tau(d_X+B_0b_X)\},
 &&\alpha\{-\tfrac14+(Aa_S+C_0c_X)/\kappa+B_0b_S\alpha^2\}.
\end{align*}
Thus, for this sufficiently small $\theta$, there is a constant $C_\Phi>0$
such that
\begin{equation}\label{eq:master-drift}
 \Phi_{k+1}-\Phi_k
 \le-\frac\alpha4\mathsf{G}_k-\frac\alpha8\mathsf{S}_k
     -\alpha\mathsf{E}_k-\alpha\mathsf{T}_k
     -\frac\alpha\tau\mathsf{X}_k+C_\Phi\alpha^3.
\end{equation}
The last term contains the oracle contributions because
$\alpha\tau\nu^2=\kappa\nu^2\alpha^3$; all constants are independent of
$k$ and $K$.

It remains to handle the one-sample initialization. Let
$N=\lfloor K/2\rfloor$. The variance assumption gives
$\mathsf{E}_0\le\nu^2$, consensus initialization gives $\mathsf{X}_0=0$, and
$\bs_0=\bd_0$ together with $\|d_{i,0}\|\le B$ gives
$\mathsf{T}_0=n^{-1}\mathbb E
\|((I-J)\otimes I_d)\bd_0\|^2\le B^2$.
Consequently, summing \eqref{eq:master-drift} from $0$ to $N$ gives
\[
 \alpha\sum_{k=0}^{N}\mathsf{E}_k
 \le\Phi_0-f_*+C_\Phi\alpha^3(N+1)=\mathcal O(\alpha^{-1}).
\]
Here $\alpha^3K=\mathcal O(1)$. Since
$N+1=\Theta(\alpha^{-3})$, there is a deterministic index $j\le N$ such
that $\mathsf{E}_j=\mathcal O(\alpha)$. Compactness, the uniform bounds on
$\mathsf{S}_j$ and $\mathsf{T}_j$, and Theorem~\ref{lem:consensus} then give
\[
 \Phi_j-f_*=\mathcal O(1):
 \quad \frac\alpha\tau\mathsf{E}_j=\mathcal O(1),\qquad
 \alpha\mathsf{T}_j=\mathcal O(\alpha),\qquad
 \frac\alpha\tau\mathsf{X}_j=\mathcal O(\alpha).
\]
Summing \eqref{eq:master-drift} again from $j$ to $K$ yields
\[
 \frac\alpha4\sum_{k=j}^{K}\mathsf{G}_k
 \le\Phi_j-f_*+C_\Phi\alpha^3(K-j+1)=\mathcal O(1).
\]
Since $K-j+1\ge K/2$, we obtain
\[
 \min_{j\le k\le K}\mathsf{G}_k
 =\mathcal O((\alpha K)^{-1})=\mathcal O(K^{-2/3}),
\]
which proves \eqref{eq:repaired-optimality-rate}.
\end{proof}

\begin{proof}[Proof of Corollary~\ref{coll}]
Choose an index $k_*$ attaining the minimum in
\eqref{eq:repaired-optimality-rate}. Theorem~\ref{lem:consensus} gives the
consensus estimate at the same index. Since Algorithm~\ref{alg:drgta} uses
two same-sample stochastic-gradient evaluations per node in each recursive
update (and one at initialization), its per-node SFO cost is $\mathcal O(K)$.
Taking $K=\mathcal O(\epsilon^{-3/2})$ completes the proof.
\end{proof}
\end{revisionblock}

%\begin{proof}
%It follows from Lemma \ref{lemma:lipsctz} and $L_g \leq L$ that 
%\bee
%\| \hat{g}_k - \grad f(\bar{x}_k) \|^2 \leq \frac{1}{n} \sum_{i=1}^n \|\grad f_i(x_k) -  \grad f_i(\bar{x}_k) \|^2 \leq \frac{L^2}{n} \|\bx_k - \bar{\bx}_k \|^2.
%\eee
%and
%\begin{equation}\label{Lip-ineq}
%    \begin{aligned}
%    f(\bar{x}_{k+1})  \leq & f(\bar{x}_k) + \left<\grad f(\bar{x}_k),\bar{x}_{k+1}-\bar{x}_k  \right>+ \frac{L}{2}\|\bar{x}_{k+1}-\bar{x}_k\|^2\\
%     = & f(\bar{x}_k) + \left<\hat{g}_k ,\bar{x}_{k+1}-\bar{x}_k  \right>+\left<\grad f(\bar{x}_k) -\hat{g}_k,\bar{x}_{k+1}-\bar{x}_k  \right> + \frac{L}{2}\|\bar{x}_{k+1}-\bar{x}_k\|^2 \\
%      \leq & f(\bar{x}_k) + \left<\hat{g}_k ,\bar{x}_{k+1}-\bar{x}_k  \right>+ \frac{3L}{4}\|\bar{x}_{k+1}-\bar{x}_k\|^2  + \frac{1}{L} \| \grad f(\bar{x}_k) -\hat{g}_k \|^2 \\
%        \leq & f(\bar{x}_k) + \left<\hat{g}_k ,\hat{x}_{k+1}-\hat{x}_k  \right>  + \left<\hat{g}_k,  \bar{x}_{k+1}- \hat{x}_{k+1}+\hat{x}_k - \bar{x}_k  \right> \\
%        &+ \frac{L}{n}  \| \bx_k - \bar{\bx}_k \|^2 + \frac{3L}{4}\|\bar{x}_{k+1}-\bar{x}_k\|^2.
%    % & \leq   \mathcal{L}_{\alpha}(\bx_k) + \left<\grad f(\bx_k) - \bs_k,\bx_{k+1}-\bx_k  \right>- ( \frac{1}{2\alpha} -\frac{L}{2})\|\bx_{k+1}-\bx_k\|^2\\
%    \end{aligned}
%\end{equation}
%By Young's inequality, we get 
%\begin{equation} \label{eq:grad-prod2}
%    \left<\hat{g}_k,  \bar{x}_{k+1}- \hat{x}_{k+1}+\hat{x}_k - \bar{x}_k  \right> \leq \frac{\alpha^2L}{2}\| \hat{g}_k  \|^2 + \frac{1}{\alpha^2L}(\| \bar{x}_{k+1}- \hat{x}_{k+1} \|^2 + \|\hat{x}_k - \bar{x}_k\|^2).
%\end{equation}
%Combining \eqref{Lip-ineq} and \eqref{eq:grad-prod2} leads to
%\begin{equation}\nonumber
%    \begin{aligned}
%     f(\bar{x}_{k+1})  \leq & f(\bar{x}_k) + \underbrace{\left<\hat{g}_k ,\hat{x}_{k+1}-\hat{x}_k  \right>}_{b_1}  +  \frac{\alpha^2L}{2}\| \hat{g}_k   \|^2 + \underbrace{ \frac{1}{\alpha^2L}(\| \bar{x}_{k+1}- \hat{x}_{k+1} \|^2 + \|\hat{x}_k - \bar{x}_k\|^2)}_{b_2}\\
%     & + \frac{L}{n}  \| \bx_k - \bar{\bx}_k \|^2 + \underbrace{ \frac{3L}{4}\|\bar{x}_{k+1}-\bar{x}_k\|^2}_{b_3}. \\
%    \end{aligned}
%\end{equation}
%Now, let us bound $b_1$, $b_2$, and $b_3$, respectively. For $b_1$, it holds that
%\begin{equation} \label{eq:b1-0}
%    \begin{aligned}
%    b_1  =& \left< \hat{g}_k , \frac{1}{n} \sum_{i=1}^n( x_{i,k+1} - x_{i,k} + \alpha v_{i,k} + \nabla \phi_i(\bx_k)   ) \right> -  \left< \hat{g}_k, \frac{1}{n} \sum_{i=1}^n(  \alpha v_{i,k} + \nabla \phi_i(\bx_k)   ) \right> \\
%    \leq &  \left< \hat{g}_k   , \frac{1}{n} \sum_{i=1}^n( x_{i,k+1} - x_{i,k} + \alpha v_{i,k} + \nabla \phi_i(\bx_k)  ) \right> +\left< \hat{g}_k, \frac{1}{n} \sum_{i=1}^n \alpha(s_{i,k} - v_{i,k}) \right> - \alpha  \|\hat{g}_k\|^2.
%    \end{aligned}
%\end{equation}
%{
%Before analyzing the term $b_1$, we need the following inequality
%\begin{equation}
%    \begin{aligned}
%    & \sum_{i=1}^n  \|\hat{g}_k - \mathcal{P}_{T_{x_{i,k}}\Mcal}( \grad f(\bar{x}_{k})) \|^2 \\
%     & \leq 2  \sum_{i=1}^n  (\|\hat{g}_k - \grad f(\bar{x}_{k}) \|^2 + \| \grad f(\bar{x}_{k})  - \mathcal{P}_{T_{x_{i,k}}\Mcal}( \grad f(\bar{x}_{k})) \|^2 ) \\
%     & \leq 2L^2\| \bx_k - \bar{\bx}_k \|^2 +2  \sum_{i=1}^n \|\mathcal{P}_{T_{\bar{x}_{k}}\Mcal} (\grad f(\bar{x}_{k}))  - \mathcal{P}_{T_{x_{i,k}}\Mcal}( \grad f(\bar{x}_{k}))
% \|^2 \\
% & \leq 2L^2\| \bx_k - \bar{\bx}_k \|^2 + 2L_2^2 \sum_{i=1}^n \| x_{i,k} - \bar{x}_k \|^2 \| \grad f(\bar{x}_{k}) \|^2 \\
% & \leq 2L^2(1+L_2^2)\| \bx_k - \bar{\bx}_k \|^2,
%    \end{aligned}
%\end{equation}
%where the third inequality is due to the Lipschitz continuity of $\Pcal_{T_{x}\Mcal}(\cdot)$ over $x \in \Mcal$, and the constant $L_2$ is given in Lemma \ref{lem:grad-consensus}. Now we are going to bound $b_1$, 
%since $s_{i,k} - v_{i,k} \in N_{x_{i,k}}\Mcal$, it follows that
%\be \label{eq:b1-1}
%\begin{aligned}
%    & \left< \hat{g}_k, \frac{1}{n} \sum_{i=1}^n \alpha(s_{i,k} - v_{i,k}) \right>
%   = \frac{\alpha}{n} \sum_{i=1}^n\left< \hat{g}_k - \mathcal{P}_{T_{x_{i,k}}\Mcal}( \grad f(\bar{x}_{k})),  s_{i,k} - v_{i,k} \right> \\
%   \leq & \frac{1}{4nL} \sum_{i=1}^n \|\hat{g}_k - \mathcal{P}_{T_{x_{i,k}}\Mcal}( \grad f(\bar{x}_{k})) \|^2 + \frac{\alpha^2L}{n}\sum_{i=1}^n \| \Pcal_{N_{x_{i,k}}\Mcal}(s_{i,k}) \|^2\\
% % \leq &  {  \frac{1}{4nL}  \|\hat{G}_k - \grad f(\bx_{k})  \|^2 + \frac{\alpha^2L}{n}\| \bs_k \|^2  }  \\
% % \leq &  {  \frac{1}{2nL} \left(  \|\hat{G}_k - \grad f(\bar{\bx}_{k})  \|^2 + \| \grad f(\bar{\bx}_{k}) -  \grad f(\bx_{k})  \|^2 \right) + \frac{\alpha^2L}{n}\| \bs_k \|^2  }  \\
% %   \leq &  {  \frac{1}{2nL} \left(  L^2\| \bx_k - \bar{\bx}_k \|^2 + L^2\| \bx_k - \bar{\bx}_k \|^2 \right) + \frac{\alpha^2L}{n}\| \bs_k \|^2  }  \\
%   \leq &  \frac{L(1+L_2^2)}{2n} \| \bx_k - \bar{\bx}_k \|^2 + \frac{\alpha^2L}{n}\| \bs_k \|^2.
%\end{aligned}
%\ee
%}
%% Since $s_{i,k} - d_{i,k} \in N_{x_{i,k}}\Mcal$, it follows that
%% \be \label{eq:b1-1}
%% \begin{aligned}
%%     & \left< \hat{g}_k, \frac{1}{n} \sum_{i=1}^n \alpha(s_{i,k} - d_{i,k}) \right>
%%    = \frac{\alpha}{n} \sum_{i=1}^n\left< \hat{g}_k - \grad f_i(x_{i,k}),  s_{i,k} - d_{i,k} \right> \\
%%    \leq & \frac{1}{4nL} \sum_{i=1}^n \|\hat{g}_k - \grad f_i(x_{i,k})  \|^2 + \frac{\alpha^2L}{n}\sum_{i=1}^n \| \Pcal_{N_{x_{i,k}}\Mcal}(s_{i,k}) \|^2\\
%%    \leq & \frac{1}{4n^2L} \sum_{i=1}^n \sum_{j=1}^n\|\grad f_j(x_{j,k}) - \grad f_i(x_{i,k})  \|^2 + \frac{\alpha^2L}{n}\| \bs_k \|^2 \\
%%    &
%%    \leq  \frac{L}{n} \| \bx_k - \bar{\bx}_k \|^2 + \frac{\alpha^2L}{n}\| \bs_k \|^2.
%% \end{aligned}
%% \ee
%Let  $\nabla \phi_i(\bx_k) + \alpha v_{i,k} = c_{i,1} + c_{i,2}, $
%where $c_{i,1} = \Pcal_{T_{x_{i,k}}\Mcal}(\nabla \phi_i(\bx_k) + \alpha v_{i,k}) , c_{i,2} = \nabla \phi_i(\bx_k)+ \alpha v_{i,k} - c_{i,1} \in N_{x_{i,k}}\Mcal$.  Therefore,
%\be \label{eq:b1-2}
%\begin{aligned}
%&\left< \hat{g}_k   , \frac{1}{n} \sum_{i=1}^n( x_{i,k+1} - x_{i,k} + \alpha v_{i,k} + \nabla \phi_i(\bx_k)  ) \right> \\
%= & \frac{1}{n} \sum_{i= 1}^n \left< \hat{g}_k,  \Pcal_{\Mcal}(x_{i,k} - c_{i,1} - c_{i,2}  ) - [x_{i,k}  - c_{i,1}    ]\right>+\frac{1}{n}\sum_{i= 1}^n \left< \hat{g}_k,  c_{i,2}\right> \\
%\leq & \frac{LQ }{n} \sum_{i=1}^n \|  v_i\|^2 +\frac{1}{n} \sum_{i= 1}^n \left< \hat{g}_k - \grad f_i(x_{i,k}),  c_{i,2}\right>\\
%\leq & \frac{LQ }{n} \sum_{i=1}^n \|  v_i\|^2 + \frac{1}{4nL}\sum_{i= 1}^n
% \|\hat{g}_k - \grad f_i(x_{i,k})\|^2 + \frac{L}{n}\sum_{i= 1}^n \|c_{i,2}\|^2 \\
% %\leq & \frac{LQ+L }{n} \sum_{i=1}^n \|  v_i\|^2 + \frac{L}{n} \|\bar{\bx}_k - \bx_k\|^2 \\
%  \leq & \frac{LQ+L }{n} \|(I_{nd}-\bW)\bx_k + \alpha \bv_k\|^2 + \frac{L^2}{n}  \|\bar{\bx}_k - \bx_k\|^2 \\
%   \leq & \frac{4LQ+3L  }{n}\|\bar{\bx}_k - \bx_k\|^2  + \frac{2LQ+2L }{n}\alpha^2 \|\bs_k\|^2.
%\end{aligned}
%\ee
%Plugging \eqref{eq:b1-1} and \eqref{eq:b1-2} into \eqref{eq:b1-0} gives
%\be \label{eq:b1}
%\begin{aligned}
%b_1 &\leq  { \frac{8LQ+ 7 L  + LL_2^2}{2n}  }\|\bx_k - \bar{\bx}_k \|^2 +\frac{2LQ+3L}{n} \alpha^2 \| \bs_k\|^2 -  \alpha  \|\hat{g}_k\|^2,
%\end{aligned}
%\ee
%For $b_2$, applying Lemma \ref{lemma:quadratic} and Lemma \ref{lemma:dpgta:consensus} yields
%\begin{equation} \label{eq:b2}
%\begin{aligned}
%      b_2 & \leq \frac{M_2^2}{\alpha^2 L n^2}(\| \bx_k - \bar{\bx}_k\|^4 + \| \bx_{k+1} - \bar{\bx}_{k+1}\|^4  ) \\
%      &
%     \leq \frac{M_2^2D L }{ n}(\| \bx_k - \bar{\bx}_k\|^2 + \| \bx_{k+1} - \bar{\bx}_{k+1}\|^2  ),
%\end{aligned}
%\end{equation}
%Since $\|\bv_k\|^2 \leq \|\bs_k\|^2 \leq 4\zeta^2L^2$ in Lemma \ref{lem:stay-neighborhood}. Since $\|\hat{y}_k\|^2 = \frac{1}{n^2} \|\sum_{i=1}^n d_{i,k}\|^2 \leq \frac{1}{n}\sum_i\|d_{i,k}\|^2 = \frac{1}{n}\|\bv_k\|^2$,  combining with \eqref{eq:sum-consen-grad} yields
%\begin{equation} \label{eq:b3}
%    \begin{aligned}
%    b_3  %  \leq &\frac{3L}{4} \left( \frac{4(8Q + \sqrt{n}L_2 + M_2)^2}{n^2} \|\bx_k - \bar{\bx}_k\|^4 +\frac{16Q^2\alpha^4}{n^2} \| \bv_k \|^4  + 4\alpha^2 \| \hat{d}_{k}  \|^2  + \frac{4M_2^2}{n^2} \| \bx_{k+1}  - \bar{\bx}_{k+1} \|^4\right)\\
%     \leq &\frac{3L}{4} \left(\frac{4(8Q + \sqrt{n}L_2 + M_2)^2}{n^2} \|\bx_k - \bar{\bx}_k\|^4 +\frac{16Q^2\alpha^4}{n^2} \| \bv_k \|^4  \right. \\
%     & + \left. \frac{4\alpha^2}{n} \| \bv_k  \|^2  + \frac{4M_2^2}{n^2} \| \bx_{k+1}  - \bar{\bx}_{k+1} \|^4\right)\\
%      \leq & \frac{4DL^3(8Q + \sqrt{n}L_2 + M_2)^2}{n}\alpha^2 \|\bx_k - \bar{\bx}_k\|^2   + \frac{4DL^3M_2^2}{n}\alpha^2 \| \bx_{k+1}  - \bar{\bx}_{k+1} \|^2 \\
%     & + \frac{192Q^2L^3\alpha^4 + 3L\alpha^2}{n}\| \bs_k \|^2.
%  %& \leq \left( \frac{36}{n} \|\bx_k - \bar{\bx}_k\|^2 + \frac{16\alpha^2}{n}\| \grad f(\bx_k) \|^2  \right) + 8M_2^2CL^2  ( \alpha^2 \| \bx_k  - \bar{\bx}_k \|^2 +\alpha^2 \| \bx_{k+1}  - \bar{\bx}_{k+1} \|^2 )
%    \end{aligned}
%\end{equation}
%Then, combining \eqref{eq:b1}, \eqref{eq:b2}, and \eqref{eq:b3} gives
%\begin{equation}\nonumber
%\begin{aligned}
%   &   f(\bar{x}_{k+1})  \leq  f(\bar{x}_k) + b_1 + \frac{\alpha^2 L^2}{2} \|\hat{g}_k\|^2 + b_2 + \frac{L}{n}\|\bx_k - \bar{\bx}_k\|^2 + b_3 \\
%      \leq & f(\bar{\bx}_k) - ( \alpha - \frac{L^2}{2} \alpha^2 ) \|\hat{g}_k \|^2 + \frac{(2Q+ 3)L + 192 Q^2\alpha^2L^3 + 3L}{n} \alpha^2 \|  \bs_k\|^2\\
%      & + \frac{{4LQ+ 3.5 L  + 0.5LL_2^2} + M_2^2DL +  4DL^3(8Q + \sqrt{n}L_2 + M_2)^2\alpha^2}{n} \|\bx_k - \bar{\bx}_k\|^2  \\
%      & + \frac{M_2^2D L + 4DL^3M_2^2\alpha^2}{ n}  \| \bx_{k+1}  - \bar{\bx}_{k+1} \|^2,
%\end{aligned}
%\end{equation}
%where the second inequality utilizes $\frac{1}{n}\| \bx_k - \bar{\bx}_k \|^2 \leq  D L^2 \alpha^2$ in Lemma \ref{lemma:dpgta:consensus}. The proof is completed. 
%%The proof is completed.
%\end{proof}



%\begin{proof}\label{eq:proof1}[Proof of Theorem \ref{thm:dprgt}].
%Since $\rho_t = 2\sigma_2 \in (0,1)$, applying \cite[Lemma 2]{xu2015augmented} to \eqref{eq:dgta:consensus} yields 
%\be\label{eq:sum:consensus}
%\frac{1}{n}\sum_{k=0}^{K+1}\| \bx_k - \bar{\bx}_k \|^2 \leq \tilde{C}_0 \frac{4\alpha^2}{n}\sum_{k=0}^{K+1}\| \bs_k \|^2 + \tilde{C}_1,
%\ee 
%where $\tilde{C}_0 = \frac{2}{(1-\rho_t)^2}$ and $\tilde{C}_1 = \frac{2}{1-\rho_t^2}\frac{1}{n}\|\bx_0 - \bar{\bx}_0\|^2$. It follows from Lemma \ref{lemma:dgta:majorized} that
%\begin{equation} \label{eq:desc}
%\begin{aligned}
%   &  f(\bar{x}_{k+1}) \\
%      \leq & f(\bar{x}_0) - (\alpha  - \frac{\alpha^2L^2}{2})\sum_{k=0}^{K} \|\hat{g}_k\|^2 + \mathcal{C}_1 \alpha^2 \frac{1}{n} \sum_{k=0}^{K} \|\bs_k \|^2 +(\mathcal{C}_2 + \mathcal{C}_3) \frac{1}{n}\sum_{k=0}^{K+1} \| \bx_k - \bar{\bx}_k\|^2\\
%      \leq & f(\bar{x}_0) - \frac{\alpha    }{2} \sum_{k=0}^{K} \|\hat{g}_k\|^2 + (\mathcal{C}_1 + 4\tilde{C}_0(\mathcal{C}_2 + \mathcal{C}_2)) \alpha^2 \frac{1}{n} \sum_{k=0}^{K+1} \|\bs_k \|^2 +(\mathcal{C}_2 + \mathcal{C}_3)\tilde{C}_1 \\
%       \leq & f(\bar{x}_0) - \frac{\alpha    }{2} \sum_{k=0}^{K} \|\hat{g}_k\|^2 + (\mathcal{C}_1 + 4\tilde{C}_0(\mathcal{C}_2 + \mathcal{C}_2)) \alpha^2 \frac{1}{n} \sum_{k=0}^{K} \|\bs_k \|^2 \\
%       & + 4(\mathcal{C}_1 + 16\tilde{C}_0(\mathcal{C}_2 + \mathcal{C}_2)) \alpha^2L^2 +(\mathcal{C}_2 + \mathcal{C}_3)\tilde{C}_1 \\
%        \leq & f(\bar{x}_0) - \frac{\alpha    }{2} \sum_{k=0}^{K} \|\hat{g}_k\|^2 + (\mathcal{C}_1 + 4\tilde{C}_0(\mathcal{C}_2 + \mathcal{C}_2)) \alpha^2 \frac{1}{n} \sum_{k=0}^{K} \|\bs_k \|^2 + \tilde{C}_4,
%     \end{aligned}
%\end{equation}
%where we use $\alpha< \frac{1}{L^2}, \frac{1}{n}\|\bs_k\|^2 \leq 16 L^2$ and $\tilde{C}_4: = 4(\mathcal{C}_1 + 16\tilde{C}_0(\mathcal{C}_2 + \mathcal{C}_2)) \alpha^2  L^2 +(\mathcal{C}_2 + \mathcal{C}_3)\tilde{C}_1$. Next, we need to associate $\|\hat{g}_k\|^2$ with $\|\bs_k\|^2$. Since  $\| \bs_k  \|  \leq \|\bs_k - \hat{\bf G}_k \| + \|\hat{\bf G}_k \|$, we have 
%\be\label{eq:tmp10}
%-\sum_{k=0}^{K} \|\hat{g}_k \|^2  = -\frac{1}{n}\sum_{k=0}^{K} \|\hat{\bf G}_k \|^2  \leq \frac{1}{n} \sum_{k=0}^{K}\|\bs_k - \hat{\bf G}_k \|^2 - \frac{1}{2n}\sum_{k=0}^{K}\|\bs_k\|^2.
%\ee
%Applying \cite[Lemma 2]{xu2015augmented} into \eqref{eq:diff-s-g} yields
%\bee
%\begin{aligned}
%\frac{1}{n}\sum_{k=0}^{K}\|\bs_k - \hat{\bf G}_k \|^2 &\leq \tilde{C}_2\frac{1}{n}\sum_{k=0}^{K}\|{\bf G}^{k+1} - {\bf G}^k \|^2 + \tilde{C}_3 \\
%& \overset{\eqref{eq:diff-g-g}}{ \leq} \tilde{C}_2\frac{1}{n}\sum_{k=0}^{K} (32L^2\|\bx_k - \bar{\bx}_k \|^2 + 8L^2\alpha^2 \|\bs_k\|^2  ) + \tilde{C}_3 \\
%& \overset{\eqref{eq:sum:consensus}}{ \leq} \tilde{C}_2(128\tilde{C}_0  + 8)L^2\alpha^2\frac{1}{n}\sum_{k=0}^{K}  \|\bs_k\|^2   + 32\tilde{C}_1\tilde{C}_2L^2  + \tilde{C}_3,
%\end{aligned}
%\eee
%where $\tilde{C}_2 = \frac{2}{(1-\sigma_2)^2}$ and $\tilde{C}_3 = \frac{2}{1-\nu^2t_2}\frac{1}{n}\|\bs_0 - \hat{\mathbf{g}}_0\|^2$.
%Plugging the above inequality into \eqref{eq:tmp10} leads to
%\bee
%-\sum_{k=0}^{K} \|\hat{g}_k \|^2    \leq  \left[ \tilde{C}_2(128\tilde{C}_0  + 8)L^2\alpha^2 - \frac{1}{2}  \right]  \frac{1}{n}\sum_{k=0}^{K}\|\bs_k\|^2  + 32\tilde{C}_1\tilde{C}_2L^2  + \tilde{C}_3.
%\eee
%Since $\alpha <\min\left\{1, \frac{1}{4\left(\tilde{C}_2(128\tilde{C}_0  + 8)L^2 + 2(\mathcal{C}_1 + 4\tilde{C}_0(\mathcal{C}_2 + \mathcal{C}_2)) \right) }\right\} $,
%it follows from \eqref{eq:desc} that
%\bee
%\begin{aligned}
%&   f(\bar{x}_{k+1})\\
%\leq &  f(\bar{x}_0) - \frac{\alpha    }{2}   \left[\frac{1}{2} -  \tilde{C}_2(128\tilde{C}_0  + 8)L^2\alpha^2 - 2(\mathcal{C}_1 + 4\tilde{C}_0(\mathcal{C}_2 + \mathcal{C}_2)) \alpha  \right] \frac{1}{n}  \sum_{k=0}^{K}\|\bs_k\|^2 \\
%      &  +(\tilde{C}_4 + 16\tilde{C}_1\tilde{C}_2L^2 \alpha + \frac{\alpha}{2}\tilde{C}_3 ) \\
%       \leq & f(\bar{x}_0) - \frac{\alpha    }{8}\frac{1}{n}     \sum_{k=0}^{K}\|\bs_k\|^2   +(\tilde{C}_4 + 16\tilde{C}_1\tilde{C}_2L^2 \alpha + \frac{\alpha}{2}\tilde{C}_3 ). 
%     \end{aligned}
%\eee
%This implies
%\bee
%\frac{\alpha}{8} \sum_{k=0}^{K} \|\hat{g}_k \|^2  \leq \frac{\alpha}{8} \frac{1}{n} \sum_{k=0}^{K} \|\bs_k\|^2 \leq f(\bar{x}_0) - f^* + \tilde{C}_5,
%\eee
%where $\tilde{C}_5 := \tilde{C}_4 + 16\tilde{C}_1\tilde{C}_2L^2 \alpha + \frac{\alpha}{2}\tilde{C}_3$ and $f^* = \min_{x\in \Mcal} f(x)$. Furthermore, 
%\bee
%\min_{k=0,\cdots,K} \| \hat{g}_k\|^2 = \min_{k=0,\cdots,K} \| \hat{s}_k\|^2 \leq  \min_{k=0,\cdots,K} \frac{1}{n}\|\bs_k\|^2 \leq \frac{8(f(\bar{x}_0) - f^* + \tilde{C}_5)}{\alpha K}.
%\eee
%Then, it follows from \eqref{eq:sum:consensus} that
%\bee
%\min_{k=0,\cdots,K} \frac{1}{n} \|\bx_k - \bar{\bx}_k \|^2  \leq \frac{32\tilde{C}_0\alpha(f(\bar{\bx}_0) - f^* + \tilde{C}_5) + \tilde{C}_1}{ K}.
%\eee
%Since 
%$
%\|\grad f(\bar{x}_k) \|^2  \leq 2\| \hat{g}_k\|^2 + 2\|\grad f(\bar{x}_k)  - \hat{g}_k \|^2  \leq 2\| \hat{g}_k\|^2 + \frac{2L^2}{n}\| \bx_k - \bar{\bx}_k\|^2. 
%$
%We conclude that 
%\bee
%\min_{k=0,\cdots,K} \|\grad f(\bar{x}_k) \|^2  \leq \frac{(16 +64\tilde{C}_0\alpha^2 )(f(\bar{x}_0) - f^* + \tilde{C}_5)+ 2L^2 \tilde{C}_1\alpha}{\alpha K}.
%\eee
%The proof is completed.
%\end{proof}

%%%%%%%%%%%%%%%%%%%%%%%%%%%%%%%%%%%%%

% \section{Decentralized projected Riemannian gradient method}
% % In the centralized setting, the iterative scheme of the projected gradient method is
% % \[ x_{k+1} = \Pcal(x_k - \alpha  \grad f(x_k)), \]
% % where $\alpha  > 0$ is a step size and $\grad f(x_k)$ represents the Riemannian gradient of $f$ at $x_k$.
% In this section, we study a decentralized projected Riemannian gradient method (DPRSRM) for solving problem \eqref{prob:original}. Let us denote $f(\bx): = \frac{1}{n}\sum_{i=1}^n f_i(x_i)$. Given the adjacency matrix $W$ of the communication network, in the $k$-th iteration, the DPRSRM does the following update
% \begin{equation}\label{eq:dpg}
% \begin{aligned}
%    \bx_{k+1}  & = \Pcal_{\Mcal^n}(\bW \bx_k - \alpha  \grad f(\bx_k)),
%    %& = \Pcal_{\Mcal^n}(\bx_k - \nabla \phi(\bx_k) - \alpha  \grad f(\bx_k)),
% \end{aligned}
% \end{equation}
% where $\alpha > 0$ is the step size, $\grad f(\bx_k)$ denotes the Riemannian gradient, and $t \geq 1$ is an integer. We present the detailed algorithm in Algorithm \ref{alg:drgta}.

% \begin{algorithm}[htbp]
% \caption{Decentralized Projected Riemannian Gradient Descent (DPRSRM) for solving \eqref{prob:original}} \label{alg:drgta}
% \begin{algorithmic}[1]
% \REQUIRE  Initial point $\bx_0\in \mathcal{M}^n$, an integer $t$, set $k = 1$.
% \WHILE {the condition is not met}
% \STATE Choose diminishing step size $\alpha = \mathcal{O}(1/\sqrt{k})$
% \STATE Update $$x_{i,k+1} = \Pcal_{\Mcal}\left(\sum_{j=1}^n W_{ij}x_{i,k} - \alpha  \grad f_i(x_{i,k}) \right),$$ for each node $i\in [n]$, in parallel.
% \STATE Set $k=k+1$.
% \ENDWHILE
% \end{algorithmic}
% \end{algorithm}

% In the following, we are going to show the convergence of DPRSRM, i.e., Algorithm \ref{alg:drgta}. The main techniques used are the linear convergence result of the projected gradient method for solving consensus problem \eqref{prob:consensus} and the Lipschitz type equalities established in Section \ref{sec:ineq-man}. 

% The following lemma is vital for proving the convergence result of Algorithm \ref{alg:drgta}.

% If the step size $\alpha$ and the integer $t$ are chosen properly, then, with proper initialization on $\bx_0$, the consensus error can be bounded. 
% \begin{lemma}\label{lemma:dpg:consensus}
% Assume that the sequence $\{\bx_k\}$ is generated by Algorithm \ref{alg:drgta}, and  $\|\hat{x}_0 - \bar{x}_0\| \leq \frac{1}{2}{\delta}$, $\| \grad f_i(x_{i,k}) \| \leq L$,  $\alpha \leq {\delta}/(8L), t \geq  \left\lceil \log_{\sigma_2}\left(\frac{{\delta}}{16\sqrt{n}\zeta}\right) \right \rceil$.  The following holds
% \begin{equation}
% \begin{aligned}
%  \|\bx_{k+1} - \bar{\bx}_{k+1}\| & \leq 2\sigma_2 \|\bx_{k} - \bar{\bx}_{k}\| + 2\sqrt{n}\alpha L .
%  %& \leq (L\sigma_2)^{k+1} \|\bx_0 - \bar{\bx}_0\| + L\sqrt{n}B \sum_{l=0}^k(L\sigma_2)^{k-l}\alpha_l.
% \end{aligned}
% \end{equation}
% \end{lemma}
% \begin{proof}
% Since $\|\hat{x}_0 - \bar{x}_0\| \leq \frac{1}{2}{\delta}$ and $\|\grad f_i(x_{i,k}) \| \leq L$, it follows from Lemma \ref{lem:stay-neighborhood} that for any $k>0$, the following holds:
% \begin{equation}\label{eq:neibohood3}
%     \begin{aligned}
%        &  \sum_{j=1}^n {W_{ij}} x_{j,k} -\alpha \grad f_i(x_{i,k})  \in U_{\Mcal} ({\delta}), ~ i = 1,\cdots,n.
%     \end{aligned}
% \end{equation}
% Based on the iterative scheme in Algorithm \ref{alg:drgta}, we have
% \[ \begin{aligned}
%     \|\bx_{k+1} - \bar{\bx}_{k+1}\| & \leq \| \bx_{k+1} - \bar{\bx}_k \| \\
%     & = \| \Pcal_{\Mcal^n}(\bW \bx_k - \alpha  \grad f(\bx_k)) - \Pcal_{\Mcal^n}( \hat{\bx}_k) \| \\
%     & \leq 2 \|\bW \bx_k - \alpha  \grad f(\bx_k) -\hat{\bx}_k  \| \\
%     & \leq 2 \| (W \otimes I_d) \bx_k - \hat{\bx}_k \| + 2\alpha\|\grad f(\bx_k)\| \\
%     & {\leq} 2 \| ((W -J ) \otimes I_d) (\bx_k - \hat{\bx}_k) \| + 2 \sqrt{n}\alpha   L \\
%     % & \leq 2 \| \bW \bx_k - \hat{\bx}_k \| + \alpha  L \\
%     & {\leq} 2 \sigma_2 \|\bx_k -\hat{\bx}_k\| + 2 \sqrt{n}\alpha L   \\
%     & \leq 2\sigma_2 \|\bx_k -\bar{\bx}_k\| + 2   \sqrt{n} \alpha L,
% \end{aligned}
% \]
% where the first inequality is from the optimality of $\bar{\bx}_{k+1}$, the second inequality utilizes \eqref{eq:neibohood3} and the $2$-Lispchitz continuity of $\Pcal$ over $U_{\Mcal}({\delta})$, the third inequality is due to the triangle inequality, the fourth inequality is from the boundedness of $\grad f_i(x_{i,k})$ over $\Mcal$ (i.e, $\|\grad f_i(x_{i,k})\|  \leq L$), and the fifth inequality comes from the assumptions on $W$. We complete the proof.
% \end{proof}

% % If $\alpha\geq L_g$, it holds that
% % \begin{equation}
% % \begin{aligned}
% %     \mathcal{L}_{\alpha}(\by) &\leq  \mathcal{L}_{\alpha,\alpha}^{\bx}(\by),\\
% %   \mathcal{L}_{\alpha,\alpha}^{\bx}(\bx)  & =  \mathcal{L}_{\alpha}(\bx), \\
% %   \grad \mathcal{L}_{\alpha,\alpha}^{\bx}(\bx)& = \grad \mathcal{L}_{\alpha}(\bx).
% % \end{aligned}
% % \end{equation}
% % When the step size $\alpha  = \alpha$ and $\alpha = \alpha^{-1}$, we have
% % \begin{equation}
% %     \begin{aligned}
% %     \bx_{k+1} &= \arg\min_{\bx\in\Mcal^n} \mathcal{L}_{\alpha,\alpha^{-1}}^{\bx_k}(\bx) \\
% %     & = \arg\min_{\bx\in\Mcal^n} \|\bx - \bx_k + \alpha \grad f(\bx_k) + (I-W)\bx_k \|^2 \\
% %     & = \arg\min_{\bx\in\Mcal^n} \|\bx  + \alpha \grad f(\bx_k) -W\bx_k \|^2\\
% %     & = \Pcal\left(W\bx_k - \alpha \grad f(\bx_k)  \right).
% %     \end{aligned}
% % \end{equation}

% With the above lemma, we can elaborate on a more detailed analysis of the consensus error.  The proof follows from \cite[Lemma D.3]{chen2021decentralized}.
% \begin{lemma} \label{lem:consensus}
%     Under Assumption \ref{assum-f}, let $\rho_t = 2\sigma_2$, assume that the sequence $\{\bx_k\}$ is generated by Algorithm \ref{alg:drgta}, $ t \geq \max\{\log_{\sigma_2}(1/2), \left\lceil \log_{\sigma_2}\left(\frac{{\delta}}{16\sqrt{n}\zeta}\right) \right \rceil \}$ and $\|\hat{x}_0 - \bar{x}_0\| \leq \frac{1}{2}{\delta}$,  $\alpha \leq {\delta}/(8L)$.  The following holds:
%     \be\label{eq:pgd:consensus:linear}
%     \|\bar{\bx}_{k+1} - \bx_{k+1} \| \leq  \rho_t^{k+1} \|  \bar{\bx}_0 - \bx_0\| + 2\sqrt{n}L \sum_{l=0}^k \rho_t^{k-l}\alpha_l.
%     \ee
%     Moreover, if the step size satisfy that $\alpha = \min\{ {\delta}/(8L), \frac{1}{(k+1)^p} \}$, where $0<p\leq 1$, then there exists a constant $C>0$ such that $\frac{1}{n}\|\bar{\bx}_k - \bx_k \|^2 \leq CL^2 \alpha^2 $, where $C$ is independent of $L$ and $n$.
% \end{lemma}

% \begin{proof}
%     It follows from Lemma \ref{lemma:dpg:consensus} that
%     \be \label{eq:temp:pgd:consensus}
%     \begin{aligned}
%         \|\bar{\bx}_{k+1} - \bx_{k+1} \| & \leq \rho_t \|\bar{\bx}_k - \bx_k\| + 2\sqrt{n}\alpha L \\
%         & \leq \rho_t^{k+1} \|  \bar{\bx}_0 - \bx_0\| + 2\sqrt{n}L \sum_{l=0}^k \rho_t^{k-l}\alpha_l,
%     \end{aligned}
%     \ee
%    which gives \eqref{eq:pgd:consensus:linear}. Let $a_k: = \frac{\|\bar{\bx}_k - \bx_k\|}{\sqrt{n} \alpha}$. For a given positive integer number $K\leq k$,  it follows from \eqref{eq:temp:pgd:consensus} that
%    \be
%    a_{k+1} \leq \rho_t a_k + 2L \frac{\alpha} 
%  {\alpha} \leq \rho_t^{k+1 - K}a_{K} + 2L \sum_{l=K}^k \rho_t^{k-l} \frac{\alpha_l}{\alpha_{l+1}}.
%    \ee
%    Since that $\alpha = \mathcal{O}(1/L)$ and $\|\bar{\bx}_0 - \bx_0 \| \leq \frac{1}{2} \sqrt{n}  {\delta}$, one have that $a_0 \leq \frac{1}{2}{\delta}/\alpha_0 = \mathcal{O}(D)$. Due to that $\lim_{k\rightarrow \infty} \frac{\alpha}{\alpha} = 1$, there exists sufficiently large $K$ such that $\alpha / \alpha \leq 2,~ \forall k\geq K$.  For $0\leq k \leq K$, there exists $C^{'} > 0$ such that $a_k^2 \leq C^{'} L^2$, where $C^{'}$ is independent of $L$ and $n$. For $k\geq K$, one has that $a_k^2 \leq CL^2$, where $C = 2C^{'} + \frac{32}{(1-\rho_t)^2}$. Hence, we get $\frac{\|\bar{\bx}_k - \bx_k\|^2}{n} \leq CL^2 \alpha^2$ for all $k\geq 0$, where $C = \mathcal{O}(\frac{1}{(1-\rho_t)^2})$.      
% \end{proof}

% By utilizing the Lipschitz-type inequalities on compact submanifolds in Section \ref{sec:ineq-man} and combining the above lemma, we can show a sufficient decrease on $f$. 
% \begin{lemma}\label{them:dpg}
% Under Assumption \ref{assum-f}, assume that the sequence $\{\bx_k\}$ is generated by Algorithm \ref{alg:drgta}, $ t \geq \max\{\log_{\sigma_2}(1/2), \left\lceil \log_{\sigma_2}\left(\frac{{\delta}}{16\sqrt{n}\zeta}\right) \right \rceil \}$ and $\|\hat{x}_0 - \bar{x}_0\| \leq \frac{1}{2}{\delta}$,  $\alpha \leq {\delta}/(8L)$.
% It holds that
% \be
% \begin{aligned}
%        f(\bar{x}_{k+1})  \leq & f(\bar{x}_k) - \frac{\alpha}{4} \|\grad f(\bar{x}_k) \|^2 +
%   \mathcal{G}_1   \alpha^3 + \mathcal{G}_2  \alpha^4,
% \end{aligned}
% \ee
% where
% \be
% \begin{aligned}
%     \mathcal{G}_1 & = (2C  + 48C^2(M_2^2+2(\sqrt{n}L_2+8Q)^2)  + 48C^2M_2^2   + 24Q^2)L^4,\\
%     \mathcal{G}_2 & = ( 32C^2(8Q + \sqrt{n}L_2 + M_2)^2 + 32C^2M_2^2 + 8Q^2  )L^5.
%    % C_3 & = \frac{\alpha L_G^2}{2}+\frac{36}{n}  +8(DQ+DM)+DM_2  +2M_2^2C_1 L^2(L_G + 2B)^2 \alpha^2.
% \end{aligned}
% \ee
% \end{lemma}

% \begin{proof}
% It follows from the Riemannian quadratic upper bound of $f$ in Lemma \ref{lemma:lipsctz} and $L_g \leq L$ that
% \be
% \label{eq:dgd:major:inequality}
% \begin{aligned}
% & f(\bar{x}_{k+1})  \leq f(\bar{x}_k) + \left<\grad f(\bar{x}_k),   \bar{x}_{k+1} - \bar{x}_{k}\right> + \frac{L}{2}\|\bar{x}_{k+1} - \bar{x}_{k} \|^2 \\
%  = & f(\bar{x}_k) - \left<\grad f(\bar{x}_k), \alpha \hat{g}_k \right> + \left<\grad f(\bar{x}_k),  \bar{x}_{k+1} - \bar{x}_{k} + \alpha \hat{g}_k \right> + \frac{L}{2}\|\bar{x}_{k+1} - \bar{x}_{k} \|^2 \\
% = & f(\bar{x}_k) - \frac{\alpha}{2}\|\grad f(\bar{x}_k) \|^2 - \frac{\alpha}{2}\| \hat{g}_k \|^2 + \frac{\alpha}{2}  \| \grad f(\bar{x}_k) - \hat{g}_k \|^2 \\
%   & +   \left<\grad f(\bar{x}_k),  \bar{x}_{k+1} - \bar{x}_{k} + \alpha \hat{g}_k \right> + \frac{L}{2} \|\bar{x}_{k+1} - \bar{x}_{k} \|^2.
%   % \leq  & f(\bar{\bx}_k) - \frac{\alpha}{4}\|\grad f(\bar{\bx}_k) \|^2 - \frac{\alpha}{2}\| \hat{\bf G}_k \|^2 + \frac{\alpha}{2} \underbrace{ \| \grad f(\bar{\bx}_k) - \hat{\bf G}_k \|^2}_{a_1} \\
%   % & +    \frac{1}{\beta} \underbrace{\| \bar{\bx}_{k+1} - \bar{\bx}_{k} + \alpha \hat{\bf G}_k \|^2 }_{a_2} + \frac{L}{2} \underbrace{\|\bar{\bx}_{k+1} - \bar{\bx}_{k} \|^2}_{a_3},
% \end{aligned}
% \ee
% By Young's inequality, we have
% \be \label{eq:grad-prod}
% \left<\grad f(\bar{x}_k),  \bar{x}_{k+1} - \bar{x}_{k} + \alpha \hat{g}_k \right> \leq \frac{\alpha}{4}\|\grad f(\bar{x}_k) \|^2 + \frac{1}{\alpha}\| \bar{x}_{k+1} - \bar{x}_{k} + \alpha \hat{g}_k\|^2.
% \ee
% Combining \eqref{eq:dgd:major:inequality} and  \eqref{eq:grad-prod} leads to
% \be\label{eq:dgd:major:inequality1}
% \begin{aligned}
%  f(\bar{x}_{k+1})  \leq &  f(\bar{x}_k) - \frac{\alpha}{4}\|\grad f(\bar{x}_k) \|^2 - \frac{\alpha}{2}\| \hat{g}_k \|^2 + \frac{\alpha}{2} \underbrace{ \| \grad f(\bar{x}_k) - \hat{g}_k \|^2}_{a_1} \\
%   & +  \frac{1}{\alpha} \underbrace{\| \bar{x}_{k+1} - \bar{x}_{k} + \alpha \hat{g}_k\|^2}_{a_2} + \frac{L}{2} \underbrace{\|\bar{x}_{k+1} - \bar{x}_{k} \|^2.}_{a_3}
% \end{aligned}
% \ee
% Now, let us bound $a_1$, $a_2$, and $a_3$, respectively. Apply Lemma \ref{lemma:lipsctz} yields
% \be
% \begin{aligned}
% a_1 & =  \left\|\grad f(\bar{x}_k) - \frac{1}{n}\sum_{i=1}^n \grad f_i(x_{i,k}) \right \|^2 \\
% & \leq  \frac{1}{n}\sum_{i=1}^n\|\grad f_i(\bar{x}_k) -  \grad f_i(x_{i,k}) \|^2 \\
% & \leq    \frac{L^2}{n} \sum_{i=1}^n \|\bar{x}_k -  x_{i,k} \|^2 = \frac{L^2}{n} \| \bar{\bx}_k - \bx_k \|^2.
% \end{aligned}
% \ee
% For $a_2$, it follows from the triangle inequality that
% \be\label{eq:tmp2}
% \begin{aligned}
%     & \| \bar{x}_{k+1} - \bar{x}_{k} + \alpha \hat{g}_k\| \\
%    \leq & \| \bar{x}_k - \hat{x}_k \| + \| \bar{x}_{k+1} - \hat{x}_{k+1} \| + \| \hat{x}_{k+1} - \hat{x}_{k} + \alpha \hat{g}_k\| \\
%    \overset{\eqref{eq:distance-an-rm}}{\leq} &  \frac{M_2}{n}(\| \bar{\bx}_k - \bx_k \|^2 + \| \bar{\bx}_{k+1} - \bx_{k+1} \|^2) + \| \hat{x}_{k+1} - \hat{x}_{k} + \alpha \hat{g}_k\|.
% \end{aligned}
% \ee
% Moreover, we have
% \be\label{eq:tmp1}
% \begin{aligned}
%   &  \| \hat{x}_{k+1} - \hat{x}_{k} + \alpha \hat{g}_k\| \\
%   = & \| \frac{1}{n}\sum_{i=1}^n (x_{i,k+1} - x_{i,k} + \alpha \grad f_i(x_{i,k}) ) \|  \\
%   \leq  &\frac{1}{n} \|  \sum_{i=1}^n (x_{i,k+1} - x_{i,k} + \alpha \grad f_i(x_{i,k}) + \grad \phi_i( \bx_k )) \| + \frac{1}{n}\|  \sum_{i=1}^n \grad \phi_i( \bx_k ) \| \\
%   \overset{\eqref{projec-second-order1} }{\leq} & \frac{Q}{n} \sum_{i=1}^n  \| \alpha \grad f_i(x_{i,k}) + \nabla \phi_i( \bx_k )  \|^2 + \frac{1}{n}\|  \sum_{i=1}^n \grad \phi_i( \bx_k ) \| \\
%   \overset{\eqref{eq:sum-consen-grad}}{\leq} & \frac{2Q\alpha^2}{n} \|{\bf G}_k \|^2 + \frac{2Q}{n} \|\nabla \phi( \bx_k ) \|^2  +\frac{L_2}{\sqrt{n}} \| \bx_k - \bar{\bx}_k \|^2 \\
%   \leq & \frac{2 Q\alpha^2}{n} \|{\bf G}_k \|^2 + \frac{(\sqrt{n}L_2 + 8Q)}{n}\| \bx_k - \bar{\bx}_k \|^2.
% \end{aligned}
% \ee
% Plugging \eqref{eq:tmp1} into \eqref{eq:tmp2} gives
% \be
% \begin{aligned}
%     a_2  \leq \frac{3M_2^2+6(\sqrt{n}L_2+8Q)^2}{n^2}\| \bx_k - \bar{\bx}_k \|^4 + \frac{3M_2^2}{n^2}\| \bar{\bx}_{k+1} - \bx_{k+1} \|^4 + \frac{24Q^2 \alpha^4}{n^2} \|{\bf G}_k \|^4.
% \end{aligned}
% \ee
% By \eqref{eq:diff-avg-x}, we obtain
% \begin{equation}
%     \begin{aligned}
%     a_3  \leq & \frac{4(8Q + \sqrt{n}L_2 + M_2)^2}{n^2} \|\bx_k - \bar{\bx}_k\|^4 +\frac{16Q^2\alpha^4}{n^2} \| {\bf G}_k \|^4  + 4\alpha^2 \| \hat{g}_{k}  \|^2   \\
%     & + \frac{4M_2^2}{n^2} \| \bx_{k+1}  - \bar{\bx}_{k+1} \|^4\\
%   %& \leq \left( \frac{36}{n} \|\bx_k - \bar{\bx}_k\|^2 + \frac{16\alpha^2}{n}\| \grad f(\bx_k) \|^2  \right) + 8M_2^2CL^2  ( \alpha^2 \| \bx_k  - \bar{\bx}_k \|^2 +\alpha^2 \| \bx_{k+1}  - \bar{\bx}_{k+1} \|^2 )
%     \end{aligned}
% \end{equation}
%  Combining $a_1, a_2, a_3$ with \eqref{eq:dgd:major:inequality1} and noting $\|{\bf G}_k\|^2 \leq n L^2$ imply that
% \be
% \begin{aligned}
%  f(\bar{x}_{k+1})  \leq &  f(\bar{x}_k) - \frac{\alpha}{4}\|\grad f(\bar{x}_k) \|^2 - \frac{\alpha}{2}\| \hat{g}_k \|^2 + \frac{\alpha}{2} a_1  +  \frac{1}{\alpha} a_2 + \frac{L}{2} a_3\\
%  \leq & f(\bar{x}_k) - \frac{\alpha}{4}\|\grad f(\bar{x}_k) \|^2 - ( \frac{\alpha}{2} - 2L\alpha^2  )\|  \hat{g}_k\|^2 + \frac{L^2\alpha}{2n} \| \bar{\bx}_k - \bx_k \|^2  \\
%  & +(\frac{3M_2^2+6(\sqrt{n}L_2+8Q)^2}{n^2\alpha}+\frac{2L(8Q + \sqrt{n}L_2 + M_2)^2}{n^2} ) \| \bx_k - \bar{\bx}_k \|^4 \\
%  & +( \frac{3M_2^2}{n^2\alpha} +  \frac{2LM_2^2}{n^2} )\| \bar{\bx}_{k+1} - \bx_{k+1} \|^4 + (24Q^2\alpha^3 + 8Q^2L  \alpha^4) L^4
% \end{aligned}
% \ee
% % Let  $\nabla \phi_i(\bx_k) + \alpha \grad f_i(x_{i,k}) = d_1 + u_2, $
% % where $d_1 = \Pcal_{T_{x_{i,k}}\Mcal}(\nabla \phi_i(\bx_k) + \alpha \grad f_i(x_{i,k})) , u_2 = \nabla \phi_i(\bx_k)+ \alpha \grad f_i(x_{i,k}) - d_1$.  Therefore,
% % \be
% % \begin{aligned}
% % &\left< \grad f_j(\bar{x}_k), x_{i,k+1} - x_{i,k} + \alpha \grad f_i(x_{i,k}) + \nabla \phi_i(\bx_k) \right>  \\
% % = &  \left< \grad f_j(\bar{x}_k),  \Pcal_{\Mcal}(x_{i,k} - d_1 - u_2  ) - [x_{i,k}  - d_1 - u_2   ]\right> \\
% %  = &   \left< \grad f_j(\bar{x}_k),  \Pcal_{\Mcal}(x_{i,k} - d_1 - u_2  ) - \Pcal_{\Mcal}(x_{i,k} - d_1  )\right>  \\
% %  &+   \left< \grad f_j(\bar{x}_k),  \Pcal_{\Mcal}(x_{i,k} - d_1   ) - [x_{i,k} - d_1 ]\right> + \left< \grad f_j(\bar{x}_k),  u_2\right>    \\
% %  \leq  & LQ\|d\|^2 + LM\|d_1\|^2 +\left< \grad f_j(\bar{x}_k),  u_2\right>
% % \\
% % \leq & (LQ+LM) \| \nabla \phi_i(\bx_k) + \alpha \grad f_i(x_{i,k}) \|^2 +\left< \grad f_j(\bar{x}_k) - \grad f_j(x_{i,k}),  u_2\right> \\
% % \leq & (LQ+LM) \| \nabla \phi_i(\bx_k) + \alpha \grad f_i(x_{i,k}) \|^2 + \frac{1}{2} \| \grad f_j(\bar{x}_k) - \grad f_j(x_{i,k})\|^2 + \frac{1}{2}\|  u_2\|^2 \\
% % \leq & (LQ+LM+1) \| \nabla \phi_i(\bx_k) + \alpha \grad f_i(x_{i,k}) \|^2 + L_G^2 \| \bar{x}_k - x_{i,k}\|^2
% % % \leq & 2(LQ+LM) \left(\| \nabla \phi_i(\bx_k)\|^2 + \alpha^2 \| \grad f_i(x_{i,k})\|^2  \right) \\
% % % \leq & 2(LQ+LM) \left(\| \nabla \phi_i(\bx_k)\|^2 + \alpha^2 \| \grad f_i(x_{i,k})\|^2  \right),
% % \end{aligned}
% % \ee
% % where the first inequality utilizes Lemma  \ref{lemma-project} , \eqref{eq:second-order}, the second inequality use that $\|d_1\| \leq \|d\|$, $u_2\in N_{x_{i,k}}\Mcal$ and $\left< \grad f_j(\bar{x}_k),  u_2\right> = 0$.
% % \be
% % \begin{aligned}
% %     a_2 =  & \sum_{j=1}^n  \left< \grad f_j(\bar{x}_k), \hat{x}_{k+1} - \hat{x}_k + \alpha \hat{g}_k   \right> + \sum_{j=1}^n \left< \grad f_j(\bar{x}_k), \bar{x}_{k+1} -\hat{x}_{k+1}  +  \bar{x}_{k} -\hat{x}_{k} \right> \\
% %     = & \sum_{j=1}^n \left< \grad f_j(\bar{x}_k), \frac{1}{n} \sum_{i=1}^n \left(x_{i,k+1} - x_{i,k} + \alpha \grad f_i(x_{i,k}) + \nabla \phi_i(\bx_k) \right)  \right>  \\
% %     & + \sum_{j=1}^n \left< \grad f_j(\bar{x}_k), \bar{x}_{k+1} -\hat{x}_{k+1}  +  \bar{x}_{k} -\hat{x}_{k} \right> \\
% %     \leq & \sum_{i=1}^n \left((LQ+LM+1) \| \nabla \phi_i(\bx_k) + \alpha \grad f_i(x_{i,k}) \|^2 + L_G^2 \| \bar{x}_k - x_{i,k}\|^2\right) \\
% %     &+ nL(\| \bar{x}_{k+1} -\hat{x}_{k+1}\| + \|\bar{x}_{k} -\hat{x}_{k}\|) \\
% %     % \leq & 2(LQ+LM+1) \left( \|(I-W)\bx_k\|^2 + \alpha^2 \|\grad f(\bx_k) \|^2 \right) + L_G^2 \|\bar{x}^k - \bx_k\|^2 \\
% %     % &+ L M_2(\| \bar{\bx}_{k+1} -\bx_{k+1}\|^2 + \|\bar{\bx}_{k} -\bx_{k}\|^2) \\
% %       \leq & 2(LQ+LM+1) \left((1-\lambda_n(W))^2 \|\bar{\bx}_k - \bx_k\|^2 + \alpha^2 \|\grad f(\bx_k) \|^2 \right) + L_G^2 \|\bar{x}^k - \bx_k\|^2 \\
% %     &+ L M_2(\| \bar{\bx}_{k+1} -\bx_{k+1}\|^2 + \|\bar{\bx}_{k} -\bx_{k}\|^2) \\
% %       \leq & 2(LQ+LM+1) \alpha^2 \|\grad f(\bx_k) \|^2 + LM_2 \| \bar{\bx}_{k+1} -\bx_{k+1}\|^2 \\
% %       & + (8(LQ+LM+1) + L_G^2 + LM_2) \|\bar{\bx}_k - \bx_k\|^2
% %     % \leq & \sum_{i=1}^n 2(DQ+DM) \left(\|(I-W)\bx_k\|^2 + \alpha^2 \| \grad f_i(x_{i,k})\|^2  \right) + nD(\| \bar{x}_{k+1} -\hat{x}_{k+1}\| + \|\bar{x}_{k} -\hat{x}_{k}\|) \\
% %     % \leq & 2(DQ+DM) \left(   (1-\lambda_n(W))^2 \| \bx_k - \bar{\bx}_k\|^2+ \alpha^2 \|\grad f(\bx_k) \|^2 \right) \\
% %     % & + D M_2(\| \bar{\bx}_{k+1} -\bx_{k+1}\|^2 + \|\bar{\bx}_{k} -\bx_{k}\|^2) \\
% % \end{aligned}
% % \ee
% % % For the third term, we have
% % % \be
% % % \begin{aligned}
% % %     \| \hat{\bx}_{k+1} - \hat{\bx}_k + \alpha \hat{\bf G}_k \| & \leq  \| \hat{\bx}_{k+1} - \hat{\bx}_k + \alpha \hat{\bf G}_k + (I-W)\bx_k \|  + \| (I - W)\bx_k \| \\
% % %     & \leq n \| \frac{1}{n}\sum_{i=1}^n\left(x_{i,k+1} - x_{i,k} + \alpha \grad f_i(x_{i,k}) + \nabla \phi_i(\bx_k)\right) \|+ \| (I - W)\bx_k \|\\
% % %     & \leq \sum_{i=1}^n \| x_{i,k+1} - x_{i,k} + \alpha \grad f_i(x_{i,k}) + \nabla \phi_i(\bx_k)\|+ \| (I - W)\bx_k \| \\
% % %     & \leq \sum_{i=1}^n \|  \alpha \grad f_i(x_{i,k}) + \nabla \phi_i(\bx_k)\|+ \| (I - W)\bx_k \| \\
% % %     &\leq \alpha \| \grad f(\bx_k) \| + 2  \| (I - W)\bx_k \|.
% % % \end{aligned}
% % % \ee
% % By \eqref{eq:distance-an-rm}, we obtain
% % \begin{equation}
% %     \begin{aligned}
% %     a_3 & \leq \left( \frac{3\sqrt{n}}{n} (1-\lambda_n(W))\|\bx_k - \bar{\bx}_k\| + \frac{4\alpha \sqrt{n}}{n}\| \grad f(\bx_k) \|  \right)^2 + 2 (  \| \hat{\bx}_k  - \bar{\bx}_k \|^2 + \| \hat{\bx}_{k+1}  - \bar{\bx}_{k+1} \|^2 ) \\
% %   &  \leq \left( \frac{36}{n} \|\bx_k - \bar{\bx}_k\|^2 + \frac{16\alpha^2}{n}\| \grad f(\bx_k) \|^2  \right) + \frac{2M_2^2 }{n^2} (  \| \bx_k  - \bar{\bx}_k \|^4 + \| \bx_{k+1}  - \bar{\bx}_{k+1} \|^4 ) \\
% %   %& \leq \left( \frac{36}{n} \|\bx_k - \bar{\bx}_k\|^2 + \frac{16\alpha^2}{n}\| \grad f(\bx_k) \|^2  \right) + 8M_2^2CL^2  ( \alpha^2 \| \bx_k  - \bar{\bx}_k \|^2 +\alpha^2 \| \bx_{k+1}  - \bar{\bx}_{k+1} \|^2 )
% %     \end{aligned}
% % \end{equation}
% % Therefore, by
% % combining $a_1,a_2,a_3$ with \eqref{eq:dgd:major:inequality}  implies that
% % \be
% % \begin{aligned}
% %     f(\bar{\bx}_{k+1})  \leq & f(\bar{\bx}_k) - \frac{\alpha}{2}\|\grad f(\bar{\bx}_k) \|^2 - \frac{\alpha}{2}\| \hat{\bf G}_k \|^2 + \left(
% %  2(LQ+LM+1) + \frac{16}{n}\right)\alpha^2\| \grad f(\bx_k) \|^2  \\
% %     & +  (8(LQ+LM+1) + 2L_G^2 + LM_2 + \frac{36}{n})\| \bx_k  - \bar{\bx}_k \|^2  + LM_2\| \bx_{k+1}  - \bar{\bx}_{k+1} \|^2 \\
% %     & +\frac{2M_2^2 }{n^2} (  \| \bx_k  - \bar{\bx}_k \|^4 + \| \bx_{k+1}  - \bar{\bx}_{k+1} \|^4 ).
% % \end{aligned}
% % \ee
% % Note that
% % \be
% % \| \grad f(\bx_k) \| \leq \| \grad f(\bar{\bx}_k) - \grad f(\bx) \| + \| \grad f(\bar{\bx}_k) \| \leq L_G\| \bx_k - \bar{\bx}_k\| + \|\grad f(\bar{\bx}_k) \|.
% % \ee
% Since $\alpha \leq \alpha \leq 1/(4L)$ and $\frac{1}{n}\|\bar{\bx}_k - \bx_k \|^2 \leq CL^2 \alpha^2 $ by Lemma \ref{lem:consensus}, it holds that
% \be
% \begin{aligned}
%     f(\bar{x}_{k+1})  \leq & f(\bar{x}_k) - \frac{\alpha}{4} \|\grad f(\bar{x}_k) \|^2 \\
%     & + (2C  + 48C^2(M_2^2+2(\sqrt{n}L_2+8Q)^2)  + 48C^2M_2^2   + 24Q^2)L^4\alpha^3 \\
%     & + ( 32C^2(8Q + \sqrt{n}L_2 + M_2)^2 + 32C^2M_2^2 + 8Q^2  )L^5\alpha^4.
% \end{aligned}
% \ee
% The proof is completed.
% \end{proof}


% \begin{theorem}
%      Under Assumption \ref{assum-f}, assume that the sequence $\{\bx_k\}$ is generated by Algorithm \ref{alg:drgta}, $ t \geq \max\{\log_{\sigma_2}(1/2), \left\lceil \log_{\sigma_2}\left(\frac{{\delta}}{16\sqrt{n}\zeta}\right) \right \rceil \}$ and $\|\hat{x}_0 - \bar{x}_0\| \leq \frac{1}{2}{\delta}$,  the step size satisfy that
%      $$\alpha  = \frac{1}{\sqrt{k+1}} \min\{{\delta}/(8L),1 \}.$$
%      The following holds:
%      \be
%   \min_{k=0,\cdots,K}\|\grad f(\bar{x}_k) \|^2 = \mathcal{O}( \frac{1}{\sqrt{K+1}}).
%      \ee
% \end{theorem}
% \begin{proof} By Lemma \ref{them:dpg}, we have
%     \be
% \begin{aligned}
%        f(\bar{x}_{k+1})  \leq & f(\bar{x}_k) - \frac{\alpha}{4}  \|\grad f(\bar{x}_k) \|^2  +
%  \mathcal{G}_1 \alpha^3 +\mathcal{G}_2 \alpha^4.
% \end{aligned}
% \ee
% Summing the above inequality over $k=0,1,\ldots, K$ and telescoping the right-hand side with any $K>0$ give
% \be
% \sum_{k=0}^K \frac{\alpha}{4}  \|\grad f(\bar{x}_k) \|^2 \leq f(\bar{x}_{0})  - f(\bar{x}_*) +  \mathcal{G}_1 \sum_{k=0}^K \alpha^3 + \mathcal{G}_2\sum_{k=0}^K \alpha^4.
% \ee
% Dividing both sides by $\sum_{k=0}^K \frac{\alpha}{4}$ yields
% \be
% \min_{k=0,\cdots,K}\|\grad f(\bar{x}_k) \|^2 \leq \frac{f(\bar{x}_{0})  - f(\bar{x}_*) +  \mathcal{G}_1 \sum_{k=0}^K \alpha^3 + \mathcal{G}_2\sum_{k=0}^K \alpha^4}{\sum_{k=0}^K \frac{\alpha}{4}}.
% \ee
% Note that $\frac{\sum_{k=0}^K \alpha^3 }{\sum_{k=0}^K \alpha } = \mathcal{O}( \frac{1}{\sqrt{K+1}})$.  The proof is completed.
% \end{proof}

% \begin{remark}
%     The critical difference between our results and those of \cite{zeng2018nonconvex} lies in the following two aspects. Firstly, the domain of each $f_i$ is restricted to a nonconvex manifold. By taking advantage of the singleton and contraction-like properties of the projection operator, we show the Euclidean consensus is sufficient to guarantee convergence. Secondly, different from Euclidean gradient-based upper bound used in \cite{zeng2018nonconvex}, we derive a new quadratic upper bound for $f_i$'s by utilizing the Riemannian geometry, which brings lots of savings in the total number of iterations of the algorithms as shown in our numerical experiments. 

%     In \cite{chen2021decentralized}, the authors analyzed the $\mathcal{O}(\log(K)/\sqrt{K})$ iteration complexity of the decentralized Riemannian gradient descent under vanishing step sizes. Still, they will obtain the same complexity results as ours result in the deterministic setting. It should be emphasized that our research is focused on problems in general compact submanifolds, while their analysis relies on the properties specific to the Stiefel manifold. Additionally, we utilize Euclidean consensus, while they use Riemannian consensus.




    
% \end{remark}


% %\section{Decentralized projected gradient method}
% % \begin{proposition}\label{proposi:consensus}
% % For each iteration $k$, define $\hat{\bx}_k = \frac{1}{n} \mathbf{1}\mathbf{1}^{\top}\bx_k$ and let $\bar{\bx}_k = \Pcal_{\Mcal}(\hat{\bx}_k)$. Then, it holds that
% % \begin{equation}
% %     \|\bx_k - \bar{\bx}_k\|\leq 2C\zeta^k\|\bx_0\| +\frac{4\alpha C \sqrt{N}L}{1-\zeta}.
% % \end{equation}
% % \end{proposition}
% % \begin{proof}
% % Let $e^k = \mathcal{P}( W\bx_k - \alpha \grad f(\bx_k)) -W\bx_k$. The iteration \eqref{eq:dpg} can be rewritten as
% % \begin{equation}
% % \begin{aligned}
% % \bx_{k+1} &= \mathcal{P}( W\bx_k - \alpha \grad f(\bx_k)) \\
% % & = W\bx_k + e^k.
% % \end{aligned}
% % \end{equation}

% % \[ \begin{aligned}
% %  \| {\bx}_{k+1} - \hat{\bx}_{k+1} \|  \leq & \| {\bx}_{k+1} - \hat{\bx}_{k} \| \\
% % \leq & \| W\bx_k - \hat{\bx}_{k} \| + 2 \alpha B + (1 - \lambda_n) \| \bx_k - \hat{\bx}_k \| \\
% % \leq & \sigma_2 \| \bx_k - \hat{\bx}_{k} \| + 2 \alpha B + (1 - \lambda_n) \| \bx_k - \hat{\bx}_k \|
% % \end{aligned}
% %  \]
% % The term $e^k$ can be bounded by
% % \begin{equation}
% % \begin{aligned}
% %  \|e^k\|
% % = & \|\mathcal{P}( W\bx_k - \alpha \grad f(\bx_k)) -W\bx_k + \alpha \grad f(\bx_k) - \alpha \grad f(\bx_k) \|  \\
% % \leq &  \|\bx_k -  W\bx_k + \alpha \grad f(\bx_k) \| + \alpha \|\grad f(\bx_k)\| \\
% % \leq & \|\bx_k - W\bx_k\| + 2\alpha \|\grad f(\bx_k)\|.
% % \end{aligned}
% % \end{equation}
% % Note that
% % \begin{equation}
% %     \begin{aligned}
% %     \bx_{k} - \hat{\bx}_{k} & =(I -  \frac{1}{n}\mathbf{1}\mathbf{1}^{\top})\bx_k  = \left( W^k - \frac{1}{n}\mathbf{1}\mathbf{1}^{\top}  \right)\bx_0  +  \sum_{j=0}^{k-1}\left(W^{k-1-j} - \frac{1}{n}\mathbf{1}\mathbf{1}^{\top}\right)e^j.
% %     \end{aligned}
% % \end{equation}
% % It follows from Lemma \ref{lemma:Wapprox} that
% % \begin{equation}\label{eq:hatx}
% % \begin{aligned}
% % \|\bx_{k} - \hat{\bx}_{k}\| & \leq C\zeta^k\|\bx_0\| + 2\alpha C \sqrt{N}L\sum_{j=0}^{k-1}\zeta^{k-1-j} \\
% % & \leq C\zeta^k\|\bx_0\| +\frac{2\alpha C \sqrt{N}L}{1-\zeta}.
% % \end{aligned}
% %     \end{equation}
% % By the definition of $\bar{\bx}$ and \eqref{eq:hatx}, we conclude that
% % \begin{equation}
% %     \begin{aligned}
% %     \|\bx_k -\bar{\bx}_k \| & \leq \|\bx_k -\hat{\bx}_k \| + \|\hat{ \bx}^k -\bar{\bx}_k \| \\
% %     & \leq 2\|\bx_k -\hat{\bx}_k \| \\
% %     & \leq 2C\zeta^k\|\bx_0\| +\frac{4\alpha C \sqrt{N}L}{1-\zeta}.
% %     \end{aligned}
% % \end{equation}
% % \end{proof}

% % \begin{lemma}[Sufficient reduction of $\mathcal{L}_{\alpha}(\bx_k)$]\label{lemma:reduction}
% % Assume that $\lambda_n(W)>0$ and $0<\alpha<\frac{\lambda_n(W)}{L}$. Then
% % \begin{equation}
% %     \mathcal{L}_{\alpha}(\bx_{k+1}) \leq \mathcal{L}_{\alpha}(\bx_k)-\frac{1}{2}(\alpha^{-1} \lambda_n(W)-L)\|\bx_{k+1} - \bx_k\|^2.
% % \end{equation}
% % \end{lemma}
% % \begin{proof}
% % Since $\sum_{i=1}^n \nabla f_i(x_i)$ is $L$-Lipschitz, where $L = \max_i\{L_i\}$. $\mathcal{L}_{\alpha}$ is Lipschitz with the constant $L_g=L + \alpha^{-1}\lambda_{\max}(I-W)=L + \alpha^{-1}(1-\lambda_n(W))$. Then
% % \begin{equation}
% %     \begin{aligned}
% %     \mathcal{L}_{\alpha}(\bx_{k+1}) \leq  & \mathcal{L}_{\alpha,1/L_g}^{\bx_k}(\bx_{k+1}) \\  =&  \mathcal{L}_{\alpha}(\bx_k) + \left<\grad f(\bx_k)+\alpha^{-1}(I-W)\bx_k,\bx_{k+1}-\bx_k  \right>\\
% %     &+ \frac{L_g}{2}\|\bx_{k+1}-\bx_k\|^2.
% %     \end{aligned}
% % \end{equation}
% % Let $d^k = \grad f(\bx_k)+\alpha^{-1}(I-W)\bx_k$. One finds that $\bx_{k+1} = \mathcal{P}_{\mathcal{M}}(\bx_k - \alpha d^k)$, it follow that
% % \begin{equation}
% %     \begin{aligned}
% %   0 & \leq   \|\bx_{k} - (\bx_k - \alpha d^k)\|^2 - \|\bx_{k+1} - (\bx_k - \alpha d^k)\|^2 \\
% %  & = -\|\bx_{k+1} - \bx_k\|^2 - 2\alpha \left< d^k,\bx_{k+1} - \bx_k \right>
% %     \end{aligned}
% % \end{equation}
% % Combining the above two inequality yield
% % \begin{equation}
% %     \begin{aligned}
% %     \mathcal{L}_{\alpha}(\bx_{k+1}) - \mathcal{L}_{\alpha}(\bx_k) & \leq -\left(\frac{1}{2\alpha} - \frac{L_g}{2}\right)\|\bx_{k+1} - \bx_k\|^2 \\
% %     & \leq -\frac{1}{2}(\alpha^{-1} \lambda_n(W)-L)\|\bx_{k+1} - \bx_k\|^2.
% %     \end{aligned}
% % \end{equation}

% % \end{proof}
% % \begin{lemma}[Bounded subgradient]\label{lemma:bounded}
% % Let $\xi^k\in N_{\bx_k}\mathcal{M}$. Then
% % \begin{equation}
% %     \|\nabla \mathcal{L}_{\alpha}(\bx_{k+1}) + \xi^{k+1}\| \leq (\alpha^{-1}(2-\lambda_n(W)) +L)\|\bx_{k+1} - \bx_k\|
% % \end{equation}
% % \end{lemma}


% % \begin{proof}
% % The following optimality condition holds
% % \begin{equation}
% %     0\in \alpha^{-1}(\bx_{k+1} - \bx_k + \alpha \nabla \mathcal{L}_{\alpha}(\bx_k)) + N_{\bx}\mathcal{M},
% % \end{equation}
% % which yields
% % \begin{equation}
% %     \begin{aligned}
% %      \| \nabla \mathcal{L}_{\alpha}(\bx_{k+1}) + \xi^{k+1} \| & \leq \alpha^{-1}\|\bx_{k+1} - \bx_k\| + \|\nabla \mathcal{L}_{\alpha}(\bx_{k+1}) - \nabla \mathcal{L}_{\alpha}(\bx_k)\| \\
% %     & \leq (\alpha^{-1}(2-\lambda_n(W)) +L)\|\bx_{k+1} - \bx_k\|.
% %     \end{aligned}
% % \end{equation}
% % \end{proof}
% % Combining Proposition \ref{proposi:consensus}, Lemma \ref{lemma:reduction} and Lemma \ref{lemma:bounded} yield the main result:
% % \begin{theorem}
% % Let $\bx_k$ be the sequence generated by DPG \eqref{eq:dpg} with the step size $0< \alpha < \frac{\lambda_n(W)}{L}$. Let Assumption \eqref{assum-f} holds. Then $\{\bx_k\}$ has at least one accumulation point $\bx_*$, and any such point  is a stationary point of $\mathcal{L}_{\alpha}(\bx)$. Furthermore, the running best rates of the sequences $\{\|\bx_{k+1}-\bx_k\|^2\}$, and $\{\|\nabla \mathcal{L}_{\alpha}(\bx_k)\|^2\}$, and $\|\frac{1}{n}\mathrm{1}^{\top}\grad f(\bx_k)\|^2$ are $o(\frac{1}{k})$. The convergence rate of the sequence $\{\frac{1}{K}\sum_{k=0}^{K-1}\|\frac{1}{n}\mathrm{1}^{\top}\grad f(\bx_k)\|^2  \}$ is $\mathcal{O}(\frac{1}{K})$.
% % \end{theorem}
% % \begin{proof}
% % By \eqref{equ:Lyap}, $\nabla \mathcal{L}_{\alpha,\alpha^{-1}}^{\bx_k}(\bx_{k+1}) = \frac{1}{\alpha}(\bx_{k+1}-\bx_k + \alpha \grad f(\bx_k) + (I-W)\bx_k)$. By the update rule of $\bx_{k+1}$, one concludes that
% % \begin{equation}
% % \begin{aligned}
% % 0 =   & \grad \mathcal{L}_{\alpha,\alpha^{-1}}^{\bx_k}(\bx_{k+1}) \\
% %   = & \Pcal_{T_{\bx_k}\Mcal^n}(\nabla \mathcal{L}_{\alpha,\alpha^{-1}}^{\bx_k}(\bx_{k+1})) \\
% %     =& \grad f(\bx_k) + \Pcal_{T_{\bx_k}\Mcal^n}(\frac{1}{\alpha}(\bx_{k+1} - \bx_k + (I-W)\bx_k))\\
% %    = & \grad  \mathcal{L}_{\alpha}(\bx_{k}) +  \Pcal_{T_{\bx_k}\Mcal^n}(\frac{1}{\alpha}(\bx_{k+1} - \bx_k))
% % \end{aligned}
% % \end{equation}
% % Thus, it follows that
% % \begin{equation}
% %     \| \grad \mathcal{L}_{\alpha}(\bx_{k}) \| \leq \alpha^{-1}\|\bx_{k+1} - \bx_{k}\|.
% % \end{equation}
% % By \eqref{equ:Lyap}, $\grad  \mathcal{L}_{\alpha}(\bx_{k}) = \grad f(\bx_k) + \Pcal_{T_{\bx_k}\Mcal^n} (\alpha^{-1}(I-W)\bx_k) $, which implies $\frac{1}{n}\mathrm{1}^{\top}(\grad f(\bx_k) + \Pcal_{T_{\bx_k}\Mcal^n} (\alpha^{-1}(I-W)\bx_k) ) = \frac{1}{n} \mathrm{1}^{\top}\grad\mathcal{L}_{\alpha}(\bx_{k})$. Thus,
% % \begin{equation*}
% % \begin{aligned}
% %   \left\| \frac{1}{n}\mathrm{1}^{\top}\grad f(\bx_k)  \right\|^2
% %  & = \left\|\frac{1}{n} \mathrm{1}^{\top} (\grad\mathcal{L}_{\alpha}(\bx_{k}) + \Pcal_{T_{\bx_k}\Mcal^n} (\alpha^{-1}(I-W)\bx_k))  \right\|^2 \\
% %  & \leq \left\|  \grad\mathcal{L}_{\alpha}(\bx_{k})  \right\|^2 + \alpha^{-2} \left\|  (I-W)\bx_k) \right\|^2 \\
% %  &\leq \left\|  \grad\mathcal{L}_{\alpha}(\bx_{k})  \right\|^2 + 2\alpha^{-2}( \left\|  \bx_k - \hat{\bx}_k \right\|^2 +  \left\|   \hat{\bx}_k - W\bx_k \right\|^2) \\
% %   &\leq \left\|  \grad\mathcal{L}_{\alpha}(\bx_{k})  \right\|^2 + 2\alpha^{-2}(1+C^2\zeta^2) \left\|  \bx_k - \hat{\bx}_k \right\|^2  \\
% % \end{aligned}
% % \end{equation*}
% % \end{proof}
% % \subsection{Convergence results with decreasing step size}
% % In this section, we discuss DPG with deceasing step sizes. In particular, we take the step size as follows:
% % \begin{equation}\label{stepzise-decrasing}
% %     \alpha = \frac{1}{L_f (k+1)^{\epsilon}}, ~0\leq \epsilon \leq 1.
% % \end{equation}

% % \begin{proposition}
% % Let DPG use step size in \eqref{stepzise-decrasing}. Then $\|\bx_k - \bar{\bx}_k\|$ convergence to 0 at the rate $\mathcal{O}(\frac{1}{(k+1)^{\epsilon}})$.
% % \end{proposition}

% % \begin{proof}
% %  By Proposition \ref{proposi:consensus}, $\|e^k\| \leq \leq 2\alpha\sqrt{N} L$ and
% %  and
% % \begin{equation}
% % \begin{aligned}
% % \|\bx_{k} - \hat{\bx}_{k}\| & \leq C\zeta^k\|\bx_0\| + 2 C \sqrt{N}L\sum_{j=0}^{k-1}\alpha_j\zeta^{k-1-j}.
% % \end{aligned}
% %     \end{equation}
% % Furthermore, by the step sizes \ref{stepzise-decrasing}, we get $\lim_{k\rightarrow \infty}\|\bx_k - \bar{\bx}_k\| = 0$.
% % \end{proof}
% % In particular, we follow the step blow
% % \begin{itemize}
% %     \item We establish the consensual bound property on $\bx$.
% %     \item
% % \end{itemize}
% % \subsection{Stationarity}
% % Define $\hat{x} = \frac{1}{n} \sum_{i=1}^n x_i$ and let $\bar{x} = \Pcal_{\Mcal}(\hat{x})$. We say $\bx = (x_1^\top, x_2^\top, \ldots, x_n^\top)^\top$ is an $\epsilon$-stationary point if
% % \[ \frac{1}{n}\sum_{i=1}^n \|x_i - \bar{x}\| \leq \epsilon, \;\; \forall \;\; i=1,2,\ldots,n \]
% % and
% % \[ \|\grad f(\bar{x})\| \leq \epsilon. \]


% \section{Decentralized Projected Gradient Tracking Method}
% In this section, we investigate a gradient tracking method for solving \eqref{prob:original}. The idea here is to adopt gradient tracking techniques in \cite{nedic2017achieving,qu2017harnessing,chen2021decentralized}  to get a better estimate for the full gradient. A crucial benefit of gradient tracking-type methods lies in the applicability of the use of constant step size.
% % the constant step size can be used and a faster convergence rate can be shown for the Euclidean algorithms. 
% In the $k$-th iteration, our decentralized projected Riemannian gradient tracking method (DPRSRM) performs the update 
%   \bee
%       \begin{aligned}
%         \bx_{k+1} & = \Pcal_{\Mcal^n}( \bW \bx_k  - \alpha  \Pcal_{T_{\bx_k}\Mcal^n}  (\bs_k)   ) \\
%         \bs_{k+1} &= \bW\bs_k + \grad f(\bx_{k+1}) - \grad f(\bx_k),
%       \end{aligned}
%   \eee
%   where $\alpha > 0$ is the step size and $\bs_{k}:=[s_{1,k}^\top, \cdots, s_{n,k}^\top]^\top$. The detailed description is presented in Algorithm \ref{alg:drgta}.

% \begin{algorithm}[htbp]
% \caption{Decentralized projected Riemannian gradient tracking method for solving \eqref{prob:original}} \label{alg:drgta}
% \begin{algorithmic}[1]
% \REQUIRE  Initial point $\bx_0\in \mathcal{M}^n$, an integer $t\geq \left\lceil \log_{\sigma_2\left(\frac{{\delta}}{8r\max\{1,L\}}\right)} \right \rceil$, set $k = 0$.
% \STATE Let $s_{i,0} = \grad f_i(x_{i,0})$ on each node $i\in [n]$.
% \WHILE {the condition is not met}
% \STATE Project onto tangent space: $d_{i,k} = \Pcal_{T_{x_{i,k}}\Mcal}( s_{i,k})$.
% \STATE Update $$x_{i,k+1} = \Pcal_{\Mcal}(\sum_{j=1}^n W_{ij}x_{i,k} - \alpha
% d_{i,k}), ~i\in [n].$$
% \STATE Riemannian gradient tracking:
% \be
% s_{i,k+1}  = \sum_{j=1}^n W_{ij}s_{j,k} + \grad f_i(x_{i,k+1}) - \grad f_i(x_{i,k}), ~i\in [n].
% \ee
% \STATE Set $k=k+1$.
% \ENDWHILE
% \end{algorithmic}
% \end{algorithm}

% In the following, we are going to show the convergence of DPRSRM. For the ease of notations, let us denote 
% $$
% \begin{aligned}
% \hat{g}_k &:=\frac{1}{n} \sum_{i=1}^n \grad f_i\left(x_{i, k}\right), \quad \hat{\mathbf{g}}_k:=\left(\mathbf{1}_n \otimes I_d\right) \hat{g}_k, \quad \bv_k:=[v_{1,k}^\top, \cdots, v_{n,k}^\top]^\top, \\
% {\bf G}_k & :=[\grad f_1(x_{1,k})^\top, \cdots, \grad f_n(x_{n,k})^\top ]^\top. 
% \end{aligned}
% $$
% Firstly, we have the following lemma on the relationship between ${\bf s}_k$ and $\hat{\bf G}_k$.
% \begin{lemma} \label{lem:diff-s-g}
% It holds that for any $k$,
% \be \|{\bf G}_{k+1} - {\bf G}_k\| \leq 4L\|\bx_k - \bar{\bx}_k\| + 2L\alpha \| \bv_k\|,
% \ee
% \be \label{eq:diff-s-g} \|{\bf s}_{k+1} - \hat{\bf G}_{k+1} \| \leq \sigma_2 \| {\bf s}_k - \hat{\bf G}_k \| + \| {\bf G}_{k+1} - {\bf G}_k \|. \ee

% \end{lemma}
% \begin{proof}
% The Lipschitz continuity of ${\bf G}$ yields
% \begin{equation}
% \begin{aligned}
%       \|{\bf G}_{k+1} - {\bf G}_k\| &  \leq L \|\bx_{k+1} - \bx_k\|.
%      % & \leq 2 \sigma_2 L \| \bx_k - \bar{\bx}_k \| + 2 \alpha L \|\bv_k\|
% \end{aligned}
% \end{equation}
% On the other hand, it holds that
% \begin{equation}
%     \begin{aligned}
%         \|  \bx_{k+1} - \bx_k \| & \leq \|  \bx_{k+1} - \bx_k + (I_{ nd}-\bW)\bx_k + \alpha \bv_k \| + \| (I_{ nd}-\bW)\bx_k + \alpha \bv_k\| \\
%         & \leq 2\| (I_{ nd}-\bW)\bx_k + \alpha \bv_k \| \leq 4\|\bx_k - \bar{\bx}_k\| + 2\alpha \| \bv_k\|,
%     \end{aligned}
% \end{equation}
% where the first inequality is from the triangle inequality and the second inequality is due to the definition of $\bx_{k+1}$. 
% It follows from the definition of $ \hat{\bf G}_{k+1}$ that
% \[ \begin{aligned}
% {\bf s}_{k+1} - \hat{\bf G}_{k+1} & = ((I_n - J) \otimes I_d) {\bf s}_{k+1} \\
% & = ((I_n -J) \otimes I_d) ( (W \otimes I_d) {\bf s}_k + {\bf G}_{k+1} - {\bf G}_{k})  \\
% % & = ((I_n - J) \otimes I_n) \left[ (W \otimes I_n) {\bf s}_{k} + {\bf G}_{k+1} - {\bf G}_{k} \right] \\
% & = ((W - J) \otimes I_d) {\bf s}_k + ((I_n - J) \otimes I_d)({\bf G}_{k+1} - {\bf G}_k).\\
% \end{aligned} \]
% Using the spectral property of $W$, we conclude that \eqref{eq:diff-s-g} holds.
% \end{proof}

% Before showing the boundedness of consensus error, we investigate the uniform boundedness of $\|\bs_k\|$ in the following lemma.
% \begin{lemma}\label{lem:stay-neighborhood}
% Under Assumption \ref{assum-f}, assume that the iterate $x_{i,k}$ is generated by Algorithm \ref{alg:drgta}, $t\geq \max\{\log_{\sigma_2}(\frac{1}{4\sqrt{n}}), \log_{\sigma_2}\frac{1}{4\sqrt{n}\zeta}, \log_{\sigma_2}\left(\frac{{\delta}}{16\sqrt{n}\zeta}\right)\}$, $\|\hat{x}_0 - \bar{x}_0\| \leq \frac{1}{2}{\delta}$. If  $\alpha<\min\{ \frac{{\delta}}{24L}, \frac{1}{8L}\}$, then for all $k$, 
% \begin{align}
% \|s_{i,k}\| & \leq 3L,~ \forall i\in [n], \label{eq:sk-bound} \\
%  \| \hat{x}_{k} - \bar{x}_{k} \| & \leq \frac{1}{2} {\delta}.
% \end{align}
% \end{lemma}

% \begin{proof}
% We prove it by induction on both $\|s_{i,k}\|$ and $ \| \hat{x}_{k} - \bar{x}_{k} \|$. Since $\|s_{i,0}\| = \| \grad f_i(x_{i,0}) \| \leq L $ for all $i\in [n]$ and $\|\hat{x}_0 - \bar{x}_0\| \leq \frac{1}{2}{\delta}$.  Suppose for some $k\geq 0$ that $\|s_{i,k}\| \leq 3L$ and $\|\hat{x}_k - \bar{x}_k\| \leq \frac{1}{2}{\delta}$.  Since $\|d_{i,k}\| \leq \|s_{i,k}\| \leq 3L$ and $\alpha< {\delta}/(24L)$, it follows from Lemma \ref{lem:stay-neighborhood} that
% \bee
%     \begin{aligned}
%          \sum_{j=1}^n {W_{ij}} x_{j,k} -\alpha  d_{i,k} \in U_{\Mcal} ({\delta}), ~ i = 1,\cdots,n, \quad
%         \| \hat{x}_{k+1} - \bar{x}_{k+1} \| \leq \frac{1}{2} {\delta}.
%     \end{aligned}
% \eee 
% Then, we have
% \be
% \begin{aligned}
% \| s_{i,k+1} - \hat{g}_k\| & = \| \sum_{j=1}^n W_{ij}s_{j,k} - \hat{g}_k + \grad f_i(x_{i,k+1}) - \grad f_i(x_{i,k}) \| \\
% & =  \| \sum_{j=1}^n (W_{ij} - \frac{1}{n})  s_{j,k} \|  + \|\grad f_i(x_{i,k+1}) - \grad f_i(x_{i,k}) \| \\
% & \leq \sigma_2 \sqrt{n} \max_i \| s_{i,k}\| + L \| x_{i,k+1} - x_{i,k} \| \\
% & \overset{\eqref{proximally-1}}{\leq} \sigma_2 \sqrt{n} \max_i\| s_{i,k}\| + 2L\|  \sum_{j=1}^n (W_{ij}  - \frac{1}{n}) x_{i,k}  - \alpha  d_{i,k}  \| \\
% & \leq (\sigma_2 \sqrt{n} + 2\alpha L) \max_i\| s_{i,k}\| + 2L \sigma_2 \sqrt{n} r \\
% & \leq \frac{1}{2}\max_i \| s_{i,k}\| + \frac{1}{2} L \\
% & \leq \frac{3}{2}L + \frac{1}{2} L \leq 2L.
% \end{aligned}
% \ee
% Hence,
% $
% \| s_{i,k+1}   \| \leq\| s_{i,k+1} - \hat{g}_k\| + \| \hat{g}_k\| \leq 3L$, where we use $\| \hat{g}_k\| \leq \frac{1}{n} \sum_{i=1}^n\| \grad f_i(x_{i,k}) \|\leq  L$. The proof is completed.
% \end{proof}

% \begin{lemma}\label{lemma:dpgta:consensus}
% Under the same conditions in Lemma \ref{lem:stay-neighborhood},  assume  the sequence $\{\bx_k\}$ is generated by Algorithm \ref{alg:drgta}.  The following holds
% \begin{equation}\label{eq:dgta:consensus}
% \begin{aligned}
%  \|\bx_{k+1} - \bar{\bx}_{k+1}\| & \leq 2\sigma_2 \|\bx_{k} - \bar{\bx}_{k}\| + 2\alpha\|\bv_k\|.
%  %& \leq (L\sigma_2)^{k+1} \|\bx_0 - \bar{\bx}_0\| + L\sqrt{n}B \sum_{l=0}^k(L\sigma_2)^{k-l}\alpha_l.
% \end{aligned}
% \end{equation}
%  Moreover, we have
% \be
% \| \bx_k - \bar{\bx}_k \|^2 \leq n 36 D L^2 \alpha^2.
% \ee
% \end{lemma}

% \begin{proof}
%     Note that
%     \[ \begin{aligned}
%     \|\bx_{k+1} - \bar{\bx}_{k+1}\| & \leq \| \bx_{k+1} - \bar{\bx}_k \| \\
%     & = \| \Pcal_{\Mcal^n}(\bW \bx_k - \alpha  \bv_k) - \Pcal_{\Mcal^n}( \hat{\bx}_k) \| \\
%     & { \leq }2 \|\bW \bx_k - \alpha  \bv_k -\hat{\bx}_k  \| \\
%     & \leq 2 \| (W \otimes I_d) \bx_k - \hat{\bx}_k \| + 2\alpha\|\bv_k\| \\
%     & {\leq} 2 \| ((W -J ) \otimes I_d) (\bx_k - \hat{\bx}_k) \| + 2\alpha\|\bv_k\| \\
%     % & \leq 2 \| \bW \bx_k - \hat{\bx}_k \| + \alpha  L_s \\
%     & \leq 2 \sigma_2 \|\bx_k -\hat{\bx}_k\| + 2\alpha\|\bv_k\|   \\
%     & \leq 2\sigma_2 \|\bx_k -\bar{\bx}_k\| + 2\alpha\|\bv_k\|,
% \end{aligned}
% \]
% where the second inequality is from 2-Lipschitz continuity of $P_{\Mcal}(\cdot)$ and the fourth inequality is due to \eqref{eq:sk-bound}.
% Since $\|\bv_k\| \leq \|\bs_k\| \leq 3\sqrt{n}L$, using the same 
% argument of Lemma \ref{lemma:dpg:consensus}, there exists $D = \mathcal{O}(\frac{1}{(1-2\sigma_2)^2})$, which is independent of $L$, such that
% \be
% \frac{1}{n} \| \bx_{k} - \bar{\bx}_k \|^2 \leq 36 D L^2 \alpha^2, \quad k\geq 0.
% \ee
% The proof is completed.
% %     \be
% % \begin{aligned}
% % \|\bx_{k+1} - \bx_k\| & \leq \|    \Pcal_{\Mcal}( \bx_k  - \alpha    \bv_k  - (I- W)\bx_k ) - \bx_k \|  \\
% % & \leq 2  \| \alpha  \bv_k  + (I- W)\bx_k \| \\
% % & \leq \alpha L \|\bs_k\| + L(1- \lambda_n(W))\|\bx_k - \bar{\bx}_k \|.
% % \end{aligned}
% % \ee
% \end{proof}

% \begin{lemma}\label{lemma:dgta:majorized}
% Under Assumption \ref{assum-f}, assume that the iterate $x_{i,k}$ is generated by Algorithm \ref{alg:drgta}. The following inequality holds:
% \begin{equation}
%      f(\bar{x}_{k+1})  \leq  f(\bar{x}_k) - (\alpha  - \frac{L^2}{2} \alpha^2) \|\hat{g}_k\|^2 + \mathcal{C}_1  \alpha^2 \frac{1}{n}\|\bs_k \|^2 +\mathcal{C}_2 \frac{1}{n} \| \bx_k - \bar{\bx}_k\|^2 +\mathcal{C}_3 \frac{1}{n} \| \bx_{k+1} - \bar{\bx}_{k+1}\|^2,
% \end{equation}
% where
% \begin{equation}
%     \begin{aligned}
%           \mathcal{C}_1 & = (2Q+ 3)L + 108Q^2\alpha^2L^3 + 3L, \\
%          \mathcal{C}_2 & = 4LQ+4L + 36M_2^2DL +  108DL^3(8Q + \sqrt{n}L_2 + M_2)^2\alpha^2, \\
%          \mathcal{C}_3 & = 36M_2^2D L + 108DL^3M_2^2\alpha^2.
%     \end{aligned}
% \end{equation}
% \end{lemma}

% \begin{proof}
% It follows Lemma \ref{lemma:lipsctz} and $L_g \leq L$ that 
% \be
% \| \hat{g}_k - \grad f(\bar{x}_k) \|^2 \leq \frac{1}{n} \sum_{i=1}^n \|\grad f_i(x_k) -  \grad f_i(\bar{x}_k) \|^2 \leq \frac{L^2}{n} \|\bx_k - \bar{\bx}_k \|^2.
% \ee
% and
% \begin{equation}\label{Lip-ineq}
%     \begin{aligned}
%     f(\bar{x}_{k+1})  \leq & f(\bar{x}_k) + \left<\grad f(\bar{x}_k),\bar{x}_{k+1}-\bar{x}_k  \right>+ \frac{L}{2}\|\bar{x}_{k+1}-\bar{x}_k\|^2\\
%      = & f(\bar{x}_k) + \left<\hat{g}_k ,\bar{x}_{k+1}-\bar{x}_k  \right>+\left<\grad f(\bar{x}_k) -\hat{g}_k,\bar{x}_{k+1}-\bar{x}_k  \right> + \frac{L}{2}\|\bar{x}_{k+1}-\bar{x}_k\|^2 \\
%       \leq & f(\bar{x}_k) + \left<\hat{g}_k ,\bar{x}_{k+1}-\bar{x}_k  \right>+ \frac{3L}{4}\|\bar{x}_{k+1}-\bar{x}_k\|^2  + \frac{1}{L} \| \grad f(\bar{x}_k) -\hat{g}_k \|^2 \\
%         \leq & f(\bar{x}_k) + \left<\hat{g}_k ,\hat{x}_{k+1}-\hat{x}_k  \right>  + \left<\hat{g}_k,  \bar{x}_{k+1}- \hat{x}_{k+1}+\hat{x}_k - \bar{x}_k  \right> \\
%         &+ \frac{L}{n}  \| \bx_k - \bar{\bx}_k \|^2 + \frac{3L}{4}\|\bar{x}_{k+1}-\bar{x}_k\|^2 \\
%     % & \leq   \mathcal{L}_{\alpha}(\bx_k) + \left<\grad f(\bx_k) - \bs_k,\bx_{k+1}-\bx_k  \right>- ( \frac{1}{2\alpha} -\frac{L}{2})\|\bx_{k+1}-\bx_k\|^2\\
%     \end{aligned}
% \end{equation}
% By Young's inequality, we have
% \begin{equation} \label{eq:grad-prod2}
%     \left<\hat{g}_k,  \bar{x}_{k+1}- \hat{x}_{k+1}+\hat{x}_k - \bar{x}_k  \right> \leq \frac{\alpha^2L}{2}\| \hat{g}_k  \|^2 + \frac{1}{\alpha^2L}(\| \bar{x}_{k+1}- \hat{x}_{k+1} \|^2 + \|\hat{x}_k - \bar{x}_k\|^2)
% \end{equation}
% Combining \eqref{Lip-ineq} and \eqref{eq:grad-prod2} leads to
% \begin{equation}
%     \begin{aligned}
%      f(\bar{x}_{k+1})  \leq & f(\bar{x}_k) + \underbrace{\left<\hat{g}_k ,\hat{x}_{k+1}-\hat{x}_k  \right>}_{b_1}  +  \frac{\alpha^2L}{2}\| \hat{g}_k   \|^2 + \underbrace{ \frac{1}{\alpha^2L}(\| \bar{x}_{k+1}- \hat{x}_{k+1} \|^2 + \|\hat{x}_k - \bar{x}_k\|^2)}_{b_2}\\
%      & + \frac{L}{n}  \| \bx_k - \bar{\bx}_k \|^2 + \underbrace{ \frac{3L}{4}\|\bar{x}_{k+1}-\bar{x}_k\|^2}_{b_3}. \\
%     \end{aligned}
% \end{equation}
% Now, let us bound $b_1$, $b_2$, and $b_3$, respectively. For $b_1$, it holds that
% \begin{equation} \label{eq:b1-0}
%     \begin{aligned}
%     b_1  =& \left< \hat{g}_k , \frac{1}{n} \sum_{i=1}^n( x_{i,k+1} - x_{i,k} + \alpha s_{i,k} + \nabla \phi_i(\bx_k)   ) \right> -  \left< \hat{g}_k, \frac{1}{n} \sum_{i=1}^n(  \alpha s_{i,k} + \nabla \phi_i(\bx_k)   ) \right> \\
%     \leq &  \left< \hat{g}_k   , \frac{1}{n} \sum_{i=1}^n( x_{i,k+1} - x_{i,k} + \alpha d_{i,k} + \nabla \phi_i(\bx_k)  ) \right> +\left< \hat{g}_k, \frac{1}{n} \sum_{i=1}^n \alpha(s_{i,k} - d_{i,k}) \right> - \alpha  \|\hat{g}_k\|^2
%     \end{aligned}
% \end{equation}
% Since $s_{i,k} - d_{i,k} \in N_{x_{i,k}}\Mcal$, it follows that
% \be \label{eq:b1-1}
% \begin{aligned}
%      \left< \hat{g}_k, \frac{1}{n} \sum_{i=1}^n \alpha(s_{i,k} - d_{i,k}) \right>
%    = &\frac{\alpha}{n} \sum_{i=1}^n\left< \hat{g}_k - \grad f_i(x_{i,k}),  s_{i,k} - d_{i,k} \right> \\
%    \leq & \frac{1}{4nL} \sum_{i=1}^n \|\hat{g}_k - \grad f_i(x_{i,k})  \|^2 + \frac{\alpha^2L}{n}\sum_{i=1}^n \| \Pcal_{N_{x_{i,k}}\Mcal}(s_{i,k}) \|^2\\
%    \leq & \frac{1}{4n^2L} \sum_{i=1}^n \sum_{j=1}^n\|\grad f_j(x_{j,k}) - \grad f_i(x_{i,k})  \|^2 + \frac{\alpha^2L}{n}\| \bs_k \|^2\\
%    \leq & \frac{L}{n} \| \bx_k - \bar{\bx}_k \|^2 + \frac{\alpha^2L}{n}\| \bs_k \|^2.
% \end{aligned}
% \ee
% Let  $\nabla \phi_i(\bx_k) + \alpha d_{i,k} = d_{i,1} + d_{i,2}, $
% where $d_{i,1} = \Pcal_{T_{x_{i,k}}\Mcal}(\nabla \phi_i(\bx_k) + \alpha d_{i,k}) , d_{i,2} = \nabla \phi_i(\bx_k)+ \alpha d_{i,k} - d_{i,1}$.  Therefore,
% \be \label{eq:b1-2}
% \begin{aligned}
% &\left< \hat{g}_k   , \frac{1}{n} \sum_{i=1}^n( x_{i,k+1} - x_{i,k} + \alpha d_{i,k} + \nabla \phi_i(\bx_k)  ) \right> \\
% = & \frac{1}{n} \sum_{i= 1}^n \left< \hat{g}_k,  \Pcal_{\Mcal}(x_{i,k} - d_{i,1} - d_{i,2}  ) - [x_{i,k}  - d_{i,1}    ]\right>+\frac{1}{n}\sum_{i= 1}^n \left< \hat{g}_k,  d_{i,2}\right> \\
% \leq & \frac{LQ }{n} \sum_{i=1}^n \|  d_i\|^2 +\frac{1}{n} \sum_{i= 1}^n \left< \hat{g}_k - \grad f_i(x_{i,k}),  d_{i,2}\right>\\
% \leq & \frac{LQ }{n} \sum_{i=1}^n \|  d_i\|^2 + \frac{1}{4nL}\sum_{i= 1}^n
%  \|\hat{g}_k - \grad f_i(x_{i,k})\|^2 + \frac{L}{n}\sum_{i= 1}^n \|d_{i,2}\|^2 \\
%  \leq & \frac{LQ+L }{n} \sum_{i=1}^n \|  d_i\|^2 + \frac{L}{n} \|\bar{\bx}_k - \bx_k\|^2 \\
%   \leq & \frac{LQ+L }{n} \|(I_{nd}-\bW)\bx_k + \alpha \bv_k\|^2 + \frac{L^2}{n}  \|\bar{\bx}_k - \bx_k\|^2 \\
%    \leq & \frac{4LQ+3L  }{n}\|\bar{\bx}_k - \bx_k\|^2  + \frac{2LQ+2L }{n}\alpha^2 \|\bs_k\|^2.
% \end{aligned}
% \ee
% % It follows Lemma \ref{lemma-project} that
% % \be
% % c_1  \leq  Q D  \sum_{i = 1}^n \| d\|^2  \leq Q D  \sum_{i = 1}^n \| \nabla \phi_i(\bx_k) + \alpha s_{i,k} \|^2. %\leq  \frac{1}{2}Q M  \sum_{i = 1}^n \| \nabla \phi_i(\bx_k) \|^2 + \alpha^2 \|s_{i,k}\|^2 \leq (1-\lambda_n(W))^2 \|\bx_k - \bar{\bx}_k \|^2 + \alpha^2 \| \bs_k\|^2.
% % \ee
% % For $c_2$, we have
% % \be
% % \begin{aligned}
% % c_2 & \leq \| \hat{g}_k\| M \sum_{i=1}^n  \| d_1 \|^2  \leq  D M \sum_{i=1}^n\| (\nabla \phi_i(\bx_k) + \alpha s_{i,k} \|^2).
% % \end{aligned}
% % \ee
% % Secondly,
% % \be
% % \begin{aligned}
% % c_3  & =  \sum_{i=1}^n \left< \hat{g}_k - \grad f_i(x_{i,k}), u_2  \right> \\
% % & \leq \frac{1}{2} \sum_{i = 1}^n \| \hat{g}_k  - \grad f_i (x_{i,k})\|^2 + \|\nabla \phi_i(\bx_k) + \alpha s_{i,k} \|^2 \\
% % & \leq L_G^2 \| \bx_k - \bar{\bx}_k \|^2 + \sum_{i = 1}^n\|\nabla \phi_i(\bx_k) + \alpha s_{i,k} \|^2.
% % \end{aligned}
% % \ee
% % \be
% % \left< \hat{g}_k   , \frac{1}{n} \sum_{i=1}^n( x_{i,k+1} - x_{i,k} + \alpha( s_{i,k} + \alpha^{-1}\nabla \phi_i(\bx_k) )  ) \right> \leq \frac{\alpha}{2} \|\hat{g}_k  \|^2 + \frac{1}{2\alpha  n}\sum_{i=1}^n \|x_{i,k+1} - x_{i,k} + \alpha( s_{i,k} + \alpha^{-1}\nabla \phi_i(\bx_k) )   \|^2.
% % \ee
% % Therefore
% % \begin{equation}
% %     \begin{aligned}
% %     a_1
% %     \leq & \frac{1}{2\alpha  }\sum_{i=1}^n \|x_{i,k+1} - x_{i,k} + \alpha( s_{i,k} + \alpha^{-1}\nabla \phi_i(\bx_k) )   \|^2 - \frac{\alpha}{2} n \|\hat{g}_k\|^2 \\
% %      \leq &   \frac{\alpha L^2}{2} \sum_{i=1}^n\| ( s_{i,k} + \alpha^{-1}\nabla \phi_i(\bx_k) )  ) \|^2 - \frac{\alpha}{2} n \|\hat{g}_k\|^2 \\
% %      \leq &  \sqrt{2} \alpha L^2  (\|\bs_k\|^2+ \alpha^{-2}\|(I-W)\bx_k\|^2) - \frac{\alpha}{2} n \|\hat{g}_k\|^2 \\
% %      \leq &  \sqrt{2} \alpha L^2   \|\bs_k\|^2+ \sqrt{2} \alpha^{-1} L^2 (1-\lambda_n(W))\|\bx_k - \bar{\bx}_k\|^2 ) - \frac{\alpha}{2} n \|\hat{g}_k\|^2 \\
% %     \end{aligned}
% % \end{equation}
% Plugging \eqref{eq:b1-1} and \eqref{eq:b1-2} into \eqref{eq:b1-0} gives
% \be \label{eq:b1}
% \begin{aligned}
% b_1 &\leq  \frac{4LQ+ 4 L}{n} \|\bx_k - \bar{\bx}_k \|^2 +\frac{2LQ+3L}{n} \alpha^2 \| \bs_k\|^2 -  \alpha  \|\hat{g}_k\|^2,
% \end{aligned}
% \ee
% For $b_2$, applying Lemma \ref{lemma:quadratic} yields
% % \begin{equation}
% %     \begin{aligned}
% %     \|\hat{\bx}_k - \bar{\bx}_k\|^2 & \leq 2 \|\hat{\bx}_k - \bx_k  \|^2 + 2\|\bx_k - \bar{\bx}_k\|^2 \leq 4\| \bx_k - \bar{\bx}_k\|^2
% %     \end{aligned}
% % \end{equation}
% % Then
% \begin{equation} \label{eq:b2}
% \begin{aligned}
%       b_2 & \leq \frac{M_2^2}{\alpha^2 L n^2}(\| \bx_k - \bar{\bx}_k\|^4 + \| \bx_{k+1} - \bar{\bx}_{k+1}\|^4  ) \\
%     & \leq \frac{36M_2^2D L }{ n}(\| \bx_k - \bar{\bx}_k\|^2 + \| \bx_{k+1} - \bar{\bx}_{k+1}\|^2  ).
% \end{aligned}
% \end{equation}
% By \eqref{eq:sum-consen-grad}, we obtain
% \begin{equation} \label{eq:b3}
%     \begin{aligned}
%     b_3  \leq &\frac{3L}{4} \left( \frac{4(8Q + \sqrt{n}L_2 + M_2)^2}{n^2} \|\bx_k - \bar{\bx}_k\|^4 +\frac{16Q^2\alpha^4}{n^2} \| \bv_k \|^4  + 4\alpha^2 \| \hat{d}_{k}  \|^2  + \frac{4M_2^2}{n^2} \| \bx_{k+1}  - \bar{\bx}_{k+1} \|^4\right)\\
%      \leq &\frac{3L}{4} \left(\frac{4(8Q + \sqrt{n}L_2 + M_2)^2}{n^2} \|\bx_k - \bar{\bx}_k\|^4 +\frac{16Q^2\alpha^4}{n^2} \| \bv_k \|^4  + \frac{4\alpha^2}{n} \| \bv_k  \|^2  + \frac{4M_2^2}{n^2} \| \bx_{k+1}  - \bar{\bx}_{k+1} \|^4\right)\\
%       \leq & \frac{108DL^3(8Q + \sqrt{n}L_2 + M_2)^2}{n}\alpha^2 \|\bx_k - \bar{\bx}_k\|^2   + \frac{108DL^3M_2^2}{n}\alpha^2 \| \bx_{k+1}  - \bar{\bx}_{k+1} \|^2 \\
%      & + \frac{108Q^2L^3\alpha^4 + 3L\alpha^2}{n}\| \bv_k \|^2.
%   %& \leq \left( \frac{36}{n} \|\bx_k - \bar{\bx}_k\|^2 + \frac{16\alpha^2}{n}\| \grad f(\bx_k) \|^2  \right) + 8M_2^2CL^2  ( \alpha^2 \| \bx_k  - \bar{\bx}_k \|^2 +\alpha^2 \| \bx_{k+1}  - \bar{\bx}_{k+1} \|^2 )
%     \end{aligned}
% \end{equation}
% % Finally,
% % \begin{equation}
% %     \begin{aligned}
% %     a_4 & \leq  \alpha^2L_G^2(\| \hat{\bf G}_k\|^2 + \|\alpha^{-1}(I-W)\bx_k  \|^2) \\
% %     &\leq  \alpha^2L_G^2 n \| \hat{g}^k\|^2 +  \alpha^2L_G^2 \alpha^{-2}(1-\lambda_n(W))^2 \|  \bx_k - \bar{\bx}_k\|^2
% %     \end{aligned}
% % \end{equation}
% Then, combining \eqref{eq:b1}, \eqref{eq:b2}, and \eqref{eq:b3} gives
% \begin{equation}
% \begin{aligned}
%    &   f(\bar{x}_{k+1})  \leq  f(\bar{x}_k) + b_1 + \frac{\alpha^2 L^2}{2} \|\hat{g}_k\|^2 + b_2 + \frac{L}{n}\|\bx_k - \bar{\bx}_k\|^2 + b_3 \\
%       \leq & f(\bar{\bx}_k) - ( \alpha - \frac{L^2}{2} \alpha^2 ) \|\hat{g}_k \|^2 + \frac{(2Q+ 3)L + 108Q^2\alpha^2L^3 + 3L}{n} \alpha^2 \|  \bs_k\|^2\\
%       & + \frac{4LQ+4L + 36M_2^2DL +  108DL^3(8Q + \sqrt{n}L_2 + M_2)^2\alpha^2}{n} \|\bx_k - \bar{\bx}_k\|^2  \\
%       & + \frac{36M_2^2D L + 108DL^3M_2^2\alpha^2}{ n}  \| \bx_{k+1}  - \bar{\bx}_{k+1} \|^2.
% \end{aligned}
% \end{equation}
% The proof is completed.
% \end{proof}

% \begin{lemma}[\cite{xu2015augmented}]\label{lemma:sequence:sum}
%     Let $\{u_k\}_{k\geq 0}$ and $\{w_k\}_{k\geq 0}$ be two positive scalar sequences such that for all $k\geq 0$
% $$
% u_{k+1} \leq \eta u_{k} + w_k,
% $$
% where $\eta\in (0,1)$ is the decaying parameter. Let $\Gamma(k)=\sum_{i=0}^ku_i^2$ and $\Omega(k) = \sum_{i=0}^k w_i^2$. Then we have
% $$
% \Gamma(k) \leq c_0 \Omega(k) + c_1,
% $$
% where $c_0 = \frac{2}{(1-\eta)^2}$ and $c_1 = \frac{2}{1-\eta^2}u_0^2$.

% \end{lemma}


% \begin{theorem}
%   Under Assumption \ref{assum-f}, let $\bx_0$ satisfy that $\| \hat{x}_0 - \bar{x}_0 \| \leq \frac{1}{2}{\delta}$,
%   \be
%   t\geq \max\left\{\log_{\sigma_2}(\frac{1}{4\sqrt{n}}), \log_{\sigma_2}\frac{1}{4\sqrt{n}\zeta}  ,\log_{\sigma_2}\left(\frac{{\delta}}{16\sqrt{n}\zeta}\right) \right\},
%   \ee
%   and $\alpha$ satisfy that $\alpha<\min\{ \frac{{\delta}}{24L}, \frac{1}{8L},\frac{1}{L^2}\}$. Then it follows that for the
% sequences generated by Algorithm \ref{alg:drgta}:
% \be
% \min_{k=0,\cdots,K} \frac{1}{n}\|\bs_k\|^2 \leq \frac{8(f(\bar{x}_0) - f^* + \tilde{C}_5)}{\alpha K},
% \ee
% \be
% \min_{k=0,\cdots,K} \frac{1}{n} \|\bx_k - \bar{\bx}_k \|^2  \leq \frac{32\tilde{C}_0\alpha(f(\bar{x}_0) - f^* + \tilde{C}_5) + \tilde{C}_1}{ K},
% \ee
% \be
% \min_{k=0,\cdots,K} \|\grad f(\bar{x}_k) \|^2  \leq \frac{(16 +64\tilde{C}_0\alpha^2 )(f(\bar{x}_0) - f^* + \tilde{C}_5)+ 2L^2 \tilde{C}_1\alpha}{\alpha K}.
% \ee
% %where
% %\be
% %\begin{aligned}
% %d
% %\end{aligned}
% %\ee
% \end{theorem}


% \begin{proof}
% Applying Lemma \ref{lemma:sequence:sum} to \eqref{eq:dgta:consensus} yields 
% \be\label{eq:sum:consensus}
% \frac{1}{n}\sum_{k=0}^{K+1}\| \bx_k - \bar{\bx}_k \|^2 \leq \tilde{C}_0 \frac{4\alpha^2}{n}\sum_{k=0}^{K+1}\| \bs_k \|^2 + \tilde{C}_1,
% \ee 
% where $\tilde{C}_0 = \frac{2}{(1-2\sigma_2)^2}$ and $\tilde{C}_1 = \frac{2}{1-4\sigma_2^{2t}}\frac{1}{n}\|\bx_0 - \bar{\bx}_0\|^2$. It follows from Lemma \ref{lemma:dgta:majorized} that
% \begin{equation} \label{eq:desc}
% \begin{aligned}
%    &  f(\bar{x}_{k+1}) \\
%       \leq & f(\bar{x}_0) - (\alpha n - \frac{\alpha^2L^2n}{2})\sum_{k=0}^{K} \|\hat{g}_k\|^2 + \mathcal{C}_1 \alpha^2 \frac{1}{n} \sum_{k=0}^{K} \|\bs_k \|^2 +(\mathcal{C}_2 + \mathcal{C}_3) \frac{1}{n}\sum_{k=0}^{K+1} \| \bx_k - \bar{\bx}_k\|^2\\
%       \leq & f(\bar{x}_0) - \frac{\alpha    }{2} \sum_{k=0}^{K} \|\hat{g}_k\|^2 + (\mathcal{C}_1 + 4\tilde{C}_0(\mathcal{C}_2 + \mathcal{C}_2)) \alpha^2 \frac{1}{n} \sum_{k=0}^{K+1} \|\bs_k \|^2 +(\mathcal{C}_2 + \mathcal{C}_3)\tilde{C}_1 \\
%        \leq & f(\bar{x}_0) - \frac{\alpha    }{2} \sum_{k=0}^{K} \|\hat{g}_k\|^2 + (\mathcal{C}_1 + 4\tilde{C}_0(\mathcal{C}_2 + \mathcal{C}_2)) \alpha^2 \frac{1}{n} \sum_{k=0}^{K} \|\bs_k \|^2 \\
%        & + 3(\mathcal{C}_1 + 4\tilde{C}_0(\mathcal{C}_2 + \mathcal{C}_2)) \alpha^2L +(\mathcal{C}_2 + \mathcal{C}_3)\tilde{C}_1 \\
%         \leq & f(\bar{x}_0) - \frac{\alpha    }{2} \sum_{k=0}^{K} \|\hat{g}_k\|^2 + (\mathcal{C}_1 + 4\tilde{C}_0(\mathcal{C}_2 + \mathcal{C}_2)) \alpha^2 \frac{1}{n} \sum_{k=0}^{K} \|\bs_k \|^2 + \tilde{C}_4.
%      \end{aligned}
% \end{equation}
% Next, we need to associate $\|\hat{g}_k\|^2$ with $\|\bs_k\|^2$. Since  $\| \bs_k  \|  \leq \|\bs_k - \hat{\bf G}_k \| + \|\hat{\bf G}_k \|$, we have 
% \be\label{eq:tmp10}
% -\sum_{k=0}^{K} \|\hat{g}_k \|^2  = -\frac{1}{n}\sum_{k=0}^{K} \|\hat{\bf G}_k \|^2  \leq \frac{1}{n} \sum_{k=0}^{K}\|\bs_k - \hat{\bf G}_k \|^2 - \frac{1}{2n}\sum_{k=0}^{K}\|\bs_k\|^2.
% \ee
% Applying Lemma \ref{lemma:sequence:sum} into \eqref{eq:diff-s-g} yields
% \bee
% \begin{aligned}
% \frac{1}{n}\sum_{k=0}^{K}\|\bs_k - \hat{\bf G}_k \|^2 &\leq \tilde{C}_2\frac{1}{n}\sum_{k=0}^{K}\|{\bf G}^{k+1} - {\bf G}^k \|^2 + \tilde{C}_3 \\
% & \leq \tilde{C}_2\frac{1}{n}\sum_{k=0}^{K} (32L^2\|\bx_k - \bar{\bx}_k \|^2 + 8L^2\alpha^2 \|\bs_k\|^2  ) + \tilde{C}_3 \\
% & \leq \tilde{C}_2(128\tilde{C}_0  + 8)L^2\alpha^2\frac{1}{n}\sum_{k=0}^{K}  \|\bs_k\|^2   + 32\tilde{C}_1\tilde{C}_2L^2  + \tilde{C}_3.
% \end{aligned}
% \eee
% Plugging the above inequality into \eqref{eq:tmp10} leads to
% \be
% -\sum_{k=0}^{K} \|\hat{g}_k \|^2    \leq  \left[ \tilde{C}_2(128\tilde{C}_0  + 8)L^2\alpha^2 - \frac{1}{2}  \right]  \frac{1}{n}\sum_{k=0}^{K}\|\bs_k\|^2  + 32\tilde{C}_1\tilde{C}_2L^2  + \tilde{C}_3.
% \ee
% Then, by \eqref{eq:desc}, it holds
% \bee
% \begin{aligned}
%    &  f(\bar{x}_{k+1}) \\
%      % \leq & f(\bar{\bx}_0) - (\alpha n - \frac{\alpha^2L^2n}{2})\sum_{k=0}^{K} \|\hat{g}_k\|^2 + \mathcal{C}_1 \alpha^2 \sum_{k=0}^{K} \|\bs_k \|^2 +(\mathcal{C}_2 + \mathcal{C}_3) \sum_{k=0}^{K+1} \| \bx_k - \bar{\bx}_k\|^2\\
%       \leq & f(\bar{x}_0) - \frac{\alpha    }{2}   \left[\frac{1}{2} -  \tilde{C}_2(128\tilde{C}_0  + 8)L^2\alpha^2 - 2(\mathcal{C}_1 + 4\tilde{C}_0(\mathcal{C}_2 + \mathcal{C}_2)) \alpha  \right] \frac{1}{n}  \sum_{k=0}^{K}\|\bs_k\|^2 \\
%       &  +(\tilde{C}_4 + 16\tilde{C}_1\tilde{C}_2L^2 \alpha + \frac{\alpha}{2}\tilde{C}_3 ) \\
%        \leq & f(\bar{x}_0) - \frac{\alpha    }{8}\frac{1}{n}     \sum_{k=0}^{K}\|\bs_k\|^2   +(\tilde{C}_4 + 16\tilde{C}_1\tilde{C}_2L^2 \alpha + \frac{\alpha}{2}\tilde{C}_3 ). 
%      \end{aligned}
% \eee
% This implies
% \bee
% \frac{\alpha}{8} \sum_{k=0}^{K} \|\hat{g}_k \|^2  \leq \frac{\alpha}{8} \frac{1}{n} \sum_{k=0}^{K} \|\bs_k\|^2 \leq f(\bar{x}_0) - f^* + \tilde{C}_5,
% \eee
% where $\tilde{C}_5 = \tilde{C}_4 + 16\tilde{C}_1\tilde{C}_2L^2 \alpha + \frac{\alpha}{2}\tilde{C}_3$ and $f^* = \min_{x\in \Mcal} f(x)$. Furthermore, 
% \bee
% \min_{k=0,\cdots,K} \| \hat{g}_k\|^2 = \min_{k=0,\cdots,K} \| \hat{s}_k\|^2 \leq  \min_{k=0,\cdots,K} \frac{1}{n}\|\bs_k\|^2 \leq \frac{8(f(\bar{x}_0) - f^* + \tilde{C}_5)}{\alpha K}.
% \eee
% Then, it follows from \eqref{eq:sum:consensus} that
% \bee
% \min_{k=0,\cdots,K} \frac{1}{n} \|\bx_k - \bar{\bx}_k \|^2  \leq \frac{32\tilde{C}_0\alpha(f(\bar{\bx}_0) - f^* + \tilde{C}_5) + \tilde{C}_1}{ K}.
% \eee
% Since
% \bee
% \|\grad f(\bar{x}_k) \|^2  \leq 2\| \hat{g}_k\|^2 + 2\|\grad f(\bar{x}_k)  - \hat{g}_k \|^2  \leq 2\| \hat{g}_k\|^2 + \frac{2L^2}{n}\| \bx_k - \bar{\bx}_k\|^2. 
% \eee
% We conclude that 
% \bee
% \min_{k=0,\cdots,K} \|\grad f(\bar{x}_k) \|^2  \leq \frac{(16 +64\tilde{C}_0\alpha^2 )(f(\bar{x}_0) - f^* + \tilde{C}_5)+ 2L^2 \tilde{C}_1\alpha}{\alpha K}.
% \eee
% \end{proof}
%  % \section{Convergence}




%%%%%%%%%%%%%%%%%%%%%%%%%%%%%%%
\section{Numerical experiments}\label{sec:exm}
In this section, we compare our proposed DPRSRM with DRSGD in \cite{chen2021decentralized} and DRPGD in \cite{deng2023decentralized} on decentralized principal component analysis. It is important to note that the original DRPGD is a deterministic algorithm that utilizes the local full gradient. To adapt it to the online setting, we replace the full gradient with a stochastic gradient for the local updates. The numerical results on decentralized low-rank matrix completion are provided in the  supplementary 
 material.

\subsection{Decentralized principal component analysis} \label{sub:pca}
The decentralized principal component analysis (PCA) problem can be expressed mathematically as follows:
\be \label{prob:pca}
\min _{\mathbf{x} \in \mathcal{M}^{n}}-\frac{1}{2 n} \sum_{i=1}^{n} \text{tr}(x_{i}^{\top} A_{i}^{\top} A_{i} x_{i}), \;\; \text { s.t. } \;\; x_{1}=\ldots=x_{n},
\ee
where \(\mathcal{M}\) is the Stiefel manifold \(\text{St}(d,r)\), \(A_{i} \in \mathbb{R}^{m_{i} \times d}\) is the data matrix corresponding to the \(i\)-th node, and \(m_{i}\) denotes the number of samples. It is worth noting that if \(x^*\) is a solution to this problem, then any transformation of \(x^*\) by an orthogonal matrix \(Q \in \mathbb{R}^{r \times r}\) is also a valid solution. The distance between two points \(x\) and \(x^*\) is then calculated as:

\[
d_s(x, x^*) := \min_{Q \in \mathbb{R}^{r \times r}, \; Q^\top Q = QQ^\top = I_r} \; \|xQ - x^*\|.
\]

\subsubsection{Synthetic dataset}
In our study, we set the parameters as follows: \(m_{1} = \ldots = m_{n} = 1000\), \(d = 10\), and \(r = 5\). A matrix \(B \in \mathbb{R}^{1000n \times d}\) is generated, and its singular value decomposition (SVD) is performed, yielding \(B = U \Sigma V^\top\), where \(U \in \mathbb{R}^{1000n \times d}\) and \(V \in \mathbb{R}^{d \times d}\) are orthogonal matrices, and \(\Sigma \in \mathbb{R}^{d \times d}\) is a diagonal matrix. To control the distribution of singular values, we define \(\tilde{\Sigma} = \text{diag}(\gamma^j)\) with \(\gamma\) chosen from the interval \((0,1)\). The matrix \(A\) is then formed as \(A = U \tilde{\Sigma} V^\top \in \mathbb{R}^{1000n \times d}\). The matrices \(A_i\) are derived by partitioning the rows of \(A\) into \(n\) equally sized subsets. It can be shown that the first \(r\) columns of \(V\) represent the solution to \eqref{prob:pca}. In our experiments, the parameters \(\gamma\) and \(n\) are set to 0.8 and 8, respectively.



We employ fixed step sizes for all algorithms. The step size is set to $\alpha = \frac{\hat{\beta}}{\sqrt{K}}$ with $K$ being the maximal number of iterations.  The grid search is utilized to find the best $\hat{\beta}$ for each algorithm. The momentum parameter is chosen as $\tau = 0.999$. The batch size in each node is set as $10$ and the maximum iteration is set $K = 2000$. The clipping constant is set as $B=10^{8}$.    We choose the polar decomposition as the retraction operator for DRSGD. We test several graph matrices to model the topology across the nodes, namely, the Erdos-Renyi (ER) network with probability $p = 0.3, 0.6$, and the Ring network. Throughout this section, we select the mixing matrix $W$ to be the Metropolis constant edge weight matrix \cite{shi2015extra}.


\begin{figure}[htp]
	\centering
	\includegraphics[width = 0.46 \textwidth]{fig/eig/DRgraph_PCAorg_consensus_8000_10_5_08_nodes_8_p_03_Ring_metropolis_hastings.pdf}
	\includegraphics[width = 0.46 \textwidth]{fig/eig/DRgraph_PCAorg_obj_8000_10_5_08_nodes_8_p_03_Ring_metropolis_hastings.pdf} \\
	\includegraphics[width = 0.46 \textwidth]{fig/eig/DRgraph_PCAorg_grad_8000_10_5_08_nodes_8_p_03_Ring_metropolis_hastings.pdf}
	\includegraphics[width = 0.46 \textwidth]{fig/eig/DRgraph_PCAorg_dist_to_grad_8000_10_5_08_nodes_8_p_03_Ring_metropolis_hastings.pdf}
	\caption{Numerical results on the synthetic dataset with different network graphs.}	
	\label{fig:num-pca-graph}
\end{figure}


\begin{figure}[htp]
	\centering
    \includegraphics[width = 0.46 \textwidth]{fig/eig/DRALL_PCAorg_consensus_8000_10_5_08_nodes_8_p_06_ER_metropolis_hastings.pdf}
    \includegraphics[width = 0.46 \textwidth]{fig/eig/DRALL_PCAorg_obj_8000_10_5_08_nodes_8_p_06_ER_metropolis_hastings.pdf} \\
    \includegraphics[width = 0.46 \textwidth]{fig/eig/DRALL_PCAorg_grad_8000_10_5_08_nodes_8_p_06_ER_metropolis_hastings.pdf}
    \includegraphics[width = 0.46 \textwidth]{fig/eig/DRALL_PCAorg_dist_to_opt_8000_10_5_08_nodes_8_p_06_ER_metropolis_hastings.pdf}
% 	\includegraphics[width = 0.32 \textwidth]{fig/DRALL_consensus_error_8000_10_5_08_nodes_8_p_03_ER_metropolis_hastings.eps}
% 	\includegraphics[width = 0.32 \textwidth]{fig/DRALL_grad_norm_8000_10_5_08_nodes_8_p_03_ER_metropolis_hastings.eps}
% % 	\includegraphics[width = 0.24 \textwidth]{fig/DRALL_obj_value_8000_10_5_08_nodes_8_p_03_ER_metropolis_hastings.eps}
% 	\includegraphics[width = 0.32 \textwidth]{fig/DRALL_dist_to_opt_8000_10_5_08_nodes_8_p_03_ER_metropolis_hastings.eps}
	\caption{Results on the synthetic dataset with ER $p=0.6$.}	
	\label{fig:num-pca-alg}
\end{figure}


Firstly, we test all algorithms with different network graphs, namely, Ring, ER $p=0.3$, and ER $p=0.6$. The results are shown in Figure \ref{fig:num-pca-graph}. We see that there is not much difference among different graphs except for the consensus error. This is consistent with the existing results for DRSGD \cite{chen2021decentralized} and DPRGD \cite{deng2023decentralized}. 
Secondly, Figure \ref{fig:num-pca-alg} presents a comparison among the three algorithms. Our DPRSRM outperforms the other two, with DRSGD and DPRGD showing comparable performance.

% \begin{figure}[htp]
% 	\centering
% 	\includegraphics[width = 0.46 \textwidth]{fig/eig/DRconsensus_PCAorg_consensus_8000_10_5_08_nodes_8_p_06_ER_metropolis_hastings.pdf}
% 	\includegraphics[width = 0.46 \textwidth]{fig/eig/DRconsensus_PCAorg_obj_8000_10_5_08_nodes_8_p_06_ER_metropolis_hastings.pdf} \\
% 	\includegraphics[width = 0.46 \textwidth]{fig/eig/DRconsensus_PCAorg_grad_8000_10_5_08_nodes_8_p_06_ER_metropolis_hastings.pdf}
% 	\includegraphics[width = 0.46 \textwidth]{fig/eig/DRconsensus_PCAorg_dist_to_opt_8000_10_5_08_nodes_8_p_06_ER_metropolis_hastings.pdf}
% 	\caption{Numerical results on synthetic dataset with different numbers of consensus steps on graph ER $p=0.6$.}	
% 	\label{fig:num-pca-consensus}
% \end{figure}





\begin{figure}[htp]
	\centering
	\includegraphics[width = 0.46 \textwidth]{fig/eig/DRALL_sto_PCA_real_consensus_60000_784_5_08_nodes_8_p_06_ER_metropolis_hastings.pdf}
	\includegraphics[width = 0.46 \textwidth]{fig/eig/DRALL_sto_PCA_real_obj_60000_784_5_08_nodes_8_p_06_ER_metropolis_hastings.pdf}\\
    \includegraphics[width = 0.46 \textwidth]{fig/eig/DRALL_sto_PCA_real_grad_60000_784_5_08_nodes_8_p_06_ER_metropolis_hastings.pdf}
	\includegraphics[width = 0.46 \textwidth]{fig/eig/DRALL_sto_PCA_real_dist_to_opt_60000_784_5_08_nodes_8_p_06_ER_metropolis_hastings.pdf}\\
	\caption{Results on Mnist dataset with ER $p=0.3$.}	
	\label{fig:num-pca-real}
\end{figure}
% \begin{figure*}
% 	\centering
% 	\includegraphics[width = 0.32 \textwidth]{fig/DRALL_consensus_error_8000_10_5_08_nodes_8_p_03_Ring_metropolis_hastings.eps}
% 	\includegraphics[width = 0.32 \textwidth]{fig/DRALL_grad_norm_8000_10_5_08_nodes_8_p_03_Ring_metropolis_hastings.eps}
% % 	\includegraphics[width = 0.32 \textwidth]{fig/DRALL_obj_value_8000_10_5_08_nodes_8_p_05_ER_metropolis_hastings.eps}
% 	\includegraphics[width = 0.32 \textwidth]{fig/DRALL_dist_to_opt_8000_10_5_08_nodes_8_p_03_Ring_metropolis_hastings.eps}
% 	\caption{Numerical results on synthetic dataset. Ring}	
% 	\label{fig:num-pca-Ring}
% \end{figure*}




% \begin{figure*}
% 	\centering
% 	\includegraphics[width = 0.32 \textwidth]{fig/DRALL_consensus_error_8000_10_5_08_nodes_8_p_08_ER_metropolis_hastings.eps}
% 	\includegraphics[width = 0.32 \textwidth]{fig/DRALL_grad_norm_8000_10_5_08_nodes_8_p_08_ER_metropolis_hastings.eps}
% % 	\includegraphics[width = 0.32 \textwidth]{fig/DRALL_obj_value_8000_10_5_08_nodes_8_p_08_ER_metropolis_hastings.eps}
% 	\includegraphics[width = 0.32 \textwidth]{fig/DRALL_dist_to_opt_8000_10_5_08_nodes_8_p_08_ER_metropolis_hastings.eps}
% 	\caption{Numerical results on synthetic dataset. ER, prob = 0.8}	
% 	\label{fig:num-pca-prob08}
% \end{figure*}



\subsubsection{Mnist dataset}
% To evaluate the efficiency of the proposed DPRSRM, 
We also conduct numerical tests on the Mnist dataset \cite{lecun1998mnist}. The training images consist of 60000 handwritten images of size $32 \times 32$ and are used to generate $A_i$'s. We first normalize the data matrix by dividing 255 and randomly split the data into $n=8$ nodes with equal cardinality. Then, each node holds a local matrix $A_i$ of dimension $\frac{60000}{n} \times 784$. We compute the first 5 principal components, i.e., $d=784, r=5$.



For all algorithms, we use the fixed step sizes $\alpha = \frac{\hat{\beta}}{60000}$ with a best-chosen $\hat{\beta}$, batch size 1500 and momemtum parameter $\tau = 0.999$. 
% The fixed step size of the form $\alpha = \frac{\hat{\alpha}}{10000}$ is taken for all algorithms.
Similar to the synthetic setting, our DPRSRM algorithm demonstrates superior performance compared to other algorithms in terms of objective function values, gradient norms, and distances to the optimal solution. It is important to note that due to the use of stochastic gradients, consensus among the algorithms may be affected by noise, which can be reduced by using a smaller step size.

% Similar to the above setting, we see from Figure \ref{fig:num-pca-real} that DPRSRM converges much faster than DRSGD and DRPGD.


% \begin{figure*}
% 	\centering
% 	\includegraphics[width = 0.32 \textwidth]{fig/DRALL_mnist_consensus_error_10000_784_1_08_nodes_8_p_03_Ring_metropolis_hastings.eps}
% 	\includegraphics[width = 0.32 \textwidth]{fig/DRALL_mnist_grad_norm_10000_784_1_08_nodes_8_p_03_Ring_metropolis_hastings.eps}
% % 	\includegraphics[width = 0.32 \textwidth]{fig/DRALL_mnist_obj_value_10000_784_1_08_nodes_8_p_03_Ring_metropolis_hastings.eps}
% 	\includegraphics[width = 0.32 \textwidth]{fig/DRALL_mnist_dist_to_opt_10000_784_1_08_nodes_8_p_03_Ring_metropolis_hastings.eps}
% 	\caption{Numerical results on Mnist dataset. Ring}	
% 	\label{fig:num-pca-real}
% \end{figure*}

% \begin{figure}
% 	\centering
% 	\includegraphics[width = 0.46 \textwidth]{fig/DRALL_dist_to_opt_60000_784_1_08_nodes_8_p_03_Ring_metropolis_hastings.pdf}
% 	\includegraphics[width = 0.46 \textwidth]{fig/DRALL_real_consensus_error_60000_784_1_08_nodes_8_p_03_Ring_metropolis_hastings.pdf}
% 	\caption{Numerical results on MNIST dataset. Left: distance to optima, right: consensus error.}
% 	\label{fig:num-pca-real}	
% \end{figure}




\subsection{Low-rank matrix completion problem}\label{sec:more}
The low-rank matrix completion (LRMC) problem aims to reconstruct a matrix \(A \in \mathbb{R}^{d \times T}\) with low rank from its partially observed entries. Let \(\Omega\) represent the set of indices corresponding to the observed entries in \(A\). The LRMC problem of rank \(r\) can be expressed as:
\be \label{prob:lrmc} \min_{X \in {\rm Gr}(d,r), V\in \R^{r \times T}} \quad \frac{1}{2} \| \Pcal_{\Omega}(XV- A) \|^2, \ee
where \(\text{Gr}(d,r)\) represents the Grassmann manifold of \(r\)-dimensional subspaces in \(\mathbb{R}^d\), and \(\mathcal{P}_{\Omega}\) is a projection operator that selects the elements of \(A\) indexed by \(\Omega\), setting the rest to zero.

In a decentralized context, assume the matrix \(\mathcal{P}_{\Omega}(A)\) is divided into \(n\) equal parts by columns, labeled as \(A_1, A_2, \ldots, A_n\), each part corresponding to a different node. Replacing the Grassmann manifold constraint with the Stiefel manifold, the decentralized LRMC problem \cite{deng2023decentralized} is therefore formulated as:
\be \label{prob:dlrmc}
\begin{aligned}
\min \quad & \frac{1}{2} \sum_{i=1}^n \| \Pcal_{\Omega_i}(X_i V_i(X) - A_i) \|^2,  \\
\st \quad & X_1 = X_2 = \cdots = X_n, \\
& X_i \in {\rm St}(d,r), \quad \forall i \in [n],
\end{aligned}
\ee
where $\Omega_i$ is the corresponding indices set of $\Omega$ and $V_i(X):= {\rm argmin}_{V} \| \Pcal_{\Omega_i} (XV - A_i)\|$.

% The goal of the low-rank matrix completion problem (LRMC) is to recover a low-rank matrix $A \in \R^{d \times T}$ from its partial observations. Let $\Omega$ be the set of indices of known entries in $A$, the rank-$r$ LRMC problem can be written as
% \be \label{prob:lrmc} \min_{X \in {\rm Gr}(d,r), V\in \R^{r \times T}} \quad \frac{1}{2} \| \Pcal_{\Omega}(XV- A) \|^2, \ee
% where ${\rm Gr}(d,r):=\{ {\rm all~} r {\rm ~dimensional~subspaces~in~}\R^d \}$ is the Grassmann manifold and the projection operator $\Pcal_{\Omega}$ is defined in an element-wise manner with $(\Pcal_{\Omega}(A))_{ij} = A_{ij}$ if $(i,j) \in \Omega$ and $0$ otherwise.

% For the decentralized setting, let us consider the case when the data matrix $\Pcal_{\Omega}(A)$ is equally divided into $n$ nodes by columns, denoted by $A_1, A_2, \ldots, A_n$. As the Grassmann manifold ${\rm Gr}(n,d)$ is a quotient manifold with its total space being ${\rm St}(n,d)$, we can also regard \eqref{prob:lrmc} as an optimization on the Stiefel manifold, which is similar to the case of the principal component analysis in Section \ref{sec:exm}. We refer to \cite{absil2009optimization,hu2022riemannian} for the equivalence of the optimization algorithms between the quotient manifold and its total space.
% Then, the decentralized LRMC problem is
% \be \label{prob:dlrmc}
% \begin{aligned}
% \min \quad & \frac{1}{2} \sum_{i=1}^n \| \Pcal_{\Omega_i}(X_i V_i(X) - A_i) \|^2,  \\
% \st \quad & X_1 = X_2 = \cdots = X_n, \\
% & X_i \in {\rm St}(d,r), \quad \forall i \in [n],
% \end{aligned}
% \ee
% where $\Omega_i$ is the corresponding indices set of $\Omega$ and $V_i(X):= {\rm argmin}_{V} \| \Pcal_{\Omega_i} (XV - A_i)\|$.

For numerical tests, we consider random generated $A$. To be specific, we first generate two random matrices $L \in \R^{d \times r}$ and $R \in \R^{r \times T}$, where each element obeys the standard Gaussian distribution. For the indices set $\Omega$, we generate a random matrix $B$ with each element following from the uniform distribution, then set $\Omega_{ij} = 1$ if $B_{ij} \leq \nu$ and $0$ otherwise. The parameter $\nu$ is set to $r(d+T-r)/(dT)$. In the implementations, we set $T=1000, d=50, r=10$, and $\alpha = \frac{\hat{\beta}}{\sqrt{K}}$ for all algorithms with $K$ being the maximal number of iterations. $\hat{\beta}$ is tuned to get the best performance for each algorithm individually. The Ring graph is used. The results are reported in Figure \ref{fig:num-mc}, where DPRGD is omitted due to its similar performance with DRSGD. We see that DPRSRM outperforms DRSGD.
\begin{figure}[htp]
	\centering
	\includegraphics[width = 0.32\textwidth]{fig/mc/DRALL_rand16_consensus_500_50_10_08_nodes_8_p_06_Ring_metropolis_hastings.pdf}
    \includegraphics[width = 0.32\textwidth]{fig/mc/DRALL_rand16_obj_500_50_10_08_nodes_8_p_06_Ring_metropolis_hastings.pdf}
     \includegraphics[width = 0.32\textwidth]{fig/mc/DRALL_rand16_grad_500_50_10_08_nodes_8_p_06_Ring_metropolis_hastings.pdf}
	\caption{Numerical results for the decentralized LRMC problem with the Ring graph.}
	\label{fig:num-mc}	
\end{figure}













\section{Conclusions and Limitations}
This work develops a decentralized projection Riemannian stochastic recursive momentum method by assuming that each
node has access to a stochastic first-order oracle.  Our algorithm
leverages local hybrid variance reduction and gradient tracking
to achieve a lower oracle complexity compared with the existing online
methods. It requires only $\mathcal{O}(1)$ gradient evaluations per iteration for each local node and does not require restarting with a large batch gradient.   %Our experiment findings confirm the superiority
%of our proposed algorithm. 

% \textbf{Limitations.} Our paper motivates several important questions for future research. The lack of closed-form solutions for the projection operator on certain manifolds suggests a need to explore methods for approximate projection calculation. Additionally, our step-size selection depends on the proximal smoothness constant $R$, emphasizing the importance of accurately estimating $R$ for specific manifolds.


\section*{Acknowledgements} 

This effort was supported by the SciAI Center, and funded by the Office of Naval Research (ONR), under Grant Number N00014-23-
1-2729 (J.H.), as well as by the National Natural Science Foundation of China
(NSFC) grants 12401419 (K.D.).







\bibliographystyle{siamplain}
\bibliography{ref}

\newpage
\appendix
\section{Technical Lemmas}\label{appen}




\begin{lemma}\label{lem:inequa}
    Given any vectors $a,b\in \mathbb{R}^n$, it holds that
    \begin{equation}
        \langle a, \frac{b}{\|b\|} \rangle \geq \|a\| - 2\|a - b\|.
    \end{equation}
\end{lemma}

\begin{proof}
  It follows from the Cauchy inequality that
\begin{equation}
\begin{aligned}
    \langle a, \frac{b}{\|b\|} \rangle & =   \langle a - b, \frac{b}{\|b\|} \rangle + \|b\| \geq -\|a - b\| + \|b\| \geq  - 2\| a-b \| + \|a\|, \\
     % \sum_{i=1}^n   \langle a, \frac{b_i}{\|b_i\|} \rangle & =   \langle a - b, \frac{b}{\|b\|} \rangle + \|b\| \geq -\|a - b\| + \|b\| \geq  - 2\| a-b \| + \|a\|, 
    \end{aligned}
\end{equation}
where the second inequality use $\|a\| \leq \|a - b\| + \|b\|$.
\end{proof}


\begin{lemma}[\cite{xu2015augmented}, Lemma 2]\label{lem:curr}

Let $u_k$ and $w_k$ be two positive scalar sequences such that for all $k\geq 1$
\begin{equation}\label{2}
  u_{k} \leq \eta u_{k-1} + w_{k-1},
\end{equation}
where $\eta\in (0,1)$ is the decaying factor.  Then we have 
\begin{equation}\label{22}
       \sum_{k=0}^K u_k \leq \frac{u_0}{1-\eta} +  \frac{1}{1-\eta}  \sum_{k=0}^{K-1} w_{k}. 
\end{equation}
\end{lemma}


% \begin{proof}
%     Applying the recurrence relation  to \eqref{2} for $k>1$ yields 
%     \begin{equation}
%         u_k \leq \eta^{k-1} u_1 + \sum_{i=1}^{k-1}\eta^{i-1} w_{k-i}.
%     \end{equation}
%     Summing the above relations over $k$ from $1$ to $K$ leads to
%     \begin{equation}
%         \begin{aligned}
%             \sum_{k=1}^K u_k  & \leq \sum_{k=1}^K \eta^{k-1} u_1 + \sum_{k=1}^K \sum_{i=1}^{k-1}\eta^{i-1} w_{k-i} \\
%             &  = \frac{u_1}{1-\eta} + \sum_{i=1}^{K-1} \sum_{k=1}^{K-i}\eta^{k-1} w_{i} \\
%              & \leq   \frac{u_1}{1-\eta} + \left(\sum_{k=1}^{K-1}\eta^k \right) \left( \sum_{k=1}^{K-1} w_{k} \right) \\
%              & \leq   \frac{u_1}{1-\eta} +  \frac{1}{1-\eta}  \sum_{k=1}^{K-1} w_{k}.   \end{aligned}
%     \end{equation}
%     The proof is completed. 
% \end{proof}

The following inequality is the control of the distance between the Euclidean mean $\hat{x}$ and the manifold mean $\bar{x}$ by the square of consensus error.
\begin{lemma}[\cite{deng2023decentralized}]\label{lemma:quadratic}
   For any $\bx\in \Mcal^n$ satisfying $\|x_i - \bar{x}\| \leq {\delta}$, {$i\in [n]$,}  we have 
   \be\label{eq:distance-an-rm}
   \|\bar{x} - \hat{x} \| \leq M_2 \frac{\|\bx - \bar{\bx}\|^2}{n},
   \ee
   where $M_2 =\max_{x\in { \bar{U}_\Mcal({\delta})}} \|D^2 \Pcal_{\Mcal}(x) \|_{\rm op}$.
\end{lemma}


The following Lipschitz-type inequality for the projection operator $\Pcal_{\Mcal}(\cdot)$ is crucial in the analysis of projection-based methods. 

\begin{lemma}[\cite{deng2023decentralized}]\label{lemma-project}
For any $x\in\Mcal, u\in \{u\in \mathbb{R}^{d\times r}:\|u\|\leq {\delta}\}$, 
there exists a constant $Q$ such that
\be\label{projec-second-order1}
\| \Pcal_{\Mcal}(x + u)  - x - \Pcal_{T_x\Mcal}(u) \| \leq Q\|u\|^2.
\ee

\end{lemma}


% In the following lemma, we show that $\left\| \sum_{i=1}^n \grad \phi_i(\bx) \right \|$ is bounded by the square of consensus error. {It can be inferred by combining \cite{chen2021decentralized} and the proximal smoothness of the submanifold.}
% \begin{lemma}[\cite{deng2023decentralized}] \label{lem:grad-consensus}
% Let $L_2 := 2 \max_{x \in {\rm conv}(\Mcal)} \| D^2 \Pcal_{T_x\Mcal}(\cdot) \|_{\rm op}$. For any $\bx \in \Mcal^n$,  it holds that
% \be\label{eq:sum-consen-grad} \| \sum_{i=1}^n \grad \phi_i(\bx)  \| \leq \sqrt{n} L_2 \|\bx - \bar{\bx}\|^2.  \ee
% %where $L_2 := 2 \max_{x \in {\rm conv}(\Mcal)} \| D^2 \Pcal_{T_x\Mcal}(\cdot) \|$. 
% \end{lemma}


% \begin{proof}
%     Using the fact that $\sum_{i=1}^n \nabla \phi_i(\bx) = \sum_{i=1}^n \left[  x_i - \sum_{j=1}^n W_{ij} x_j \right] = 0$, we have
%     \[ \begin{aligned}
%        & \left\| \sum_{i=1}^n \grad \phi_i(\bx) \right \|  = \left \|  \sum_{i=1}^n \left[ \Pcal_{T_{x_i}\Mcal}(\nabla \phi_i(\bx)) - \Pcal_{T_{\bar{x}}\Mcal}(\nabla \phi_i(\bx)) \right]   \right \| \\
%         & \leq \frac{L_2}{2} \sum_{i=1}^n \left[ \| x_i - \bar{x}\| \|\nabla \phi_i (\bx)\| \right] 
%          \leq \frac{L_2}{2}  \left( \max_{i\in [n]} \|x_i - \bar{x} \|\right) \sum_{i=1}^n \| x_i - \sum_{j=1} W_{ij}x_j \| \\
%         & \leq \frac{L_2}{2}  \left( \max_{i\in [n]} \|x_i - \bar{x} \|\right) \left[ \sum_{i=1}^n \| x_i -\bar{x}\| + \| \sum_{j=1} W_{ij}x_j -\bar{x} \| \right]\\
%         & \leq L_2 \left( \max_{i\in [n]} \|x_i - \bar{x} \|\right) \sum_{i=1}^n \| x_i - \bar{x} \| 
%          \leq \sqrt{n} L_2  \|\bx - \bar{\bx}\|^2,
%     \end{aligned}\]
%     where the first inequality is due to the Lipschitz continuity of $\Pcal_{T_{x}\Mcal}(\cdot)$ over $x \in \Mcal$, the third inequality is due to the triangle inequality of $\|\cdot \|$, the fourth inequality uses the convexity of $\|\cdot\|$, and the last inequality comes from $\|a\|_1 \leq \sqrt{n}\|a\|$ and $\|a\|_{\infty} \leq \|a\|$ for any $a \in \R^{n}$. 
% \end{proof}

% The next technical result bounds the distance between $\bar{x}_{k+1}$ and $\bar{x}_{k}$ for a given iterative process.
% \begin{lemma}[\cite{deng2023decentralized}] \label{lem:diff-avg-x}
% Let $x_{i,k+1} = \Pcal_{\Mcal}\left(\sum_{j=1}^n W_{ij}x_{i,k} - \alpha  v_{i,k} \right)$, where $v_{i,k} \in T_{x_{i,k}}\Mcal$. Denote $\hat{v}_k: = \frac{1}{n}\sum_{i=1}^n v_{i,k}$. Suppose that Assumption \ref{assum-w} holds. 
% It holds that
% \be \label{eq:diff-avg-x}
% \begin{aligned}
%  \| \bar{x}_{k+1} - \bar{x}_k \| \leq & \frac{8Q + \sqrt{n}L_2 + M_2}{n} \|\bx_k - \bar{\bx}_k\|^2 +\frac{2Q\alpha^2}{n} \| \bv_k \|^2  \\
%  &+ \alpha \| \hat{u}_{k}  \|  + \frac{M_2}{n} \| \bx_{k+1}  - \bar{\bx}_{k+1} \|^2. \end{aligned} \ee
% % Moreover, it holds that
% % \be \label{eq:diff-x}
% % \begin{aligned}
% %  \| \bx_{k+1} - \bx_k \| \leq & \frac{8Q + \sqrt{n}(L_2+1) + M_2}{\sqrt{n}} \|\bx_k - \bar{\bx}_k\|^2 +\frac{2Q\alpha^2}{\sqrt{n}} \| \bu_k \|^2  \\
% %  &+ \alpha \sqrt{n} \| \hat{u}_{k}  \|  + (\frac{M_2}{\sqrt{n}}+1) \| \bx_{k+1}  - \bar{\bx}_{k+1} \|^2. \end{aligned} \ee
% \end{lemma}
%\section*{Acknowledgements} The authors are grateful to Prof. Samuel Burer and three  anonymous referees for their valuable comments and suggestions.



% The next lemma is a technical observation and follows from Lemma 3 in \cite{cutkosky2019momentum}. 
% \begin{lemma}\label{lemma:expect}
%  Let us denote $\Delta_k = \bd_k - \grad f(\bx_k)$. Then we have
% $$
% \begin{aligned}
% & \mathbb{E}\left[\left(\grad f\left(\bx_k, \mathbf{\xi}_k\right)-\grad f\left(\bx_k\right)\right) \cdot \alpha^{-1}\left(1-\tau\right)^2 \Delta_{k-1}\right]=0 \\
% & \mathbb{E}\left[\left(\grad f\left(\bx_k, \mathbf{\xi}_k\right)-\grad f\left(\bx_{k-1}, \mathbf{\xi}_k\right)-\grad f\left(\bx_k\right)+\grad f\left(\bx_{k-1}\right)\right) \cdot \alpha^{-1}\left(1-\tau\right)^2 \Delta_{k-1}\right]=0.
% \end{aligned}
% $$
% \end{lemma}




\section{Proof of Section \ref{sec4}}



% We first investigate the uniform boundedness of $\|\bs_k\|$ in the following lemma.
% \begin{lemma}\label{lemma:neibohood}
% Let $\{\bx_k\}_k$ be the sequence generated by Algorithm \ref{alg:drgta}. Suppose that Assumptions \ref{assum-w} and \ref{assum-f} hold. If $\bx_0 \in \mathcal{N}(\delta)$, $\alpha \leq \min\left\{ \delta,  \frac{(R(1-\sigma_2) - 2\delta)\delta}{R\sqrt{n}} \right\}$, and $\delta \leq  \min\{ \frac{R}{2}, \frac{R(1-\sigma_2)}{2}\}$, it follows that for all $k$, $\bx_k \in \mathcal{N}(\delta)$  and
% \be\label{eq:sk-bound}
% \|s_{i,k}\| \leq 4 L,~ \forall i\in [n].
% \ee
% \end{lemma}

% \begin{proof}






% We prove it by induction on both $\|s_{i,k}\|$ and $ \max_i\| x_{i,k} - \bar{x}_{k} \|$. Since $\|s_{i,0}\| = \| \grad f_i(x_{i,0}) \| \leq L $ for all $i\in [n]$ and $\max_i\|x_{i,0} - \bar{x}_0\| \leq \frac{1}{2}\revise{\delta}$.  Suppose for some $k\geq 0$ that $\|s_{i,k}\| \leq 4 L$ and $\max_i\|x_{i,k} - \bar{x}_k\| \leq \frac{1}{2}\revise{\delta}$.  Since $\|v_{i,k}\| \leq \|s_{i,k}\| \leq 4 L$ and $\alpha< \revise{\delta}/(96L)$, it follows from Lemma \ref{lem:stay-neighborhood} that
% \bee
%     \begin{aligned}
%          \sum_{j=1}^n {W_{ij}^t} x_{j,k} -\alpha  v_{i,k} \in \bar{U}_{\Mcal} (\revise{\delta}), ~ i = 1,\cdots,n, \quad \max_i\| x_{i,k+1} - \bar{x}_{k+1} \| \leq \frac{1}{2} \revise{\delta}.
%     \end{aligned}
% \eee 
% Then, we have
% \bee
% \begin{aligned}
% \| s_{i,k+1} - \hat{g}_k\| & = \| \sum_{j=1}^n W^t_{ij}s_{j,k} - \hat{d}_k + d_{i,k+1} - d_{i,k} \| \\
% & =  \| \sum_{j=1}^n (W^t_{ij} - \frac{1}{n})  s_{j,k} \|  + \| d_{i,k+1} - d_{i,k} \| \\
% & \leq \sigma_2^t \sqrt{n} \max_i \| s_{i,k}\| + 2L  \\
% & \leq  \sigma_2^t \sqrt{n} \max_i\| s_{i,k}\| + 2 L \\
% & \leq \frac{1}{4} \max_i\| s_{i,k}\| + 2 L \leq L + 2 L \leq 3L. 
% \end{aligned}
% \eee
% Hence, $
% \| s_{i,k+1}   \| \leq\| s_{i,k+1} - \hat{g}_k\| + \| \hat{g}_k\| \leq 3L + L \leq 4 L$, where we use $\| \hat{g}_k\| \leq \frac{1}{n} \sum_{i=1}^n\| \grad f_i(x_{i,k}) \|\leq  L$. The proof is completed.
% \end{proof}

\subsection{Proof of Section \ref{sec:linear-consensus}}

\begin{revisionblock}
\begin{proof}[Proof of Lemma \ref{lemma:neibohood}]

The tracking identity and double stochasticity give
$\hat{\bs}_k=\hat{\bd}_k$ for every $k$. By the bounded initialization,
$\|\bs_0\|=\|\bd_0\|\leq\sqrt nB$. Suppose that
$\|\bs_k\|\leq3\sqrt nB/(1-\sigma_2)$.
Then, we have
\bee
\begin{aligned}
\| \bs_{k+1} - \hat{\bd}_k\| & = \| \bW \bs_k - \hat{\bd}_k + \bd_{k+1} - \bd_k \| \\
& =   \| \bW \bs_k - \hat{\bs}_k + \bd_{k+1} - \bd_k \|\\
& =  \| [(W-J)\otimes I_d ]  \bs_k  + \bd_{k+1} - \bd_k \|\\
& \leq  \sigma_2 \| \bs_k \| + 2 \sqrt{n} B.
\end{aligned}
\eee
Since $\|\hat{\bd}_k\|\leq\sqrt nB$, it follows that
\[
 \|\bs_{k+1}\|\leq\frac{3\sqrt nB\sigma_2}{1-\sigma_2}
 +3\sqrt nB=\frac{3\sqrt nB}{1-\sigma_2}.
\]
Induction proves \eqref{eq:sk-bound}.
\end{proof}


\begin{proof}[Proof of Lemma \ref{lem:stay-neighborhood}]
  % We first prove that $\|\bv_k\|$ is bounded. We note that by the initialization strategy $s_{i,0} = d_{i,0}$, it holds that 
  % \begin{equation}
  %     \hat{s}_k = \hat{d}_k. 
  % \end{equation}
  % Then we have that
  % \begin{equation}
  %     \|\bs_k\|^2 = \sum
  % \end{equation}
    We prove the two claims by induction. Assume that
    $\bx_k\in\mathcal{N}(\delta)$. For any $i\in[n]$,
\[ \begin{aligned}
\| \sum_{j=1}^n W_{ij} x_{j,k} - \alpha v_{i,k} - \bar{x}_k  \| 
 \leq &   \| \sum_{j=1}^n W_{ij} (x_{j,k} - \bar{x}_k ) - \alpha  v_{i,k}   \| \\
 \leq &   \sum_{j=1}^n W_{ij} \|x_{j,k} - \bar{x}_k  \| + \alpha \|v_{i,k} \| \\
 \leq   &  \delta +   \sqrt{nD_B}\alpha \leq 2\delta.
\end{aligned}\]
By $\delta<R/2$, we have
$\sum_{j=1}^nW_{ij}x_{j,k}-\alpha v_{i,k}\in\bar U_{\Mcal}(2\delta)$.
Moreover, $\max_i\|x_{i,k}-\bar x_k\|\leq\delta$, so
$\hat x_k\in\bar U_{\Mcal}(\delta)\subset\bar U_{\Mcal}(2\delta)$.
The $\kappa_M$-Lipschitz property of $\Pcal_{\Mcal}$ therefore gives
    \[ \begin{aligned}
      &  \| \bx_{k+1} - \bar{\bx}_{k+1} \|  \leq \| \bx_{k+1} - \bar{\bx}_k \| \\
        &= \|\Pcal_{\Mcal^n}( \bW \bx_k - \alpha \bv_k) - \Pcal_{\Mcal^n}(\hat{\bx}_k) \| \\
        & \leq \frac{R}{R- 2\delta} \| \bW \bx_k - \alpha \bv_k - \hat{\bx}_k\| \\
        & \leq \frac{R}{R- 2\delta} \| [ (W - J) \otimes I_d] ( \bx_k -\hat{\bx}_k) -\alpha \bv_k  \| \\
        & \leq \frac{R \sigma_2}{R- 2\delta} \|\bx_k - \bar{\bx}_k\| + \frac{R \alpha}{R- 2\delta} \|\bv_k\| \\
        & \leq \frac{R \sigma_2}{R- 2\delta} \delta + \frac{R \alpha}{R- 2\delta} \sqrt{nD_B}.
        %& \leq \frac{\delta}{2} + \frac{\delta}{2} = \delta,
    \end{aligned} \]
 For a fixed $\delta$, the inequality
 $\|\bx_{k+1}-\bar{\bx}_{k+1}\|\leq\delta$ holds provided that
 \begin{equation}
     \alpha \leq \frac{(R(1-\sigma_2) - 2\delta)\delta}{R\sqrt{nD_B}}.
 \end{equation}
 The right-hand side is positive whenever
 $0<\delta<R(1-\sigma_2)/2$, as required by \eqref{cond-alpha-delta}.
 This closes the induction.
\end{proof}
\end{revisionblock}

\subsection{Proof of Section \ref{sec:proof:dprgt}}

\begin{revisionblock}
\begin{proof}[Proof of Lemma~\ref{lem:coupled-one-step}]
We prove the four inequalities in turn. For the transition from $k$ to
$k+1$, all conditional expectations are taken with respect to
$\mathcal F_{k+1}$, i.e., immediately before drawing
$\boldsymbol\xi_{k+1}$. Thus, $\bx_k,\bx_{k+1}$ and $\bd_k$ are measurable
with respect to the conditioning sigma-algebra.

\emph{Step 1: estimator error.}
Set
\[
 \Delta_k:=\bd_k-\grad f(\bx_k),\qquad
 \mathbf g_k^\xi:=\grad f(\bx_k,\boldsymbol\xi_{k+1}),\qquad
 \mathbf g_k:=\grad f(\bx_k).
\]
Since $B\geq G_{\max}$, every block of $\mathbf g_{k+1}$ belongs to the clipping
ball. Blockwise nonexpansiveness of the Euclidean projection therefore gives
\begin{equation}\label{eq:clip-estimator-comparison}
 \|\bd_{k+1}-\mathbf g_{k+1}\|^2
 \leq\|\bq_{k+1}-\mathbf g_{k+1}\|^2.
\end{equation}
The un-clipped error has the exact decomposition
\begin{equation}\label{eq:estimator-centered-decomp}
 \bq_{k+1}-\mathbf g_{k+1}=(1-\tau)\Delta_k
 +\boldsymbol\varepsilon_{k+1},
\end{equation}
where
\begin{align*}
 \boldsymbol\varepsilon_{k+1}:={}&
 \tau(\mathbf g_{k+1}^\xi-\mathbf g_{k+1})\\
 &+(1-\tau)\{\mathbf g_{k+1}^\xi-\mathbf g_k^\xi
                  -(\mathbf g_{k+1}-\mathbf g_k)\}.
\end{align*}
Both terms in braces are evaluated using the same fresh sample. Their sum is
conditionally centered, so
$\mathbb E[\boldsymbol\varepsilon_{k+1}\mid\mathcal F_{k+1}]=0$ and the cross term with
$\Delta_k$ in \eqref{eq:estimator-centered-decomp} vanishes. We do not assume
that the two fresh-sample terms are mutually independent. Instead,
$\|a+b\|^2\leq2\|a\|^2+2\|b\|^2$, conditional variance is bounded by the
corresponding second moment, and Assumption~\ref{assum:stochatis} give
\begin{align}
 \mathbb E[\|\boldsymbol\varepsilon_{k+1}\|^2\mid\mathcal F_{k+1}]
 \leq{}&2n\tau^2\nu^2
 +2(1-\tau)^2\bar L^2\|\bx_{k+1}-\bx_k\|^2.       \label{eq:centered-noise-bound}
\end{align}
Taking full expectation in \eqref{eq:clip-estimator-comparison} and using
\eqref{eq:estimator-centered-decomp}--\eqref{eq:centered-noise-bound} yields
\begin{equation}\label{eq:repaired-estimator-step}
 \mathsf{E}_{k+1}\leq(1-\tau)^2\mathsf{E}_k+2\tau^2\nu^2
 +\frac{2\bar L^2}{n}\mathbb{E}\|\bx_{k+1}-\bx_k\|^2.
\end{equation}
This is the point at which the original proof introduced the spurious factor
$\alpha^2$: mean-square Lipschitz continuity gives the last squared
iterate-difference directly.

For completeness, let $\by_k:=\bW\bx_k-\alpha\bv_k$. Since
$\bx_{k+1}=\Pcal_{\Mcal^n}(\by_k)$ and $\bx_k\in\Mcal^n$, optimality of the
projection gives $\|\bx_{k+1}-\by_k\|\leq\|\bx_k-\by_k\|$. Hence,
\begin{align*}
 \|\bx_{k+1}-\bx_k\|
 &\leq2\|\by_k-\bx_k\|\\
 &\leq2\|(I_{nd}-\bW)(\bx_k-\bar\bx_k)+\alpha\bv_k\|.
\end{align*}
Because $\|I_{nd}-\bW\|\leq2$ and
$\|\bv_k\|\leq\|\bs_k\|$, squaring the preceding display gives
\begin{equation}\label{eq:repaired-iterate-difference}
 \frac1n\mathbb{E}\|\bx_{k+1}-\bx_k\|^2
 \leq32\mathsf{X}_k+8\alpha^2\mathsf{S}_k.
\end{equation}
Combining \eqref{eq:repaired-estimator-step} and
\eqref{eq:repaired-iterate-difference}, and using $(1-\tau)^2\leq1$ in the
forcing term, proves \eqref{eq:coupled-E-main}.

\emph{Step 2: tracking error.}
The identity $d_{i,k+1}=q_{i,k+1}$ need not hold when clipping is active.
However, $d_{i,k}\in B\mathbb B$ for every $k\geq0$ (at $k=0$ this follows
from Assumption~\ref{assum:stochatis}), and therefore
\begin{align}
 \|d_{i,k+1}-d_{i,k}\|
 &=\|\Pcal_{B\mathbb B}(q_{i,k+1})-
       \Pcal_{B\mathbb B}(d_{i,k})\|\nonumber\\
 &\leq\|q_{i,k+1}-d_{i,k}\|.                 \label{eq:clip-increment}
\end{align}
Expanding the last vector gives three terms:
\begin{align*}
 q_{i,k+1}-d_{i,k}
 ={}&\grad f_i(x_{i,k+1},\xi_{i,k+1})
       -\grad f_i(x_{i,k},\xi_{i,k+1})\\
 &+\tau\{\grad f_i(x_{i,k},\xi_{i,k+1})-\grad f_i(x_{i,k})\}
   +\tau\{\grad f_i(x_{i,k})-d_{i,k}\}.
\end{align*}
Thus, by $\|a+b+c\|^2\leq3(\|a\|^2+\|b\|^2+\|c\|^2)$, the oracle
assumptions, and \eqref{eq:repaired-iterate-difference},
\begin{equation}\label{eq:repaired-difference-d}
 \frac1n\mathbb{E}\|\bd_{k+1}-\bd_k\|^2
 \leq96\bar L^2\mathsf{X}_k+24\bar L^2\alpha^2\mathsf{S}_k
 +3\tau^2\mathsf{E}_k+3\tau^2\nu^2.
\end{equation}

Double stochasticity and $\hat\bs_k=\hat\bd_k$ imply
\begin{equation}\label{eq:tracking-exact-decomp}
 \bs_{k+1}-\hat{\bd}_{k+1}
 =((W-J)\otimes I_d)(\bs_k-\hat{\bd}_k)
  +((I-J)\otimes I_d)(\bd_{k+1}-\bd_k).
\end{equation}
For $\sigma_2>0$, apply
$\|a+b\|^2\leq(1+\zeta)\|a\|^2+(1+1/\zeta)\|b\|^2$ with
\[
 \zeta_T:=\frac{1-\sigma_2^2}{2\sigma_2^2}.
\]
Then
\[
 (1+\zeta_T)\sigma_2^2=q_T,
 \qquad 1+\frac1{\zeta_T}=h_T.
\]
Using $\|I-J\|=1$ in \eqref{eq:tracking-exact-decomp} and then
\eqref{eq:repaired-difference-d} proves \eqref{eq:coupled-T-main}. If
$\sigma_2=0$, the first term in \eqref{eq:tracking-exact-decomp} vanishes;
the same displayed bound remains valid because $q_T=1/2$ and $h_T=1$.

\emph{Step 3: dynamic consensus error.}
The pathwise projection estimate already established in the proof of
Theorem~\ref{lem:consensus} gives
\begin{equation}\label{eq:pathwise-dynamic-consensus}
 \|\bx_{k+1}-\bar\bx_{k+1}\|
 \leq\rho\|\bx_k-\bar\bx_k\|+\kappa_M\alpha\|\bs_k\|.
\end{equation}
Minkowski's inequality in $L^2$, divided by $\sqrt n$, yields
\[
 \sqrt{\mathsf{X}_{k+1}}
 \leq\rho\sqrt{\mathsf{X}_k}+\kappa_M\alpha\sqrt{\mathsf{S}_k}.
\]
When $\rho>0$, Young's inequality with
$\zeta_X=(1-\rho^2)/(2\rho^2)$ gives
$(1+\zeta_X)\rho^2=q_X$ and
$(1+1/\zeta_X)\kappa_M^2=c_X$, proving
\eqref{eq:coupled-X-main}. When $\rho=0$, direct squaring gives
$\mathsf X_{k+1}\leq\kappa_M^2\alpha^2\mathsf S_k$, which is stronger than
\eqref{eq:coupled-X-main} with $q_X=1/2$ and $c_X=\kappa_M^2$.

\emph{Step 4: descent with explicit projection remainders.}
We now derive \eqref{eq:coupled-D-main} without hiding the geometric terms.
Define the ambient consensus residual and its tangent projection by
\[
 c_{i,k}:=x_{i,k}-\sum_{j=1}^nW_{ij}x_{j,k},\qquad
 \grad\phi_i(\bx_k):=\Pcal_{T_{x_{i,k}}\Mcal}c_{i,k}.
\]
Let $\mathbf c_k$ denote the stack of the $c_{i,k}$. Then
$\sum_i c_{i,k}=0$ and
$\|\mathbf c_k\|=\|((I-W)\otimes I_d)\bx_k\|
\leq2\|\bx_k-\bar\bx_k\|$. Since the tangent projector is Lipschitz on the
compact submanifold, there is $H_T>0$ such that
\begin{equation}\label{eq:tangent-projector-lipschitz}
 \|\Pcal_{T_x\Mcal}-\Pcal_{T_y\Mcal}\|_{\rm op}
 \leq H_T\|x-y\|,\qquad x,y\in\Mcal.
\end{equation}
Consequently,
\begin{align}
 \left\|\sum_{i=1}^n\grad\phi_i(\bx_k)\right\|
 &=\left\|\sum_{i=1}^n
   (\Pcal_{T_{x_{i,k}}\Mcal}-\Pcal_{T_{\bar x_k}\Mcal})c_{i,k}\right\|
 \nonumber\\
 &\leq2H_T\|\bx_k-\bar\bx_k\|^2.             \label{eq:sum-consensus-grad}
\end{align}

Put $w_{i,k}:=c_{i,k}+\alpha v_{i,k}$. By
\eqref{eq:xk-barxk-bound} and \eqref{eq:sk-bound},
\[
 \max_{1\leq i\leq n}\|w_{i,k}\|
 \leq C_w\alpha,
 \qquad C_w:=2\sqrt{n C_B}+\sqrt{n D_B}.
\]
After reducing $\theta$ if necessary, this quantity lies in the uniform
radius of the second-order projection estimate
\eqref{projec-second-order1}. Define
\[
 r_{i,k}:=x_{i,k+1}-x_{i,k}+\grad\phi_i(\bx_k)+\alpha v_{i,k}.
\]
Because $v_{i,k}\in T_{x_{i,k}}\Mcal$,
\begin{equation}\label{eq:projection-remainder-block}
 \|r_{i,k}\|\leq Q\|w_{i,k}\|^2.
\end{equation}
Writing $\hat r_k=n^{-1}\sum_i r_{i,k}$, using
\eqref{eq:sum-consensus-grad}, \eqref{eq:projection-remainder-block}, and
$\|\bv_k\|\leq\|\bs_k\|$, we obtain
\begin{align}
 \|\hat x_{k+1}-\hat x_k+\alpha\hat v_k\|
 &\leq\frac1n\left\|\sum_i\grad\phi_i(\bx_k)\right\|
       +\frac Qn\sum_i\|w_{i,k}\|^2\nonumber\\
 &\leq\frac{C_m}{n}\|\bx_k-\bar\bx_k\|^2
       +\frac{2Q\alpha^2}{n}\|\bs_k\|^2,       \label{eq:mean-step-remainder}
\end{align}
where $C_m:=2H_T+8Q$; here we used
$\sum_i\|w_{i,k}\|^2\leq8\|\bx_k-\bar\bx_k\|^2
+2\alpha^2\|\bs_k\|^2$.

Let
\[
 \delta_k:=\bar x_{k+1}-\bar x_k+\alpha\hat v_k,
 \qquad \mathcal C_k:=\|\bx_k-\bar\bx_k\|^2.
\]
The mean-versus-projected-mean estimate \eqref{eq:distance-an-rm} and
\eqref{eq:mean-step-remainder} give
\[
 \|\delta_k\|\leq
 \frac{M_2+C_m}{n}\mathcal C_k
 +\frac{M_2}{n}\mathcal C_{k+1}
 +\frac{2Q\alpha^2}{n}\|\bs_k\|^2.
\]
Therefore,
\begin{equation}\label{eq:delta-k-explicit}
 \|\delta_k\|^2\leq
 \frac{\Gamma_1}{n^2}\mathcal C_k^2
 +\frac{\Gamma_2}{n^2}\mathcal C_{k+1}^2
 +\frac{\Gamma_3\alpha^4}{n^2}\|\bs_k\|^4,
\end{equation}
where
\[
 \Gamma_1:=3(M_2+C_m)^2,\qquad
 \Gamma_2:=3M_2^2,\qquad \Gamma_3:=12Q^2.
\]
Moreover,
\begin{equation}\label{eq:bar-step-explicit}
 \|\bar x_{k+1}-\bar x_k\|^2
 \leq\frac{2\alpha^2}{n}\|\bs_k\|^2+2\|\delta_k\|^2.
\end{equation}

Let $g_k:=\grad f(\bar x_k)$ and
$\mathbf g_k^{\rm av}:=\mathbf1_n\otimes g_k$. The Riemannian quadratic
upper bound and the identity
\begin{equation}\label{eq:polarization-full-direction}
 -\alpha\langle g_k,\hat s_k\rangle
 =-\frac\alpha2\|g_k\|^2-\frac{\alpha}{2n}\|\bs_k\|^2
  +\frac{\alpha}{2n}\|\mathbf g_k^{\rm av}-\bs_k\|^2
\end{equation}
give
\begin{align}
 f(\bar x_{k+1})
 \leq{}&f(\bar x_k)-\frac\alpha4\|g_k\|^2
 -\frac{\alpha}{2n}\|\bs_k\|^2
 +\frac{\alpha}{2n}\|\mathbf g_k^{\rm av}-\bs_k\|^2\nonumber\\
 &+\frac1\alpha\|\delta_k\|^2
 +\frac L2\|\bar x_{k+1}-\bar x_k\|^2
 +\frac\alpha n\sum_{i=1}^n
   \langle g_k,s_{i,k}-v_{i,k}\rangle.             \label{eq:descent-before-bounds}
\end{align}
Indeed, \eqref{eq:descent-before-bounds} follows by writing
$\bar x_{k+1}-\bar x_k+\alpha\hat s_k
=\delta_k+\alpha(\hat s_k-\hat v_k)$ and applying Young's inequality only
to $\langle g_k,\delta_k\rangle$.

We next bound the two error terms in \eqref{eq:descent-before-bounds}. Define
$\hat g_k=n^{-1}\sum_i\grad f_i(x_{i,k})$ and
$\hat d_k=n^{-1}\sum_i d_{i,k}$. Jensen's inequality gives
\begin{align}
 \frac1n\mathbb E\|\mathbf g_k^{\rm av}-\bs_k\|^2
 &\leq3L_G^2\mathsf X_k+3\mathsf E_k+3\mathsf T_k.   \label{eq:average-gradient-error}
\end{align}
To see this, insert $\mathbf1_n\otimes\hat g_k$ and
$\hat\bd_k=\mathbf1_n\otimes\hat d_k$ between
$\mathbf g_k^{\rm av}$ and $\bs_k$; the three squared terms are bounded by
$L_G^2\mathsf X_k$, $\mathsf E_k$, and $\mathsf T_k$, respectively.

For the normal term, set
\[
 A_{i,k}:=g_k-\Pcal_{T_{x_{i,k}}\Mcal}g_k.
\]
Since $g_k\in T_{\bar x_k}\Mcal$, compactness and
\eqref{eq:tangent-projector-lipschitz} give
$\|A_{i,k}\|\leq H_N\|x_{i,k}-\bar x_k\|$ for a constant $H_N>0$.
Moreover, $s_{i,k}-v_{i,k}\in N_{x_{i,k}}\Mcal$ and
$\|s_{i,k}-v_{i,k}\|\leq\|s_{i,k}\|$. Thus
$\langle g_k,s_{i,k}-v_{i,k}\rangle
=\langle A_{i,k},s_{i,k}-v_{i,k}\rangle$, and blockwise Young's inequality
gives, for every fixed $\eta>0$,
\begin{equation}\label{eq:repaired-normal-term}
 \frac\alpha n\sum_{i=1}^n
 \mathbb E\langle g_k,s_{i,k}-v_{i,k}\rangle
 \leq\frac{\alpha\eta}{2}\mathsf S_k
 +\frac{\alpha H_N^2}{2\eta}\mathsf X_k.
\end{equation}
The factor $\alpha$ appears exactly once in this bound.

Finally, the pathwise estimates \eqref{eq:xk-barxk-bound} and
\eqref{eq:sk-bound} imply from \eqref{eq:delta-k-explicit} that
\begin{equation}\label{eq:delta-order-explicit}
 \mathbb E\|\delta_k\|^2\leq C_\delta\alpha^4,
 \quad
 C_\delta:= (\Gamma_1+\Gamma_2)C_B^2+\Gamma_3D_B^2.
\end{equation}
Choose $\eta=1/4$, $\alpha\leq1$, and $L\alpha\leq1/8$. Substituting
\eqref{eq:bar-step-explicit}, \eqref{eq:average-gradient-error},
\eqref{eq:repaired-normal-term}, and \eqref{eq:delta-order-explicit} into
\eqref{eq:descent-before-bounds} gives \eqref{eq:coupled-D-main} with, for
example,
\[
 d_X:=\frac32L_G^2+2H_N^2,
 \qquad C_D:=(1+L)C_\delta.
\]
All these constants depend only on the problem data and the fixed tube, and
are independent of $k,K,\alpha$, and $\tau$.
\end{proof}
\end{revisionblock}

\end{document}
